\documentclass[11pt,a4paper]{article}
\usepackage[T1]{fontenc}
\usepackage[utf8]{inputenc}
\usepackage{lmodern}
\usepackage[margin=25mm]{geometry}
\usepackage{amsmath,amssymb,amsthm,mathtools}
\usepackage{microtype,booktabs,longtable,array}
\usepackage{placeins}
\usepackage{xcolor}
\usepackage{enumitem}
\usepackage{fancyhdr}
\usepackage{listings}
\definecolor{linkblue}{RGB}{27,57,94}
\usepackage[colorlinks=true,linkcolor=linkblue,citecolor=linkblue,urlcolor=linkblue]{hyperref}
\usepackage{bookmark}
\usepackage[nameinlink,noabbrev]{cleveref}
\hypersetup{pdftitle={The Baker--Campbell--Hausdorff Formula: A Unified Treatment with Complete Proofs, All-Order Expansions, and Analytic Qualifications},
 pdfauthor={},
 pdfsubject={BCH formula, Dynkin formula, free Lie algebras, Zassenhaus formula, exponential coordinates, Weyl relations},
 pdfkeywords={Baker--Campbell--Hausdorff, Dynkin, Zassenhaus, Bernoulli numbers, Lie groups, Weyl relations}}
\pagestyle{fancy}\fancyhf{}
\fancyhead[L]{\small The Baker--Campbell--Hausdorff formula}
\fancyhead[R]{\small A unified treatment}
\fancyfoot[C]{\thepage}
\renewcommand{\headrulewidth}{0.3pt}
\setlength{\parskip}{3pt}
\setlength{\emergencystretch}{2em}
\setlist{itemsep=2pt,topsep=4pt}
\allowdisplaybreaks[2]
\setcounter{tocdepth}{2}
\numberwithin{equation}{section}
\newtheorem{theorem}{Theorem}[section]
\newtheorem{lemma}[theorem]{Lemma}
\newtheorem{proposition}[theorem]{Proposition}
\newtheorem{corollary}[theorem]{Corollary}
\theoremstyle{definition}
\newtheorem{definition}[theorem]{Definition}
\newtheorem{example}[theorem]{Example}
\theoremstyle{remark}
\newtheorem{remark}[theorem]{Remark}
\DeclareMathOperator{\ad}{ad}
\DeclareMathOperator{\Ad}{Ad}
\DeclareMathOperator{\tr}{tr}
\DeclareMathOperator{\BCH}{BCH}
\DeclareMathOperator{\GL}{GL}
\DeclareMathOperator{\End}{End}
\DeclareMathOperator{\spec}{spec}
\DeclareMathOperator{\id}{id}
\DeclareMathOperator{\rev}{rev}
\DeclareMathOperator{\Prim}{Prim}
\DeclareMathOperator{\Dom}{Dom}
\newcommand{\kk}{\mathbb{K}}
\newcommand{\R}{\mathbb{R}}
\newcommand{\C}{\mathbb{C}}
\newcommand{\Q}{\mathbb{Q}}
\newcommand{\N}{\mathbb{N}}
\newcommand{\Z}{\mathbb{Z}}
\newcommand{\g}{\mathfrak{g}}
\newcommand{\h}{\mathfrak{h}}
\newcommand{\cL}{\mathcal{L}}
\newcommand{\cA}{\mathcal{A}}
\newcommand{\cT}{\mathcal{T}}
\newcommand{\cF}{\mathcal{F}}
\newcommand{\cS}{\mathcal{S}}
\newcommand{\eps}{\varepsilon}
\newcommand{\ii}{\mathrm{i}}
\newcommand{\dd}{\mathrm{d}}
\newcommand{\ot}{\,\widehat{\otimes}\,}
\newcommand{\norm}[1]{\left\lVert#1\right\rVert}
\newcommand{\abs}[1]{\left\lvert#1\right\rvert}
\newcommand{\coeff}[1]{\bigl[#1\bigr]}
\newcommand{\Rb}[1]{R(\mathrm{#1})}
\newcommand{\Lb}[1]{L(\mathrm{#1})}
\newcommand{\ket}[1]{\lvert#1\rangle}
\newcommand{\bra}[1]{\langle#1\rvert}
\newcommand{\braket}[2]{\langle#1\mid#2\rangle}
\newcommand{\bplus}[1]{b^{+}_{#1}}
\newcommand{\bminus}[1]{b^{-}_{#1}}
\newcommand{\word}[1]{\texttt{#1}}
\lstset{basicstyle=\small\ttfamily,columns=fullflexible,breaklines=true,frame=single,
 framerule=0.3pt,rulecolor=\color{black!30},showstringspaces=false,
 commentstyle=\itshape\color{black!60}}
\title{\textbf{The Baker--Campbell--Hausdorff Formula}\\[0.4em]
\large A Unified Treatment with Complete Proofs,\\All-Order Expansions, and Analytic Qualifications}
\author{}
\date{16 September 2026}
\begin{document}
\maketitle
\begin{abstract}
This article consolidates three independent proof-based treatments of the
mathematical content of the Wikipedia article on the Baker--Campbell--Hausdorff
formula into one self-contained account, keeping the strongest argument
available for each result. The formal logarithm of $e^Xe^Y$ is first computed
coefficient by coefficient in a completed free associative algebra, giving a
finite nonrecursive formula for every associative word coefficient. A direct
proof of the Dynkin--Specht--Wever identity establishes Friedrichs'
characterization of primitive elements and yields Dynkin's complete commutator
expansion together with a finite formula for every bidegree. An independent
differential route proves the Campbell conjugation identity, the differential
of the exponential, the Poincar\'e integral formula, complete Bernoulli
recurrences and partial fractions, an all-orders expansion in the number of
occurrences of one variable with an explicit bivariate kernel for its quadratic
part, and a finite-at-each-order analytic expansion around a regular base
point. We prove quantitative convergence and remainder bounds, the trace
identity with its branch restriction, and a complete classification of the
logarithms in the article's nonconvergence example. Lie-group consequences
include exponential coordinates, the local Lie correspondence, integration
from simply connected groups, and the polynomial group law of nilpotent
algebras. Central and eigenvector commutator relations are resummed globally.
For the Zassenhaus factorization we prove existence and uniqueness, three
recursive constructions, a finite nonrecursive weighted-tree formula for every
exponent, a globally convergent ordered-simplex expansion, and a quantitative
convergence criterion. Trotter, symmetric, and recursively higher-order
product formulas follow with complete logarithmic error series. The
infinitesimal formulas are developed into all-order coframe, pullback-metric,
and volume-density expansions, with the corrections needed in the source's
geometric paragraph. Finally the Weyl and displacement-operator identities are
proved in the Schr\"odinger and Bargmann--Fock representations, with operator
domains made explicit and a counterexample showing why a dense-domain
commutator alone does not suffice. Exact rational computations through degree
twelve accompany the proofs.
\end{abstract}
\vspace{0.5em}
\noindent\textbf{Keywords.} Free Lie algebra; noncommutative power series;
Dynkin formula; Bernoulli numbers; Zassenhaus factorization; Trotter and
Suzuki product formulas; exponential coordinates; Killing form; Weyl relations;
displacement operators.
\begingroup
\small
\setlength{\parskip}{0pt}
\tableofcontents
\endgroup
\newpage
\section{Scope, conventions, and the three meanings of BCH}\label{sec:scope}
The question that organizes this article is elementary to state: given two
quantities $X,Y$ whose products need not commute, can the product $e^Xe^Y$ be
written as a single exponential, and what exactly is its exponent? The central
answer is
\begin{equation}\label{eq:main}
 e^Xe^Y=e^{Z(X,Y)},\qquad
 Z(X,Y)=X+Y+\tfrac12[X,Y]+\tfrac1{12}[X,[X,Y]]+\tfrac1{12}[Y,[Y,X]]+\cdots.
\end{equation}
The substantive assertion is not the existence of a formal logarithm, which is
immediate in a completed free associative algebra. It is that every homogeneous
part of that logarithm is a \emph{Lie polynomial}, a linear combination of
nested commutators, and that its coefficients can be specified completely.

The equation $e^Xe^Y=e^Z$ is not, by itself, a definition of a unique
logarithm. The exponential can fail to be injective, and a product of
exponentials need not be an exponential \emph{in the specified Lie algebra}.
What is universal is the formal series
\begin{equation}\label{eq:formalbch}
 Z(X,Y)=\BCH(X,Y)=\log(e^Xe^Y)=\sum_{n\geq1}Z_n(X,Y),
\end{equation}
where $Z_n$ has total degree $n$. A principal matrix logarithm, a continued
branch of a logarithm, and a convergent BCH series are related but distinct
objects.

For definiteness we state at the outset, in both settings, the assertion that
the rest of the article proves and extends. Here $Z_n(X,Y)$ is the
homogeneous component of total degree $n$ produced by the degree recursion
\eqref{eq:degreerec}, and a \emph{closed Lie subalgebra} of a Banach algebra
is a closed linear subspace closed under the commutator.

\begin{theorem}[Baker--Campbell--Hausdorff]\label{thm:bch}
\begin{enumerate}[label=(\alph*)]
\item \emph{(Formal form, \eqref{eq:formalbch}.)} In the free associative
algebra $\kk\langle X,Y\rangle$ over a field $\kk$ of characteristic zero,
for every $n\geq1$ the element $Z_n$ is homogeneous of degree $n$, is a Lie
polynomial in $X$ and $Y$, and satisfies
\[
 Z_n=\sum_{|w|=n}\coeff{w}Z_n\;w
 =\frac1n\sum_{|w|=n}\coeff{w}Z_n\;R(w),
\]
where $R$ is the right-nested bracket \eqref{eq:rightbracket} and the
coefficients $\coeff{w}Z_n=c(w)$ are given by the finite sum
\eqref{eq:wordblocks}.
\item \emph{(Analytic form, \eqref{eq:main}.)} Let $\cA$ be a unital Banach
algebra over $\kk=\R$ or $\C$ with submultiplicative norm, and let
$X,Y\in\cA$ satisfy $\norm X+\norm Y<\log2$. Then $\sum_{n\geq1}Z_n(X,Y)$
converges absolutely; its sum $Z$ is the Mercator logarithm
$\log(e^Xe^Y)=\sum_{k\geq1}\frac{(-1)^{k-1}}k(e^Xe^Y-I)^k$ and satisfies
$e^Z=e^Xe^Y$; an element $Z_0$
with $\norm{Z_0}<\log2$ satisfies $e^{Z_0}=e^Xe^Y$ if and only if $Z_0=Z$;
every $Z_n(X,Y)$, and therefore $Z$ itself, lies in every closed Lie
subalgebra of $\cA$ that contains $X$ and $Y$; and
\[
 Z_1=X+Y,\qquad Z_2=\tfrac12[X,Y],\qquad
 Z_3=\tfrac1{12}[X,[X,Y]]+\tfrac1{12}[Y,[Y,X]].
\]
\end{enumerate}
\end{theorem}

Part (a) is \cref{thm:wordcut,thm:formalbch,thm:dynkin}, part (b) is
\cref{thm:convergence,prop:uniquelog,prop:closedlie,prop:lowdegree}; the
proofs are given there. Neither part asserts that $Z$ is a principal
logarithm, and neither is claimed outside its stated setting. In
particular part (b) does not assert $\norm Z<\log2$: the bound proved in
\cref{cor:bchlog} is $\norm Z\leq-\log(2-e^{\norm X+\norm Y})$, which
exceeds $\log2$ as $\norm X+\norm Y$ approaches $\log2$, and the
uniqueness clause is stated accordingly.

\subsection{Coverage target and what is not claimed}
The coverage target is the English Wikipedia article
\emph{Baker--Campbell--Hausdorff formula} in the revision with identifier
\texttt{1368909113}, last edited on 11 August 2026 and consulted on
16 September 2026 \cite{wiki}. Every displayed mathematical identity and every
mathematical assertion stated in its main text or its mathematical notes is
addressed in this article; \cref{app:coverage} gives a claim-by-claim
concordance. The article is used as a list of questions, not as a substitute
for proofs. Historical priority statements and application surveys are
bibliographical claims rather than theorems. A theorem merely named in a link,
such as the Stone--von Neumann uniqueness theorem or the Golden--Thompson
inequality, is not thereby stated on the page and is not part of the coverage
target; the actual quantum-mechanical exponential formulas on the page are
proved directly in \cref{sec:quantum}.

Several assertions on the page require correction or explicit hypotheses
rather than proof as written. The coframe matrix in exponential coordinates
is $W=\phi(\ad_X)$ and not the adjoint matrix $\ad_X$, which is always
singular in positive dimension; a Killing form need not be a metric; an
arbitrary pullback can be degenerate; the trace identity concerns the BCH
branch of the logarithm; and a commutator relation on a dense domain does not
by itself imply exponentiated Weyl relations. In each case we prove the valid
statement and exhibit a counterexample to the unqualified one.

The algebraic proofs include the structural results they use, in
particular the ordered Poincar\'e--Birkhoff--Witt theorem. Standard
foundations of analysis are assumed without proof: the inverse function
theorem, existence and uniqueness for ordinary differential equations with
analytic or bounded-linear right-hand sides, the identity theorem for
analytic functions, the contraction mapping theorem, Lebesgue's dominated
and monotone convergence theorems, uniqueness of Fourier coefficients for
continuous functions, and, in \cref{sec:quantum}, the spectral calculus of
self-adjoint operators, Stone's theorem on one-parameter unitary groups,
and the basic criterion for essential self-adjointness. No result about
the BCH formula itself is assumed.

The present article combines three earlier treatments of the same coverage
target that were prepared independently. Where those treatments proved the
same result by different routes, we keep the route that is most complete or
most elementary; where one contained a construction absent from the others,
that construction is included here. All proofs are written out in full.

\subsection{Three distinct settings}
Conflating the following settings is the source of most apparent paradoxes
surrounding the formula.
\begin{description}[leftmargin=1.5em,style=nextline]
\item[Formal algebra.] $X$ and $Y$ have positive degree in a graded algebra
over a field $\kk$ of characteristic zero. Every identity is an identity of
formal series, and every fixed-degree coefficient is a finite sum. No numerical
convergence is intended or asserted.
\item[Local analysis.] $X,Y$ belong to a real or complex unital Banach algebra
$\cA$ with a submultiplicative norm, or to a finite-dimensional Lie algebra,
and are sufficiently small. The series converges and specifies the logarithm
branch issuing from zero. Finite matrices are the principal example.
\item[Global identities under special relations.] Commutation relations may
allow the series to terminate or to be replaced by a closed expression. An
independent differential argument can then prove an exponential identity far
beyond the convergence domain of its Taylor series. Such an identity need not
identify a principal logarithm, and for unbounded operators domains cannot be
suppressed.
\end{description}
A matrix may have logarithms even though its BCH series diverges; a closed
expression may continue the local BCH logarithm without making the original
series convergent; and a logarithm need not belong to the Lie algebra generated
by its two inputs. Each of these phenomena is exhibited by an explicit example
below.

\subsection{Conventions}
Throughout, $[A,B]=AB-BA$ in an associative algebra, $\ad_AB=[A,B]$, and
compositions of adjoint maps are read from right to left. Associativity alone
gives antisymmetry and the Jacobi identity: expanding
$[A,[B,C]]+[B,[C,A]]+[C,[A,B]]$ produces each of the six ordered products
$ABC,ACB,BAC,BCA,CAB,CBA$ twice with opposite signs. Two consequences used
constantly are the derivation rule and its adjoint form,
\begin{equation}\label{eq:derivation}
 [A,[B,C]]=[[A,B],C]+[B,[A,C]],\qquad [\ad_A,\ad_B]=\ad_{[A,B]}.
\end{equation}
For a nonempty word $w=a_1\cdots a_n$ in the letters $X,Y$ its
\emph{right-nested bracket} is
\begin{equation}\label{eq:rightbracket}
 \Rb{a_1\cdots a_n}=[a_1,[a_2,\ldots,[a_{n-1},a_n]\ldots]],
 \qquad \Rb{a}=a,\qquad R(1)=0.
\end{equation}
Thus $\Rb{XXY}=[X,[X,Y]]=\ad_X^2Y$ and $\Rb{YYX}=\ad_Y^2X$. A symbol such as
$X^rY^s$ inside $R(\cdot)$ describes a word with $r$ letters $X$ followed by
$s$ letters $Y$, not a commutator of associative powers; empty runs are
omitted. We write $Z_n$ for the part of $Z$ of total degree $n$, and $Z_{p,q}$
for the part of degree $p$ in $X$ and $q$ in $Y$. For an associative
polynomial $P$ and a word $w$, $\coeff{w}P$ denotes the coefficient of $w$ in
$P$; it is never a commutator. The notation $O_{\deg}(m)$ denotes terms of
total degree at least $m$, and $O(Y^m)$ denotes terms containing at least $m$
occurrences of $Y$, with no claim of uniformity as $X$ varies.

The scalar functions that organize the entire subject are
\begin{equation}\label{eq:scalarfunctions}
\begin{gathered}
 \phi(z)=\frac{1-e^{-z}}{z},\qquad
 \phi(-z)=\frac{e^{z}-1}{z},\qquad
 \beta(z)=\frac{z}{1-e^{-z}},\\
 \gamma(z)=\beta(-z)=\frac{z}{e^z-1},\qquad
 \psi(x)=\frac{x\log x}{x-1}.
\end{gathered}
\end{equation}
The first two are entire, with value $1$ at zero. The functions $\beta$ and
$\gamma$ are meromorphic with a removable singularity at zero and simple poles
at the nonzero points of $2\pi\ii\Z$. The function $\psi$ is used near $x=1$,
with the logarithm branch analytic there and $\psi(1)=1$. An expression such
as $\phi(M)$ or $\beta(\ad_X)$ always denotes functional calculus, formal or
analytic, and never presupposes that $M$ or $\ad_X$ is invertible.

Bernoulli numbers appear in the convention $B_1^+=+\tfrac12$ of the source
page; we write $\bplus{n}=B_n^+/n!$ for the Taylor coefficients of $\beta$ and
$\bminus{n}=(-1)^n\bplus{n}=B_n^-/n!$ for those of $\gamma$.

\subsection{A guide to the results}
Sections \ref{sec:formal}--\ref{sec:dynkin} establish the algebraic theorem
and every coefficient: a finite formula for each associative word coefficient,
Dynkin's commutator formula, a finite formula for each bidegree, and the
complete polynomial through degree six with a finite certificate. Sections
\ref{sec:differential}--\ref{sec:integral} give the differential
constructions: the Campbell identity, the differential of the exponential,
Bernoulli coefficients, the Poincar\'e integral, every order in the number of
occurrences of $Y$, and two total-degree recursions.
\Cref{sec:convergence} proves quantitative convergence estimates, the trace
identity, and the nonconvergence example. \Cref{sec:liegroups} treats Lie
groups, and \cref{sec:special} the exactly summable cases.
\Cref{sec:zassenhaus,sec:splitting} treat the Zassenhaus factorization and the
product formulas. \Cref{sec:geometry,sec:quantum} treat exponential
coordinates and the quantum-mechanical identities. \Cref{sec:computation}
describes the exact computations. The appendices contain the ordered
Poincar\'e--Birkhoff--Witt theorem, coefficient certificates and further
complete degrees, and the coverage concordance.

A reader interested chiefly in coefficients can use the nonrecursive formulas
\eqref{eq:wordcut}, \eqref{eq:dynkin}, \eqref{eq:bidegree},
\eqref{eq:operatorblock}, and \eqref{eq:treeformula} directly.

\section{Formal exponentials and the complete word expansion}\label{sec:formal}
\subsection{The degree completion}
Let $V$ be the vector space over $\kk$ with basis $X,Y$, and put
\[
 T(V)=\bigoplus_{n\geq0}V^{\otimes n}=\kk\langle X,Y\rangle,
 \qquad
 \widehat T(V)=\prod_{n\geq0}V^{\otimes n}=\kk\langle\!\langle X,Y\rangle\!\rangle.
\]
The space $V^{\otimes n}=T_n$ has as basis the $2^n$ words of length $n$, the
empty word being the unit $1$. Multiplication concatenates words and is
extended to $\widehat T$ by the Cauchy product: the degree-$n$ part of a
product $AB$ is $\sum_{p+q=n}A_pB_q$, a finite sum. The filtration
$F^m\widehat T=\prod_{n\geq m}V^{\otimes n}$ satisfies
$F^mF^n\subseteq F^{m+n}$. A sum $\sum_j A_j$ of elements of $\widehat T$
\emph{converges formally} if for every degree $n$ only finitely many $A_j$
have a nonzero degree-$n$ part; its value is then defined degree by degree. No
norm is involved. In particular, if $A\in F^1\widehat T$, then $A^k\in
F^k\widehat T$, so for any scalar sequence $(c_k)$ the series
$\sum_{k\geq0}c_kA^k$ converges formally. Consequently substitution of
positive-degree series into any one-variable formal power series is well
defined, and
\begin{equation}\label{eq:explog}
 \exp A=\sum_{k\geq0}\frac{A^k}{k!},\qquad
 \log(1+U)=\sum_{k\geq1}\frac{(-1)^{k-1}}{k}U^k
 \qquad(A,U\in F^1\widehat T)
\end{equation}
are well-defined elements of $\widehat T$.

\begin{lemma}[Formal exponential and logarithm]\label{lem:explog}
The maps $\exp:F^1\widehat T\to1+F^1\widehat T$ and
$\log:1+F^1\widehat T\to F^1\widehat T$ of \eqref{eq:explog} are mutually
inverse bijections. If $A,B\in F^1\widehat T$ commute, then $e^{A+B}=e^Ae^B$.
If $g,h\in1+F^1\widehat T$ commute, then $\log(gh)=\log g+\log h$ and
$\log(g^{-1})=-\log g$.
\end{lemma}
\begin{proof}
We first prove the scalar identities $\log(e^t)=t$ and $e^{\log(1+t)}=1+t$
in $\Q[[t]]$. Put $E(t)=\sum_{k\geq0}t^k/k!$ and
$L(u)=\sum_{k\geq1}(-1)^{k-1}u^k/k$. Formal differentiation gives $E'=E$ and
$L'(u)=\sum_{k\geq1}(-1)^{k-1}u^{k-1}=(1+u)^{-1}$. The composition
$L(E(t)-1)$ has derivative $E'(t)/(1+E(t)-1)=E(t)/E(t)=1$ and vanishes at
$t=0$; a formal power series is determined by its derivative and constant
term, so $L(E(t)-1)=t$. Likewise $F(u)=E(L(u))$ satisfies
$F'(u)=E(L(u))L'(u)=F(u)/(1+u)$; the series $1+u$ satisfies the same
first-order equation with the same constant term $1$, and a solution of
$(1+u)F'=F$ is determined by its constant term because comparing the
coefficient of $u^n$ gives $(n+1)F_{n+1}=(1-n)F_n$ recursively. Hence
$E(L(u))=1+u$.

Now let $A\in F^1\widehat T$. The powers of the single element $A$ commute
with each other, so the subalgebra $\kk[[A]]$ of $\widehat T$ generated by $A$
is commutative, and the substitution $t\mapsto A$ defines a ring homomorphism
from $\Q[[t]]$ into $\widehat T$, well defined because each coefficient of
the image is a finite sum. Applying this homomorphism to $L(E(t)-1)=t$ gives
$\log(\exp A)=A$; applying the analogous homomorphism $u\mapsto U$ for
$U\in F^1\widehat T$ to $E(L(u))=1+u$ gives $\exp(\log(1+U))=1+U$. This
proves that the two maps are mutually inverse bijections.

If $A$ and $B$ commute, the binomial theorem holds for $(A+B)^n$, and
\[
 \sum_{n\geq0}\frac{(A+B)^n}{n!}
 =\sum_{n\geq0}\sum_{r+s=n}\frac{A^rB^s}{r!s!}
 =\sum_{r,s\geq0}\frac{A^rB^s}{r!s!}=e^Ae^B,
\]
all rearrangements being finite in each degree. If $g,h\in1+F^1\widehat T$
commute, then $\log g$ and $\log h$ commute, since each is a series in one
of the commuting elements $g-1$, $h-1$; hence
$e^{\log g+\log h}=e^{\log g}e^{\log h}=gh$, and applying $\log$ gives
$\log(gh)=\log g+\log h$. Taking $h=g^{-1}$, which exists as the geometric
series $\sum_{k\geq0}(1-g)^k$ and commutes with $g$, gives
$\log(g^{-1})=-\log g$. No commutativity of unrelated elements was used.
\end{proof}

In a Banach algebra the exponential series converges absolutely for every
$A$, since $\norm{A^n/n!}\leq\norm A^n/n!$, and the logarithm series converges
absolutely whenever $\norm U<1$. The scalar identities of the proof then hold
in the closed commutative subalgebra generated by $A$ or $U$, by the same
substitution argument with absolutely convergent series, so the analytic
statements corresponding to \cref{lem:explog} hold on those domains.

\subsection{Definition of the BCH series and the block expansion}
Define
\begin{equation}\label{eq:Zdef}
 Z(X,Y)=\BCH(X,Y)=\log(e^Xe^Y)=\sum_{n\geq1}Z_n(X,Y),\qquad Z_n\in V^{\otimes n}.
\end{equation}
By \cref{lem:explog}, $Z$ is the unique element of $F^1\widehat T$ with
$e^Z=e^Xe^Y$; thus \eqref{eq:main} already holds as an associative formal
identity. Expanding the two exponentials first and then the logarithm gives
\begin{equation}\label{eq:blocks}
 Z(X,Y)=\sum_{k\geq1}\frac{(-1)^{k-1}}{k}
 \sum_{\substack{r_i,s_i\geq0\\r_i+s_i>0\ (1\leq i\leq k)}}
 \frac{X^{r_1}Y^{s_1}\cdots X^{r_k}Y^{s_k}}
 {\prod_{i=1}^k r_i!s_i!}.
\end{equation}
Indeed $e^Xe^Y-1=\sum_{r+s>0}X^rY^s/(r!s!)$, and the $k$th power of this sum
is the sum over all $k$-tuples of nonempty blocks $X^{r_i}Y^{s_i}$. This
proves the associative expansion displayed on the source page. At this stage
the inner sum is noncommutative multiplication, not a sum over independent
commutators, and nothing has yet been proved about the Lie nature of the
homogeneous parts.

\subsection{A finite nonrecursive formula for every word}
For a nonempty word $w$ define its \emph{block weight}
\begin{equation}\label{eq:weight}
 a(w)=
 \begin{cases}
 \dfrac1{r!s!},&w=X^rY^s,\quad r,s\geq0,\quad r+s>0,\\[6pt]
 0,&\text{otherwise}.
 \end{cases}
\end{equation}
Thus $a(w)=0$ precisely when $w$ contains the adjacent pattern $YX$; pure
powers of one letter are permitted. A \emph{cut} of $w$ into $k$ pieces is a
factorization $w=w_1\cdots w_k$ into consecutive nonempty subwords.

\begin{theorem}[Every associative BCH coefficient]\label{thm:wordcut}
For every nonempty word $w=a_1\cdots a_n$,
\begin{equation}\label{eq:wordcut}
 c(w):=\coeff{w}Z(X,Y)
 =\sum_{k=1}^{n}\frac{(-1)^{k-1}}{k}
   \sum_{\substack{w=w_1\cdots w_k\\|w_i|>0}}\prod_{i=1}^k a(w_i)
 =\sum_{k=1}^{n}\frac{(-1)^{k-1}}{k}
   \sum_{0=i_0<i_1<\cdots<i_k=n}
   \prod_{j=1}^{k}a\bigl(a_{i_{j-1}+1}\cdots a_{i_j}\bigr).
\end{equation}
Consequently
\begin{equation}\label{eq:homblocks}
 Z_n=\sum_{w\in\{X,Y\}^n}c(w)\,w
 =\sum_{k=1}^{n}\frac{(-1)^{k-1}}{k}
 \sum_{\substack{r_i+s_i>0\\\sum_i(r_i+s_i)=n}}
 \frac{X^{r_1}Y^{s_1}\cdots X^{r_k}Y^{s_k}}
 {\prod_i r_i!s_i!}.
\end{equation}
Both formulas are finite and nonrecursive, and every coefficient is rational.
For a word of length $n$ there are exactly $2^{n-1}$ cut patterns before zero
terms are discarded.
\end{theorem}
\begin{proof}
Write $U=e^Xe^Y-1\in F^1\widehat T$. Its coefficient on a nonempty word $v$
is $a(v)$: the only words with nonzero coefficient in $e^Xe^Y$ are the
products $X^rY^s$, which arise once each, with coefficient $1/(r!s!)$. The
coefficient of $w$ in $U^k$ is obtained by expanding the product of $k$
copies of $U$: a word appears in the product exactly when it is the
concatenation $w_1\cdots w_k$ of one word from each factor, and its
coefficient is the sum over all such factorizations of the product of the
factor coefficients. Since every word in $U$ is nonempty, these
factorizations are precisely the cuts of $w$ into $k$ nonempty pieces, and
the coefficient is $\sum_{w=w_1\cdots w_k}\prod_ia(w_i)$. Because $w$ has
length $n$, no cut into more than $n$ nonempty pieces exists, so $U^k$ has
zero coefficient on $w$ for $k>n$. Substituting into
$Z=\sum_{k\geq1}(-1)^{k-1}U^k/k$ gives \eqref{eq:wordcut}; the second form
just parametrizes a cut by the positions $i_1<\cdots<i_{k-1}$ at which it
breaks. Choosing a subset of the $n-1$ gaps between consecutive letters is
the same as choosing a cut pattern, which gives the count $2^{n-1}$.
Expanding each permitted block as $X^{r_i}Y^{s_i}$ in \eqref{eq:wordcut}
and summing over words of length $n$ gives \eqref{eq:homblocks}, which is
the degree-$n$ part of \eqref{eq:blocks}.
\end{proof}

\begin{example}
For $w=XY$ the one-piece contribution is $a(XY)=1$ and the two-piece
contribution is $-\tfrac12a(X)a(Y)=-\tfrac12$, so $c(XY)=\tfrac12$. For
$w=YX$ the one-piece contribution is $a(YX)=0$, so $c(YX)=-\tfrac12$. For
$w=XXY$ the one-piece contribution is $a(XXY)=\tfrac12$; the two cuts
$X\mid XY$ and $XX\mid Y$ contribute $-\tfrac12(1\cdot1+\tfrac12\cdot1)
=-\tfrac34$; the three-piece cut contributes $\tfrac13$. Their sum is
$\tfrac12-\tfrac34+\tfrac13=\tfrac1{12}$. For $w=XYX$ the one-piece
contribution is zero, the two cuts $X\mid YX$ and $XY\mid X$ contribute
$-\tfrac12(0+1)=-\tfrac12$, and the three-piece cut contributes $\tfrac13$,
giving $c(XYX)=-\tfrac16$. These small computations already show why the
coefficients are rational and why cancellation between different cut
patterns matters.
\end{example}

\subsection{Any finite number of exponential factors}\label{sec:manyexp}
The same construction works without change for several factors.
\begin{corollary}[Several exponentials]\label{cor:manyexp}
Let $A_1,\ldots,A_m$ be letters. For a $k\times m$ array
$\rho=(r_{ji})$ of nonnegative integers in which every row has positive
sum, let $w(\rho)$ be the word obtained by concatenating the row words
$A_1^{r_{j1}}\cdots A_m^{r_{jm}}$ for $j=1,\ldots,k$, and put
$|\rho|=\sum_{j,i}r_{ji}$. Then
\begin{equation}\label{eq:manyexp}
 \log(e^{A_1}\cdots e^{A_m})
 =\sum_{k\geq1}\frac{(-1)^{k-1}}{k}
   \sum_{\rho}\frac{w(\rho)}{\prod_{j,i}r_{ji}!},
\end{equation}
and the coefficient of a given word $w$ is the finite sum over cuts of $w$
into $k$ nonempty pieces of the product of the block weights
$a_m(w_i)$, where $a_m(A_1^{r_1}\cdots A_m^{r_m})=(r_1!\cdots r_m!)^{-1}$
and $a_m$ vanishes on every other word.
Equivalently, with $t$ a formal scalar,
\begin{equation}\label{eq:multibch}
 \log\Bigl(\prod_{i=1}^{m}e^{tA_i}\Bigr)
 =\sum_{k\geq1}\frac{(-1)^{k-1}}{k}
 \sum_{\rho}\frac{t^{|\rho|}}{\prod_{j,i}r_{ji}!}\,w(\rho).
\end{equation}
\end{corollary}
\begin{proof}
Repeat the proof of \cref{thm:wordcut} with $U=e^{A_1}\cdots e^{A_m}-1$,
whose nonempty words are exactly the products $A_1^{r_1}\cdots A_m^{r_m}$
with coefficient $(r_1!\cdots r_m!)^{-1}$. The scalar $t$ simply counts
total degree.
\end{proof}
If the factors are of the form $e^{\alpha_jX}$ or $e^{\beta_jY}$ with the
letters $X,Y$ repeated in different positions, one first uses a distinct
letter for each factor and then substitutes; the resulting coefficients are
finite polynomials in the scalar weights $\alpha_j,\beta_j$. After
\cref{sec:hopf} it will follow that each homogeneous part of
\eqref{eq:manyexp} is a Lie polynomial and that each word $w(\rho)$ may be
replaced by $R(w(\rho))/|\rho|$. This observation supplies every order, not
only the leading error, for the splitting formulas of \cref{sec:splitting}.

\section{Hopf structure, primitive elements, and why the logarithm is a Lie series}\label{sec:hopf}
We prove the primitive-element characterization used by the source page
rather than treating it as a black box. The proof simultaneously constructs
the projection that produces Dynkin's formula.

\subsection{The coproduct must take values in a completion}
The relevant target of the coproduct on infinite series is the
\emph{completed} tensor product
\[
 \widehat T\ot\widehat T
 =\prod_{n\geq0}\bigoplus_{p+q=n}V^{\otimes p}\otimes V^{\otimes q},
\]
with multiplication and infinite sums interpreted by total degree exactly as
in $\widehat T$. The ordinary algebraic tensor product
$\widehat T\otimes\widehat T$ is generally too small: we show in
\cref{app:pbw} that already the coproduct of the one-variable series
$\sum_{n\geq0}X^n$ is not a finite sum of separated tensors. Define the
continuous algebra homomorphism
\begin{equation}\label{eq:coproduct}
 \Delta:\widehat T\longrightarrow\widehat T\ot\widehat T,
 \qquad
 \Delta X=X\otimes1+1\otimes X,\quad
 \Delta Y=Y\otimes1+1\otimes Y.
\end{equation}
Since $\widehat T$ is the completion of the free algebra on $X,Y$, an
algebra homomorphism on $T(V)$ is determined freely by the images of $X$ and
$Y$; and because $\Delta$ neither raises nor lowers total degree, it
extends uniquely by continuity to $\widehat T$. For a word
$w=a_1\cdots a_n$ and a subset $I\subseteq\{1,\ldots,n\}$ let $w_I$ be the
subword formed by the letters at the positions in $I$, in their original
order, with $w_\varnothing=1$. Then
\begin{equation}\label{eq:unshuffle}
 \Delta w=\sum_{I\subseteq\{1,\ldots,n\}}w_I\otimes w_{I^c}.
\end{equation}
This \emph{unshuffle} formula follows by multiplying out the $n$ factors
$a_j\otimes1+1\otimes a_j$: each choice of one term from each factor
assigns every letter to the left or right tensor factor, the letters
assigned to the left concatenate in their original order to $w_I$, and those
assigned to the right concatenate to $w_{I^c}$. Coassociativity,
$(\Delta\otimes\id)\Delta=(\id\otimes\Delta)\Delta$, follows in the same way:
either iterated coproduct assigns each letter independently to one of three
tensor factors, and the resulting sum over assignments does not depend on
the order in which the two splittings are performed. The counit
$\eps:\widehat T\to\kk$ extracts the constant term; the identities
$(\eps\otimes\id)\Delta=\id=(\id\otimes\eps)\Delta$ hold on words by
\eqref{eq:unshuffle}, since the only term with $w_I=1$ is $I=\varnothing$.

Define the antipode on words by
\begin{equation}\label{eq:antipode}
 S(a_1\cdots a_n)=(-1)^n a_n\cdots a_1,\qquad S(1)=1,
\end{equation}
extended linearly and continuously. It is an antiautomorphism:
$S(uv)=S(v)S(u)$, since reversing a concatenation reverses the factors in
the opposite order and the signs multiply. For linear maps $f,g$ on
$\widehat T$ the \emph{convolution} is $f*g=m(f\otimes g)\Delta$, where $m$
is multiplication; by coassociativity convolution is associative, with unit
$\eta\eps$, where $\eta$ inserts scalars.

\begin{lemma}[Antipode identities]\label{lem:antipode}
$\id*S=S*\id=\eta\eps$, that is, for every word $w$,
\begin{equation}\label{eq:antipodeidentities}
 m(\id\otimes S)\Delta w=m(S\otimes\id)\Delta w=\eps(w)1.
\end{equation}
\end{lemma}
\begin{proof}
Write $\Delta v=\sum v_{(1)}\otimes v_{(2)}$ in Sweedler notation for the
finite sum \eqref{eq:unshuffle} and put $K(v)=\sum v_{(1)}S(v_{(2)})$. For a letter $a$
and a word $v$, the unshuffle formula for $av$ splits according to whether
the first letter is assigned to the left or the right factor:
\[
 \Delta(av)=\sum av_{(1)}\otimes v_{(2)}+\sum v_{(1)}\otimes av_{(2)}.
\]
Using $S(av_{(2)})=-S(v_{(2)})a$ we obtain
\[
 K(av)=a\sum v_{(1)}S(v_{(2)})-\sum v_{(1)}S(v_{(2)})a=aK(v)-K(v)a.
\]
Since $K(1)=1\cdot S(1)=1$, we get $K(a)=a-a=0$ for every letter, and then
inductively $K(w)=0$ for every nonempty word. This proves the first identity.
For the second, put $K'(v)=\sum S(v_{(1)})v_{(2)}$ and split the coproduct of
$va$ according to the position of the last letter $a$:
$\Delta(va)=\sum v_{(1)}a\otimes v_{(2)}+\sum v_{(1)}\otimes v_{(2)}a$, whence
$K'(va)=-aK'(v)+K'(v)a$, and the same induction applies.
\end{proof}
Thus $T(V)$, and by continuity its completion, carries the asserted Hopf
algebra structure, with no appeal to an external structure theorem.

\begin{definition}
A series $P\in\widehat T$ is \emph{primitive} if
$\Delta P=P\otimes1+1\otimes P$. A series $G\in1+F^1\widehat T$ is
\emph{group-like} if $\Delta G=G\otimes G$.
\end{definition}
The normalization of the constant term is part of the definition of a
group-like element: the zero series satisfies $\Delta0=0\otimes0$ but does not
belong to a group. A primitive element has zero constant term, since a
primitive scalar $c$ would satisfy $c\otimes1=2c\otimes1$.

\begin{lemma}[Exponentials of primitives]\label{lem:grouplike}
Let $P\in F^1\widehat T$. Then $P$ is primitive if and only if $e^P$ is
group-like. The group-like elements form a group under multiplication.
\end{lemma}
\begin{proof}
Since $\Delta$ is a continuous algebra homomorphism, it commutes with the
exponential series: $\Delta(e^P)=e^{\Delta P}$. If $P$ is primitive, then
$\Delta P=P\otimes1+1\otimes P$ is a sum of two commuting elements, because
$(P\otimes1)(1\otimes P)=P\otimes P=(1\otimes P)(P\otimes1)$. By
\cref{lem:explog} applied in $\widehat T\ot\widehat T$,
\[
 \Delta(e^P)=e^{P\otimes1+1\otimes P}=e^{P\otimes1}e^{1\otimes P}
 =(e^P\otimes1)(1\otimes e^P)=e^P\otimes e^P.
\]
Conversely, if $G=e^P$ is group-like, then
$\Delta G=G\otimes G=(G\otimes1)(1\otimes G)$ is a product of two commuting
elements of $1+F^1(\widehat T\ot\widehat T)$, and the commuting-logarithm
rule of \cref{lem:explog} gives
\[
 \Delta P=\log(\Delta G)=\log(G\otimes1)+\log(1\otimes G)
 =P\otimes1+1\otimes P,
\]
where $\log(G\otimes1)=(\log G)\otimes1$ because $A\mapsto A\otimes1$ is a
continuous algebra homomorphism. Finally, if $G,H$ are group-like then
$\Delta(GH)=(G\otimes G)(H\otimes H)=GH\otimes GH$, and $G$ is invertible by
its geometric series, with
$\Delta(G^{-1})=(\Delta G)^{-1}=(G\otimes G)^{-1}=G^{-1}\otimes G^{-1}$.
\end{proof}

\subsection{An explicit projection: the Dynkin--Specht--Wever identity}
Let $N$ be the degree operator, $N(w)=|w|\,w$ and $N(1)=0$. It is a
derivation of $T(V)$, because degrees add under concatenation. Consider the
convolution
\begin{equation}\label{eq:dynkinoperator}
 D=N*S=m(N\otimes S)\Delta.
\end{equation}
The order of the two convolution factors is what produces
\emph{right}-nested rather than left-nested brackets. This convolution
operator is a special case of the generalized Dynkin operators studied in
\cite{menouspatras}; the elementary computation below contains the entire
argument required here.

\begin{lemma}[Right Dynkin operator]\label{lem:dsw}
For every nonempty word $w$, $D(w)=R(w)$. If $P$ is a homogeneous primitive
element of degree $n\geq1$, then
\begin{equation}\label{eq:dsw}
 R(P)=D(P)=nP,\qquad\text{that is,}\qquad
 P=\frac1n\sum_{|w|=n}\coeff{w}P\,R(w).
\end{equation}
\end{lemma}
\begin{proof}
From the definitions, $D(1)=N(1)S(1)=0$ and, for a letter $a$,
$D(a)=N(a)S(1)+N(1)S(a)=a$. Let $a$ be a letter and $w$ a nonempty word, and
split $\Delta(aw)$ as in the proof of \cref{lem:antipode}:
$\Delta(aw)=\sum aw_{(1)}\otimes w_{(2)}+\sum w_{(1)}\otimes aw_{(2)}$.
Using $N(aw_{(1)})=aw_{(1)}+aN(w_{(1)})$ (derivation property with
$N(a)=a$) and $S(aw_{(2)})=-S(w_{(2)})a$, we get
\begin{align*}
 D(aw)
 &=\sum aw_{(1)}S(w_{(2)})+\sum aN(w_{(1)})S(w_{(2)})
   -\sum N(w_{(1)})S(w_{(2)})a\\
 &=aK(w)+aD(w)-D(w)a,
\end{align*}
where $K(w)=\sum w_{(1)}S(w_{(2)})=\eps(w)1=0$ for nonempty $w$ by
\cref{lem:antipode}. Hence $D(aw)=[a,D(w)]$, and induction on the length
starting from $D(a)=a$ gives $D(a_1\cdots a_n)=[a_1,D(a_2\cdots a_n)]
=\cdots=R(a_1\cdots a_n)$.

If $P$ is homogeneous primitive of degree $n$, then its coproduct has only
the two terms $P\otimes1$ and $1\otimes P$, and
\[
 D(P)=m(N\otimes S)(P\otimes1+1\otimes P)=N(P)S(1)+N(1)S(P)=nP+0.
\]
Since $D$ is linear and $D(w)=R(w)$ on words, $D(P)=\sum_w\coeff{w}P\,R(w)$,
which gives the second form of \eqref{eq:dsw} after division by $n$.
\end{proof}

Let $\cL=\cL(X,Y)\subseteq T(V)$ be the Lie subalgebra generated by $X,Y$
under the commutator bracket, and $\cL_n=\cL\cap V^{\otimes n}$. Since
$X,Y$ are homogeneous, $\cL$ is spanned by homogeneous bracket monomials, so
$\cL=\bigoplus_n\cL_n$; we write $\widehat\cL=\prod_{n\geq1}\cL_n$.

\begin{theorem}[Friedrichs' characterization; Dynkin--Specht--Wever]\label{thm:friedrichs}
A homogeneous element $P\in V^{\otimes n}$, $n\geq1$, is primitive if and
only if $P\in\cL_n$. For such $P$, \eqref{eq:dsw} holds, so $D/n$ is an
idempotent projection of $V^{\otimes n}$ onto $\cL_n$ whose kernel consists
of the degree-$n$ elements annihilated by $D$. The primitive elements of
$\widehat T$ are exactly the elements of $\widehat\cL$, the formal
degreewise sums of Lie polynomials:
\begin{equation}\label{eq:friedrichs}
 \Prim(\widehat T)=\prod_{n\geq1}\cL_n.
\end{equation}
\end{theorem}
\begin{proof}
The generators $X,Y$ are primitive by definition of $\Delta$. If $P,Q$ are
primitive, then
\[
 \Delta[P,Q]=[\Delta P,\Delta Q]
 =[P\otimes1+1\otimes P,\ Q\otimes1+1\otimes Q]
 =[P,Q]\otimes1+1\otimes[P,Q],
\]
because the mixed terms cancel: $(P\otimes1)(1\otimes Q)=P\otimes Q
=(1\otimes Q)(P\otimes1)$, and similarly with $P$ and $Q$ interchanged.
Hence every bracket monomial in $X,Y$ is primitive, every Lie polynomial is
primitive by linearity, and every degreewise-convergent sum of Lie
polynomials is primitive by continuity of $\Delta$.

Conversely, let $P\in\widehat T$ be primitive. Its constant term vanishes,
as noted above. Since $\Delta$ preserves total degree, comparing the
degree-$n$ parts of $\Delta P=P\otimes1+1\otimes P$ shows that each
homogeneous component $P_n$ is primitive. By \cref{lem:dsw},
$P_n=\frac1n\sum_{|w|=n}\coeff{w}P_n\,R(w)$ is a linear combination of
right-nested brackets, hence an element of $\cL_n$. Division by $n$ is
exactly where characteristic zero is used. This proves the homogeneous
characterization and \eqref{eq:friedrichs}.

Finally, since $D$ always takes values in $\cL$ (its value on any word is a
nested bracket) and every element of $\cL_n$ is primitive with $D(P)=nP$,
we have $D(D(Q))=nD(Q)$ for every $Q\in V^{\otimes n}$; so $D/n$ is
idempotent with image $\cL_n$.
\end{proof}

\begin{theorem}[Formal BCH theorem]\label{thm:formalbch}
There is a unique $Z\in F^1\widehat T$ with $e^Z=e^Xe^Y$. Every homogeneous
component $Z_n$ is a Lie polynomial with rational coefficients, $Z_1=X+Y$,
and $Z-X-Y$ belongs degreewise to the Lie ideal of $\cL$ generated by
$[X,Y]$.
\end{theorem}
\begin{proof}
Existence and uniqueness follow from \cref{lem:explog}. Since $X$ and $Y$
are primitive, $e^X$ and $e^Y$ are group-like by \cref{lem:grouplike}, their
product $e^Xe^Y$ is group-like because group-like elements form a group, and
$Z=\log(e^Xe^Y)$ is primitive by \cref{lem:grouplike} again. By
\cref{thm:friedrichs} each $Z_n$ lies in $\cL_n$. Rationality follows from
\eqref{eq:wordcut}, whose ingredients are factorials and the coefficients
$(-1)^{k-1}/k$. Taking degree one in \eqref{eq:blocks} gives $Z_1=X+Y$.

For the ideal assertion, let $J$ be the Lie ideal of $\cL$ generated by
$[X,Y]$. A right-nested bracket $R(w)$ with $|w|\geq2$ has innermost
bracket $[a_{n-1},a_n]$, which is either zero or $\pm[X,Y]$, and an iterated
bracket one of whose factors lies in the ideal $J$ lies in $J$; since each
$Z_n$ with $n\geq2$ is a combination of such brackets by \eqref{eq:dsw}, it
lies in $J$. (Equivalently, in the quotient Lie algebra $\cL/J$ the two
generators commute, so $\cL/J$ is abelian and the image of $Z$ under the
Lie homomorphism $\cL\to\cL/J$ is $\bar X+\bar Y$ by \cref{prop:symmetry}
applied in an abelian Lie algebra, which again says $Z_n\in J$ for
$n\geq2$.)
\end{proof}
The phrase ``every term beyond $X+Y$ contains $[X,Y]$'' is thus an
ideal-theoretic statement; it does not mean that an arbitrary display of
$Z_n$ must visibly show $[X,Y]$ as the innermost bracket of every summand.

\subsection{Universality and enveloping algebras}\label{sec:universal}
A homogeneous Lie polynomial can be evaluated at any pair of elements of any
Lie algebra $\g$ over $\kk$. The identification of $\cL(X,Y)$ with the
abstract free Lie algebra, and the fact that the free associative algebra is
its universal enveloping algebra, rest on the injectivity of $\g\to U(\g)$,
which is proved in \cref{app:pbw} by the ordered Poincar\'e--Birkhoff--Witt
theorem (\cref{thm:pbw}). The argument is as follows. Given $x,y\in\g$, the
associative homomorphism $T(V)\to U(\g)$ determined by $X\mapsto x$,
$Y\mapsto y$ sends every iterated commutator of $X,Y$ to the corresponding
iterated bracket of $x,y$, which lies in the image of $\g$. Its restriction
to $\cL$ is therefore a Lie homomorphism $\cL\to\g$, and it is the unique
one sending $X,Y$ to $x,y$ because $X,Y$ generate $\cL$. This is exactly the
universal property of the free Lie algebra. Consequently an associative
polynomial identity between Lie polynomials in $X,Y$, verified as an
equality of word coefficients in $T(V)$, is an identity valid under
substitution in \emph{every} Lie algebra.

\begin{proposition}[Enveloping algebra of the free Lie algebra]\label{prop:enveloping}
The universal enveloping algebra of $\cL(X,Y)$ is canonically $T(V)$, and
the Hopf structure of \cref{sec:hopf} is the enveloping Hopf structure
determined by $\Delta u=u\otimes1+1\otimes u$, $\eps(u)=0$, $S(u)=-u$ for
$u\in\cL$. The same formulas define a Hopf algebra structure on $U(\g)$ for
every Lie algebra $\g$.
\end{proposition}
\begin{proof}
The inclusion $\cL\to T(V)$ is a Lie homomorphism into the commutator
algebra of $T(V)$, so by the universal property of $U(\cL)$ it extends to an
associative homomorphism $\alpha:U(\cL)\to T(V)$. The assignment
$X\mapsto X$, $Y\mapsto Y$ into $U(\cL)$ extends, by freeness of $T(V)$, to
an associative homomorphism $\beta:T(V)\to U(\cL)$. Both composites
$\alpha\beta$ and $\beta\alpha$ fix $X$ and $Y$; since $X,Y$ generate
$T(V)$ as an algebra and generate $U(\cL)$ as an algebra (because they
generate $\cL$ as a Lie algebra and $\cL$ generates $U(\cL)$), both
composites are identities. The Hopf assertion for general $\g$ is proved in
\cref{app:pbw}; on generators the two structures agree, and each is
determined by its values on generators.
\end{proof}
For an arbitrary Lie algebra $\g$ without a topology, the unambiguous
universal statement is
\[
 \BCH(tx,ty)=\sum_{n\geq1}t^nZ_n(x,y)\in t\,\g[[t]],
\]
obtained by evaluating in $U(\g)[[t]]$. An infinite sum at $t=1$ in the
untopologized vector space $\g$ has no meaning until nilpotence
(\cref{sec:nilpotent}), a complete filtration, or a convergence argument
(\cref{sec:convergence}) has been supplied.

\section{Dynkin's formula, every bidegree, and the polynomial through degree six}\label{sec:dynkin}
\subsection{Dynkin's complete commutator formula}
The structural theorem becomes completely explicit by applying the Dynkin
operator to the associative logarithm. The classical coefficient formula is
due to Dynkin \cite{dynkin}; modern computational treatments include
\cite{casasmurua}. The proof below derives its coefficients directly.

\begin{theorem}[Dynkin's all-terms formula]\label{thm:dynkin}
With $d=\sum_{i=1}^k(r_i+s_i)$,
\begin{equation}\label{eq:dynkin}
 \boxed{\displaystyle
 Z(X,Y)=\sum_{k\geq1}\frac{(-1)^{k-1}}{k}
 \sum_{\substack{r_i,s_i\geq0\\r_i+s_i>0\ (1\leq i\leq k)}}
 \frac{R\bigl(X^{r_1}Y^{s_1}\cdots X^{r_k}Y^{s_k}\bigr)}
      {d\,\prod_{i=1}^k r_i!s_i!}.}
\end{equation}
Equivalently, in terms of the word coefficients \eqref{eq:wordcut},
\begin{equation}\label{eq:dynkinword}
 Z_n=\frac1n\sum_{|w|=n}c(w)\,R(w),
\end{equation}
and, separating the bidegrees,
\begin{equation}\label{eq:bidegreeproj}
 \boxed{\displaystyle
 Z_{p,q}=\frac1{p+q}
 \sum_{\substack{w\in\{X,Y\}^{p+q}\\\#_Xw=p,\ \#_Yw=q}}
 c(w)\,R(w),\qquad p+q\geq1.}
\end{equation}
In each total degree the sums are finite. A summand of \eqref{eq:dynkin}
vanishes if $s_k>1$, or if $s_k=0$ and $r_k>1$.
\end{theorem}
\begin{proof}
By \cref{thm:formalbch}, each $Z_n$ is a homogeneous primitive element of
degree $n$. \Cref{lem:dsw} therefore gives $Z_n=\frac1nD(Z_n)$. Since $D$ is
linear and $D(w)=R(w)$ on words, applying $D/n$ to the word expansion
$Z_n=\sum_{|w|=n}c(w)w$ of \eqref{eq:homblocks} gives \eqref{eq:dynkinword}.
Applying it instead to the block form of \eqref{eq:homblocks} replaces each
block word $X^{r_1}Y^{s_1}\cdots X^{r_k}Y^{s_k}$ by its right-nested bracket
and divides by its length $d=n$; summing over $n$ gives \eqref{eq:dynkin}.
Finally, $Z_n=\sum_{p+q=n}Z_{p,q}$ and the operator $D$ preserves bidegree
(it maps a word to a combination of words with the same letters), so the
bidegree-$(p,q)$ part of \eqref{eq:dynkinword} is \eqref{eq:bidegreeproj}.
If $s_k>1$ the two innermost letters of the bracketed word are $Y,Y$, and
if $s_k=0$, $r_k>1$ they are $X,X$; in either case the innermost bracket
$[a_{d-1},a_d]$ vanishes and so does the whole nested bracket.
\end{proof}
The two vanishing tests are sufficient but not necessary: a run of equal
letters may continue across a block boundary, for instance for the three
blocks $X$, $XY$, $Y$, whose concatenation $XXYY$ has innermost bracket
$[Y,Y]=0$ although $s_k=1$. Notice also that one must project a
\emph{whole homogeneous Lie polynomial}; applying $D/n$ to an arbitrary
associative polynomial generally changes it. The right-nested brackets form a
spanning family of $\cL_n$, not a basis, so different displays can represent
the same Lie element by Jacobi identities. This is redundancy in the
representation, not ambiguity in $Z$. Formula \eqref{eq:bidegreeproj},
combined with the finite cut formula \eqref{eq:wordcut}, is a particularly
compact answer to the request for \emph{all} terms: there are no unspecified
coefficients or auxiliary recurrences.

\subsection{Symmetries, associativity, and functoriality}
\begin{proposition}\label{prop:symmetry}
As universal formal series, with $T$ a third letter,
\begin{align}
 Z(Y,X)&=-Z(-X,-Y),\label{eq:swap}\\
 Z_n(Y,X)&=(-1)^{n+1}Z_n(X,Y),\label{eq:parity}\\
 \rev\bigl(Z_n(X,Y)\bigr)&=(-1)^{n+1}Z_n(X,Y),
 \qquad\text{that is,}\qquad c(\rev w)=(-1)^{|w|+1}c(w),\label{eq:reversal}\\
 Z\bigl(Z(X,Y),T\bigr)&=Z\bigl(X,Z(Y,T)\bigr),\label{eq:assoc}\\
 Z(X,0)&=Z(0,X)=X,\qquad Z(X,-X)=0.\label{eq:unitinverse}
\end{align}
Here $\rev$ reverses every word. Moreover, every Lie-algebra homomorphism
$f:\g\to\h$ satisfies $f(Z_n(x,y))=Z_n(fx,fy)$ for all $n$, and hence
commutes with the BCH series wherever formal substitution or analytic
convergence is valid.
\end{proposition}
\begin{proof}
Taking inverses in $e^{-X}e^{-Y}=e^{Z(-X,-Y)}$ gives
$e^Ye^X=(e^{Z(-X,-Y)})^{-1}=e^{-Z(-X,-Y)}$, using $(e^A)^{-1}=e^{-A}$ from
\cref{lem:explog}. Since also $e^Ye^X=e^{Z(Y,X)}$, uniqueness of the formal
logarithm gives \eqref{eq:swap}. Replacing $(X,Y)$ by $(-X,-Y)$ multiplies
the degree-$n$ part by $(-1)^n$, so comparing homogeneous parts in
\eqref{eq:swap} gives \eqref{eq:parity}. Word reversal $\rev$ is an
antiautomorphism of $\widehat T$: $\rev(uv)=\rev(v)\rev(u)$. It fixes $X$
and $Y$, hence $\rev(e^X)=e^X$, $\rev(e^Y)=e^Y$, and
$\rev(e^Xe^Y)=e^Ye^X$. Because $\rev$ is linear, continuous, and
multiplicative on powers of a single element, it commutes with the
logarithm series: $\rev(\log(1+U))=\log(1+\rev U)$. Therefore
$\rev(Z(X,Y))=\log(e^Ye^X)=Z(Y,X)$, and \eqref{eq:parity} gives
\eqref{eq:reversal}; the coefficient statement is its coordinate form.

For \eqref{eq:assoc}, work in the completion of the free algebra on three
letters $X,Y,T$. Both sides exponentiate to $e^Xe^Ye^T$: the left side
because $e^{Z(Z(X,Y),T)}=e^{Z(X,Y)}e^T=e^Xe^Ye^T$, the right side
similarly; here we use that $Z(A,B)$ for positive-degree $A,B$ is defined
by substitution into the universal series and satisfies $e^{Z(A,B)}=e^Ae^B$
by the substitution homomorphism. Uniqueness of the logarithm proves
\eqref{eq:assoc}. The identities \eqref{eq:unitinverse} follow in the same
way from $e^Xe^0=e^X$ and $e^Xe^{-X}=1$.

A Lie homomorphism $f$ preserves every iterated bracket, hence every
right-nested bracket $R(w)$ evaluated at $(x,y)$, hence every $Z_n$ by
\eqref{eq:dynkinword}. In a convergent analytic setting a continuous linear
map preserves the limit of the series; a linear map between
finite-dimensional spaces is automatically continuous.
\end{proof}
Both the swap symmetry and the reversal symmetry provide useful independent
checks on large coefficient tables, and they are used as such in
\cref{sec:computation}.

\subsection{Every term through degree six}
With the right-nesting convention \eqref{eq:rightbracket}, the terms printed
on the source page are exactly
\begin{align}
 Z_1={}&X+Y,\label{eq:z1}\\
 Z_2={}&\tfrac12\Rb{XY}=\tfrac12[X,Y],\label{eq:z2}\\
 Z_3={}&\tfrac1{12}\bigl(\Rb{XXY}+\Rb{YYX}\bigr)
       =\tfrac1{12}[X,[X,Y]]+\tfrac1{12}[Y,[Y,X]],\label{eq:z3}\\
 Z_4={}&-\tfrac1{24}\Rb{YXXY}=-\tfrac1{24}[Y,[X,[X,Y]]],\label{eq:z4}\\
 Z_5={}&-\tfrac1{720}\bigl(\Rb{YYYYX}+\Rb{XXXXY}\bigr)
       +\tfrac1{360}\bigl(\Rb{XYYYX}+\Rb{YXXXY}\bigr)
       \notag\\
     &+\tfrac1{120}\bigl(\Rb{YXYXY}+\Rb{XYXYX}\bigr),\label{eq:z5}\\
 Z_6={}&\tfrac1{240}\Rb{XYXYXY}
       +\tfrac1{720}\bigl(\Rb{XYXXXY}-\Rb{XXYYXY}\bigr)
       \notag\\
     &+\tfrac1{1440}\bigl(\Rb{XYYYXY}-\Rb{XXYXXY}\bigr).\label{eq:z6}
\end{align}
For example $\Rb{XYXXXY}=[X,[Y,[X,[X,[X,Y]]]]]$. There is no omitted term
of any of these degrees; equations \eqref{eq:wordcut} and
\eqref{eq:dynkin} define every subsequent degree.

For direct comparison with associative matrix multiplication, the first
four degrees are
\begin{align}
 Z_2={}&\tfrac12(XY-YX),\label{eq:assoc2}\\
 Z_3={}&\tfrac1{12}\bigl(X^2Y+XY^2-2XYX+Y^2X+YX^2-2YXY\bigr),\label{eq:assoc3}\\
 Z_4={}&\tfrac1{24}\bigl(X^2Y^2-2XYXY-Y^2X^2+2YXYX\bigr).\label{eq:assoc4}
\end{align}

To turn displayed brackets into word coefficients we use the following
explicit expansion. For $w=a_1\cdots a_n$ and a subset
$S\subseteq\{1,\ldots,n-1\}$, let $w_S$ be the subword of $a_1\cdots a_{n-1}$
at the positions in $S$ and $w_{S^c}$ the subword at the remaining positions
of $\{1,\ldots,n-1\}$; in contrast with \eqref{eq:unshuffle}, the index set
here omits the last position. Then
\begin{equation}\label{eq:bracketexpand}
 R(w)=\sum_{S\subseteq\{1,\ldots,n-1\}}(-1)^{n-1-|S|}\,w_S\,a_n\,\rev(w_{S^c}).
\end{equation}
\begin{proof}[Proof of \eqref{eq:bracketexpand}]
Induct on $n$; the case $n=1$ reads $R(a_1)=a_1$. For $n\geq2$,
$R(w)=a_1R(a_2\cdots a_n)-R(a_2\cdots a_n)a_1$. By the induction hypothesis
each term of $R(a_2\cdots a_n)$ has the form
$\pm(\text{subword of }a_2\cdots a_{n-1})\,a_n\,(\text{reversed complementary
subword})$. Placing $a_1$ on the left adds position $1$ to $S$ and keeps the
sign; placing $a_1$ on the right, with a minus sign, appends $a_1$ at the end
of the reversed complement, which is exactly where position $1$ belongs in
$\rev(w_{S^c})$ when $1\notin S$, and changes the sign, matching the factor
$(-1)^{n-1-|S|}$ because $|S^c|$ has increased by one.
\end{proof}

\begin{proposition}\label{prop:lowdegree}
Equations \eqref{eq:z1}--\eqref{eq:z6} are the homogeneous components of
degrees one through six of $Z(X,Y)$, and \eqref{eq:assoc2}--\eqref{eq:assoc4}
are their associative expansions.
\end{proposition}
\begin{proof}
Degree one is \cref{thm:formalbch}. For degree two, \eqref{eq:wordcut}
gives $c(XY)=\tfrac12$ and $c(YX)=-\tfrac12$, while
$c(XX)=a(XX)-\tfrac12a(X)a(X)=\tfrac12-\tfrac12=0$ and likewise
$c(YY)=0$. Hence $Z_2=\tfrac12(XY-YX)=\tfrac12[X,Y]$, which is
\eqref{eq:z2} and \eqref{eq:assoc2}.

For degree three the cut formula gives the following values. For $XXY$ we
computed $\tfrac1{12}$; by reversal symmetry \eqref{eq:reversal} with the
sign $(-1)^{4}=+1$, $c(YXX)=\tfrac1{12}$, and by swap symmetry
\eqref{eq:parity}, $c(YYX)=c(XYY)=\tfrac1{12}$. For $XYX$ we computed
$-\tfrac16$, and by swap, $c(YXY)=-\tfrac16$. For $XXX$: the one-piece
contribution is $\tfrac16$, the two cuts contribute
$-\tfrac12(\tfrac12+\tfrac12)=-\tfrac12$, and the three-piece cut
contributes $\tfrac13$; the sum is zero, and likewise $c(YYY)=0$. Thus
\[
 Z_3=\tfrac1{12}\bigl(XXY+YXX+YYX+XYY-2XYX-2YXY\bigr),
\]
which is \eqref{eq:assoc3}. On the other hand, expanding the brackets,
\[
 [X,[X,Y]]=X^2Y-2XYX+YX^2,\qquad
 [Y,[Y,X]]=Y^2X-2YXY+XY^2,
\]
and $\tfrac1{12}$ times their sum is the same polynomial. This proves
\eqref{eq:z3}.

For degree four, direct evaluation of \eqref{eq:wordcut} on the sixteen
words gives nonzero values only for
\[
 c(XXYY)=\tfrac1{24},\quad c(XYXY)=-\tfrac1{12},\quad
 c(YXYX)=\tfrac1{12},\quad c(YYXX)=-\tfrac1{24},
\]
and zero on the remaining twelve words. We verify one representative of
each type. For $XXYY$: one piece gives $a(XXYY)=\tfrac14$; the three
two-piece cuts $X\mid XYY,\ XX\mid YY,\ XXY\mid Y$ give
$-\tfrac12(\tfrac12+\tfrac14+\tfrac12)=-\tfrac58$; the three cuts into three pieces,
$X\mid X\mid YY,\ X\mid XY\mid Y,\ XX\mid Y\mid Y$, give
$\tfrac13(\tfrac12+1+\tfrac12)=\tfrac23$; the four-piece cut gives
$-\tfrac14$; the sum is $\tfrac14-\tfrac58+\tfrac23-\tfrac14=\tfrac1{24}$.
For $XYXY$: one piece gives $0$; two-piece cuts $X\mid YXY$ ($0$),
$XY\mid XY$ ($1$), $XYX\mid Y$ ($0$) give $-\tfrac12$; three-piece cuts
$X\mid Y\mid XY$, $X\mid YX\mid Y$, $XY\mid X\mid Y$ give
$\tfrac13(1+0+1)=\tfrac23$; four pieces give $-\tfrac14$; the sum is
$-\tfrac12+\tfrac23-\tfrac14=-\tfrac1{12}$. The values at $YYXX$ and $YXYX$
follow by reversal symmetry with sign $(-1)^5=-1$. For a word with zero
coefficient, take $XXXY$: one piece gives $a(XXXY)=\tfrac16$; the two-piece
cuts $X\mid XXY$, $XX\mid XY$, $XXX\mid Y$ give
$-\tfrac12(\tfrac12+\tfrac12+\tfrac16)=-\tfrac7{12}$; the three-piece cuts
$X\mid X\mid XY$ ($1$), $X\mid XX\mid Y$ ($\tfrac12$), $XX\mid X\mid Y$
($\tfrac12$) give $\tfrac13\cdot2=\tfrac23$; four pieces give $-\tfrac14$;
the sum is $\tfrac16-\tfrac7{12}+\tfrac23-\tfrac14=0$. The
remaining zero coefficients are checked in the same way; the eleven other
words are $XXXX$, $YYYY$, $XXYX$, $XYXX$, $YXXX$, $XYYY$, $YXYY$, $YYXY$,
$YYYX$, $XYYX$, $YXXY$, and each cut sum vanishes (for the last two, for
example, the contributions are $0-\tfrac12(0+0+\tfrac12)+\tfrac13(0+\tfrac12+1)
-\tfrac14=0$ and $0-\tfrac12(\tfrac12+0+0)+\tfrac13(1+\tfrac12+0)-\tfrac14=0$,
the cuts being listed from left to right by the position of the first break).
Thus \eqref{eq:assoc4} holds. Expanding the bracket,
\[
 [Y,[X,[X,Y]]]=[Y,\,X^2Y-2XYX+YX^2]
 =2XYXY-2YXYX+Y^2X^2-X^2Y^2,
\]
and $-\tfrac1{24}$ times this is \eqref{eq:assoc4}. This proves
\eqref{eq:z4}.

For degrees five and six we give a finite coefficientwise certificate.
Expand each right-nested bracket in \eqref{eq:z5} and \eqref{eq:z6} by
\eqref{eq:bracketexpand} and collect the coefficient of every word $w$ of
length five or six. Independently, evaluate the finite cut sum
\eqref{eq:wordcut} for every such word. \Cref{app:tables} prints all nonzero
values of $c(w)$ through degree six (there are $30$ nonzero coefficients in
degree five and $28$ in degree six), and both computations give zero on
every word absent from the table. Thus the two polynomials agree in every
one of the $32$ and $64$ coordinates respectively. An equality of
noncommutative polynomials is fully determined by these word coefficients,
so the displayed formulas are identities in the free associative algebra,
hence in every associative specialization. Only the cut-formula side is
printed in \cref{app:tables}; the bracket-expansion side is reproduced,
coefficient by coefficient including the zeros, by the table generator and
by the verifiers described in \cref{sec:computation}, which additionally
recompute both sides from the differential recursions of
\cref{sec:integral}.
\end{proof}
These displays do not select a linearly independent basis of the free Lie
algebra; for instance $[X,[Y,[X,Y]]]=[Y,[X,[X,Y]]]$, since the difference is
$[[X,Y],[X,Y]]=0$ by \eqref{eq:derivation}. The coefficient of $X^4Y$ in
$Z_5$ is $-\tfrac1{720}=B_4^+/4!$, and that of $X^5Y$ in $Z_6$ is
$0=B_5^+/5!$, as the single-$Y$ Bernoulli coefficients of
\cref{cor:linearY} predict.

\section{Conjugation, the differential of the exponential, and Bernoulli series}\label{sec:differential}
This section opens a second route to the BCH series, independent of the Hopf
algebra argument. It also proves the Campbell identity and every
infinitesimal identity used later. Parameter differentiation of exponentials
is treated systematically in \cite{wilcox}; all formulas required here are
derived in full.

\subsection{Campbell's conjugation identity and the braiding formulas}
\begin{theorem}[Campbell--Hadamard identity]\label{thm:campbell}
Let either $X,Y$ have positive degree in $\widehat T$ with $s$ a formal or
scalar parameter, or let $X,Y$ lie in a Banach algebra with $s$ an
arbitrary scalar. Then
\begin{equation}\label{eq:campbell}
 e^{sX}Ye^{-sX}=e^{s\ad_X}Y=\sum_{n\geq0}\frac{s^n}{n!}\ad_X^nY.
\end{equation}
For invertible elements write $\Ad_AY=AYA^{-1}$; then
$\Ad_{e^X}=e^{\ad_X}$. The Banach-algebra series converges absolutely for
every $X,Y,s$, with
$\sum_n|s|^n\norm{\ad_X^nY}/n!\leq e^{2|s|\norm X}\norm Y$.
An entirely associative version is
\begin{equation}\label{eq:adbinomial}
 \ad_X^nY=\sum_{j=0}^{n}(-1)^j\binom nj X^{n-j}YX^j.
\end{equation}
\end{theorem}
\begin{proof}
Let $L_X$ and $R_X$ denote left and right multiplication by $X$. They
commute, since $L_XR_XA=XAX=R_XL_XA$, and $\ad_X=L_X-R_X$. The binomial
theorem for commuting operators gives
$\ad_X^n=\sum_j(-1)^j\binom njL_X^{n-j}R_X^j$, which applied to $Y$ is
\eqref{eq:adbinomial}.

Set $F(s)=e^{sX}Ye^{-sX}$. In a Banach algebra $F$ is differentiable with
$F'(s)=Xe^{sX}Ye^{-sX}-e^{sX}Ye^{-sX}X=\ad_XF(s)$ and $F(0)=Y$, by the
product rule and $\frac{\dd}{\dd s}e^{sX}=Xe^{sX}=e^{sX}X$. The map
$\ad_X$ is a bounded linear operator with $\norm{\ad_X}\leq2\norm X$, so the
linear differential equation $F'=\ad_XF$ has the unique solution
$F(s)=e^{s\ad_X}Y=\sum_ns^n\ad_X^nY/n!$; uniqueness holds because a
difference $G$ of two solutions satisfies $G(0)=0$ and
$\norm{G(s)}\leq2\norm X\bigl|\int_0^s\norm{G(\sigma)}\,\dd\sigma\bigr|$ on
every compact interval, hence vanishes by Gr\"onwall's inequality. The norm bound follows from $\norm{\ad_X^nY}\leq(2\norm X)^n\norm Y$.
Formally, both sides of \eqref{eq:campbell} are power series in $s$ with
coefficients in $\widehat T$; using \eqref{eq:adbinomial}, the coefficient
of $s^n$ on the left is
$\sum_{j}X^{n-j}Y(-X)^j/((n-j)!j!)=\frac1{n!}\sum_j\binom nj(-1)^jX^{n-j}YX^j
=\ad_X^nY/n!$, which is the coefficient on the right. (This direct comparison
also gives a second, purely algebraic proof in the Banach case.) Setting
$s=1$ gives $\Ad_{e^X}=e^{\ad_X}$.
\end{proof}

Conjugation $A\mapsto e^XAe^{-X}$ is an algebra automorphism; it is
continuous in the Banach case, and in the formal case it preserves the
filtration ($e^XF^m\widehat Te^{-X}\subseteq F^m\widehat T$, since $e^{\pm X}$
have constant term $1$) and is therefore continuous for the degree
topology. Hence it commutes with the exponential series. Therefore
\begin{align}
 e^Xe^Ye^{-X}&=\exp\bigl(e^XYe^{-X}\bigr)=\exp\bigl(e^{\ad_X}Y\bigr),\label{eq:conjexp}\\
 e^Xe^Y&=\exp\Bigl(\sum_{n\geq0}\frac{\ad_X^nY}{n!}\Bigr)e^X.\label{eq:braiding}
\end{align}
The second follows from the first by multiplying on the right by $e^X$. These
are the general braiding identities of the source page, including the
iterated-commutator notation $[X,Y]_n=\ad_X^nY$ there. Unlike a general BCH
series, the adjoint exponential in \eqref{eq:braiding} is an entire series
for bounded operators and requires no smallness.

\subsection{The differential of the exponential}
\begin{theorem}[Duhamel differential]\label{thm:dexp}
Let $A(t)$ be a differentiable curve in a Banach algebra, with $H=A'(t)$,
or a formal power series in $t$ with coefficients in $F^1\widehat T$. Then
\begin{align}
 (\dd\exp)_A(H):=\frac{\dd}{\dd t}e^{A(t)}
 &=\sum_{p,q\geq0}\frac{A^pHA^q}{(p+q+1)!}
 =\int_0^1e^{(1-s)A}He^{sA}\,\dd s,\label{eq:duhamel}\\
 e^{-A}\frac{\dd}{\dd t}e^{A}
 &=\int_0^1e^{-s\ad_A}H\,\dd s
 =\sum_{n\geq0}\frac{(-1)^n}{(n+1)!}\ad_A^nH=\phi(\ad_A)H,\label{eq:leftdexp}\\
 \Bigl(\frac{\dd}{\dd t}e^{A}\Bigr)e^{-A}
 &=\int_0^1e^{s\ad_A}H\,\dd s
 =\sum_{n\geq0}\frac{1}{(n+1)!}\ad_A^nH=\phi(-\ad_A)H.\label{eq:rightdexp}
\end{align}
All series converge absolutely for arbitrary bounded $A,H$.
\end{theorem}
\begin{proof}
The product rule, applied without moving any factor, gives
$\frac{\dd}{\dd t}A^n=\sum_{j=0}^{n-1}A^jHA^{n-1-j}$; in the Banach case the
exponential series may be differentiated term by term because
$\sum_n\norm{\frac{\dd}{\dd t}A^n}/n!\leq\sum_nn\norm A^{n-1}\norm H/n!$
converges uniformly on compact parameter sets, and in the formal case the
statement is coefficientwise. Collecting terms, the coefficient of $A^pHA^q$
in $\frac{\dd}{\dd t}e^A$ is $1/(p+q+1)!$, which is the first equality in
\eqref{eq:duhamel}. Expanding the two exponentials in the integral, the
coefficient of $A^pHA^q$ is
\[
 \frac1{p!q!}\int_0^1(1-s)^ps^q\,\dd s=\frac{1}{(p+q+1)!}.
\]
The beta integral is evaluated by induction on $p$: for $p=0$ it is
$1/(q+1)$, and integration by parts gives
$\int_0^1(1-s)^ps^q\dd s=\frac{p}{q+1}\int_0^1(1-s)^{p-1}s^{q+1}\dd s$, so
$\int_0^1(1-s)^ps^q\dd s=\frac{p!}{(q+1)(q+2)\cdots(q+p)}\cdot\frac1{q+p+1}
=\frac{p!q!}{(p+q+1)!}$. This proves \eqref{eq:duhamel}; interchanging sum
and integral is justified by the absolutely convergent majorant
$e^{(1-s)\norm A}\norm He^{s\norm A}$.

Multiplying \eqref{eq:duhamel} on the left by $e^{-A}$ gives
$\int_0^1e^{-sA}He^{sA}\dd s$, and \cref{thm:campbell} with $X=A$ and
parameter $-s$ turns the integrand into $e^{-s\ad_A}H$. Integrating the
absolutely convergent series $\sum_n(-s)^n\ad_A^nH/n!$ term by term gives
$\sum_n(-1)^n\ad_A^nH/(n+1)!$, which is \eqref{eq:leftdexp}. Multiplying
\eqref{eq:duhamel} on the right by $e^{-A}$ gives
$\int_0^1e^{(1-s)A}He^{-(1-s)A}\dd s=\int_0^1e^{u\ad_A}H\,\dd u$, and the same
integration gives \eqref{eq:rightdexp}.
\end{proof}
The quotient in $\phi(z)=(1-e^{-z})/z$ is notation for an entire function
with $\phi(0)=1$; \eqref{eq:leftdexp} does not require $\ad_A$ to be
invertible. The same statements hold intrinsically for finite-dimensional
Lie groups, where multiplication by $e^{-A}$ means left or right translation
of tangent vectors; an intrinsic derivation is given in \cref{sec:liegroups}.
In particular $(\dd\exp)_0=\id$.

\subsection{Bernoulli coefficients, including the convention at degree one}
The formal reciprocal of $\phi$ is
\begin{equation}\label{eq:betadef}
 \beta(z)=\frac{z}{1-e^{-z}}=\sum_{n\geq0}\bplus{n}z^n,\qquad
 \bplus{n}=\frac{B_n^+}{n!},
\end{equation}
which defines the Bernoulli numbers $B_n^+$ in the convention
$B_1^+=+\tfrac12$ used in the note of the source page. The other common
convention is generated by $\gamma(z)=z/(e^z-1)=\beta(-z)=\sum\bminus{n}z^n$,
with $\bminus{n}=(-1)^n\bplus{n}$ and $B_1^-=-\tfrac12$.

\begin{proposition}[Complete Bernoulli coefficients]\label{prop:bernoulli}
$\bplus{0}=1$, and for $n\geq1$ the coefficients are determined by the
recursion
\begin{equation}\label{eq:bernrec}
 \bplus{n}=-\sum_{j=1}^{n}\frac{(-1)^j}{(j+1)!}\,\bplus{n-j},
\end{equation}
and by the finite nonrecursive formula
\begin{equation}\label{eq:bernfinite}
 \bplus{n}=\sum_{k=1}^{n}(-1)^{n+k}
 \sum_{\substack{j_1+\cdots+j_k=n\\j_i\geq1}}\ \prod_{i=1}^{k}\frac1{(j_i+1)!}.
\end{equation}
Moreover
\begin{equation}\label{eq:coth}
 \beta(z)=\frac z2\Bigl(1+\coth\frac z2\Bigr)
 =1+\frac z2+\frac{z^2}{12}-\frac{z^4}{720}+\frac{z^6}{30240}
 -\frac{z^8}{1209600}+\cdots,
\end{equation}
so $\bplus{2m+1}=0$ for $m\geq1$; the even coefficients are
\begin{equation}\label{eq:bernzeta}
 \bplus{2m}=\frac{2(-1)^{m-1}\zeta(2m)}{(2\pi)^{2m}}\qquad(m\geq1),
\end{equation}
where $\zeta(2m)=\sum_{k\geq1}k^{-2m}$; and the Taylor series of $\beta$ has
radius of convergence exactly $2\pi$.
\end{proposition}
\begin{proof}
Write $\phi(z)=\sum_{j\geq0}\phi_jz^j$ with $\phi_j=(-1)^j/(j+1)!$, so
$\phi_0=1$. Comparing the coefficient of $z^n$, $n\geq1$, in
$\phi(z)\beta(z)=1$ gives $\sum_{j=0}^n\phi_j\bplus{n-j}=0$, which is
\eqref{eq:bernrec}. For \eqref{eq:bernfinite}, write $\phi=1+A$ with
$A(z)=\sum_{j\geq1}\phi_jz^j$ of positive order and expand
$\beta=(1+A)^{-1}=\sum_{k\geq0}(-A)^k$, a formally convergent geometric
series. The coefficient of $z^n$ in $(-A)^k$ is
$(-1)^k\sum_{j_1+\cdots+j_k=n}\prod_i\phi_{j_i}
=(-1)^k(-1)^n\sum\prod_i1/(j_i+1)!$, and summing over $k$ from $1$ to $n$
(no $k>n$ contributes because each $j_i\geq1$) gives \eqref{eq:bernfinite}.

Multiplying numerator and denominator of $\beta$ by $e^{z/2}$,
\[
 \beta(z)=\frac{ze^{z/2}}{e^{z/2}-e^{-z/2}}
 =\frac z2\cdot\frac{e^{z/2}+e^{-z/2}+(e^{z/2}-e^{-z/2})}{e^{z/2}-e^{-z/2}}
 =\frac z2\Bigl(\coth\frac z2+1\Bigr).
\]
The function $\frac z2\coth\frac z2=\beta(z)-\frac z2$ is even in $z$ and
regular at zero, so all its odd Taylor coefficients vanish; this proves
$\bplus{2m+1}=0$ for $m\geq1$. The displayed numerical coefficients follow
from \eqref{eq:bernrec}: $\bplus1=\tfrac12$, $\bplus2=\tfrac1{12}$,
$\bplus4=-\tfrac1{720}$, $\bplus6=\tfrac1{30240}$,
$\bplus8=-\tfrac1{1209600}$; for instance
$\bplus1=-\phi_1\bplus0=\tfrac12$ and
$\bplus2=-\bigl(\phi_1\bplus1+\phi_2\bplus0\bigr)
=-\bigl(-\tfrac12\cdot\tfrac12+\tfrac16\bigr)=\tfrac1{12}$.

For the even coefficients we derive the partial fraction expansion of
$\coth$ from an elementary Fourier series. Fix a real $z>0$ and let $g$ be
the $1$-periodic extension of $u\mapsto\cosh(z(u-\tfrac12))$ from $[0,1]$;
it is continuous, since $g(0)=g(1)=\cosh(z/2)$, and piecewise smooth. Its
Fourier coefficients are
\begin{align*}
 c_k&=\int_0^1\cosh\bigl(z(u-\tfrac12)\bigr)e^{-2\pi\ii ku}\,\dd u
 =\frac12e^{-z/2}\int_0^1e^{(z-2\pi\ii k)u}\dd u
  +\frac12e^{z/2}\int_0^1e^{-(z+2\pi\ii k)u}\dd u\\
 &=\frac{e^{z/2}-e^{-z/2}}{2(z-2\pi\ii k)}
  +\frac{e^{z/2}-e^{-z/2}}{2(z+2\pi\ii k)}
 =\frac{2z\sinh(z/2)}{z^2+4\pi^2k^2},
\end{align*}
using $e^{\pm2\pi\ii k}=1$. Since $c_k=O(k^{-2})$, the Fourier series
$\sum_{k\in\Z}c_ke^{2\pi\ii ku}$ converges absolutely and uniformly to a
continuous function with the same Fourier coefficients as $g$; by uniqueness
of Fourier coefficients for continuous functions, its sum is $g$. Evaluating
at $u=0$,
\[
 \cosh\frac z2=\frac{2\sinh(z/2)}{z}+\sum_{k\geq1}\frac{4z\sinh(z/2)}{z^2+4\pi^2k^2},
 \qquad\text{so}\qquad
 \coth\frac z2=\frac2z+4z\sum_{k\geq1}\frac1{z^2+4\pi^2k^2}.
\]
Both sides are meromorphic functions of $z\in\C$; the right side converges
normally on compact sets avoiding the points $\pm2\pi\ii k$, so the identity
proved for real $z>0$ extends to all $z$ away from the poles by the identity
theorem. Substituting into \eqref{eq:coth},
\begin{equation}\label{eq:partialfractions}
 \beta(z)=1+\frac z2+2z^2\sum_{k\geq1}\frac1{z^2+4\pi^2k^2}.
\end{equation}
For $|z|<2\pi$ expand each term geometrically,
$\frac{1}{z^2+4\pi^2k^2}=\sum_{m\geq0}\frac{(-1)^mz^{2m}}{(2\pi k)^{2m+2}}$;
the double series converges absolutely because
$\sum_k\sum_m|z|^{2m}/(2\pi k)^{2m+2}=\sum_k1/(4\pi^2k^2-|z|^2)<\infty$, so
the order of summation may be interchanged, giving
$\beta(z)=1+\frac z2+\sum_{m\geq0}\frac{2(-1)^m\zeta(2m+2)}{(2\pi)^{2m+2}}z^{2m+2}$,
which is \eqref{eq:bernzeta}. Finally, $1-e^{-z}$ vanishes exactly at
$2\pi\ii\Z$; at $z=0$ the zero of the denominator is cancelled by the
numerator, while at $\pm2\pi\ii$ the numerator is nonzero, so $\beta$ has
genuine poles there and nowhere closer to the origin. The Taylor radius is
therefore exactly $2\pi$. (Equivalently, \eqref{eq:bernzeta} gives
$|\bplus{2m}|^{1/2m}\to1/(2\pi)$ because $\zeta(2m)\to1$.)
\end{proof}
Functional notation such as $\beta(\ad_X)$ does not require $\ad_X$ to be
invertible. It means the formal series, or the analytic functional calculus in
a domain avoiding the poles; in a finite-dimensional space the latter is
defined whenever $\spec(\ad_X)$ avoids the nonzero points $2\pi\ii\ell$
(\cref{sec:regularX}).

\subsection{The function \texorpdfstring{$\psi$}{psi} and its complete expansion}
The other scalar function on the source page is
\begin{equation}\label{eq:psi}
 \psi(x)=\frac{x\log x}{x-1}=1-\sum_{n\geq1}\frac{(1-x)^n}{n(n+1)},
 \qquad|1-x|<1,\qquad \psi(1)=1,
\end{equation}
and, with the logarithm branch satisfying $\log(e^z)=z$ near $z=0$,
\begin{equation}\label{eq:psibeta}
 \psi(e^z)=\beta(z)=\sum_{n\geq0}B_n^+\frac{z^n}{n!}.
\end{equation}
\begin{proof}[Proof of \eqref{eq:psi} and \eqref{eq:psibeta}]
Put $u=1-x$ and use $\log(1-u)=-\sum_{j\geq1}u^j/j$ for $|u|<1$. Then
\[
 \psi(1-u)=\frac{(1-u)\log(1-u)}{-u}
 =(1-u)\sum_{j\geq1}\frac{u^{j-1}}{j}
 =\sum_{j\geq1}\frac{u^{j-1}}{j}-\sum_{j\geq1}\frac{u^{j}}{j}
 =1+\sum_{n\geq1}\Bigl(\frac1{n+1}-\frac1n\Bigr)u^n,
\]
and $\frac1{n+1}-\frac1n=-\frac1{n(n+1)}$. The same computation is valid as
an identity of formal power series in $u$. Substituting $x=e^z$ into
$x\log x/(x-1)$ gives $ze^z/(e^z-1)=z/(1-e^{-z})=\beta(z)$; as a formal
identity, substitute the positive-order series $u=1-e^z$ into
\eqref{eq:psi}.
\end{proof}
Equation \eqref{eq:psibeta} must not be read as a branch-independent
identity for arbitrary complex $z$; it is an identity of power series at
$z=0$, and analytically it holds wherever the logarithm branch used in
$\psi$ satisfies $\log(e^z)=z$.

\section{The integral formula and every order in either variable}\label{sec:integral}
\subsection{The differential equation and the Poincar\'e integral}
\begin{theorem}[Logarithmic differential equation]\label{thm:ode}
In the degree completion, and analytically for $X,Y$ in a sufficiently
small neighborhood of $(0,0)$ in a Banach algebra or a finite-dimensional
Lie algebra, the function $Z(t)=\BCH(X,tY)=\log(e^Xe^{tY})$ is
characterized by
\begin{equation}\label{eq:ode}
 Z'(t)=\beta(\ad_{Z(t)})Y,\qquad Z(0)=X.
\end{equation}
\end{theorem}
\begin{proof}
Let $G(t)=e^Xe^{tY}$. Its left logarithmic derivative is
$G^{-1}G'=e^{-tY}e^{-X}e^Xe^{tY}Y=Y$. Since $G(t)=e^{Z(t)}$,
\eqref{eq:leftdexp} gives $G^{-1}G'=\phi(\ad_{Z})Z'$, hence
$\phi(\ad_Z)Z'=Y$. Formally, $\ad_Z$ raises total degree in $X,Y$, so
$\beta(\ad_Z)=\sum_n\bplus{n}\ad_Z^n$ is a well-defined operator on
$\widehat T$ and $\beta(\ad_Z)\phi(\ad_Z)=\id$ because $\beta\phi=1$ as
scalar power series. Applying $\beta(\ad_Z)$ to both sides gives
\eqref{eq:ode}; $Z(0)=\log e^X=X$. Analytically, for $Z$ small
$\norm{\ad_Z}\leq2\norm Z<2\pi$ and $\phi(\ad_Z)$ is invertible with inverse
$\beta(\ad_Z)$ given by its convergent Taylor series, so the same
manipulation is valid; in a finite-dimensional Lie algebra the intrinsic
version of \eqref{eq:leftdexp} proved in \cref{sec:liegroups} is used.

Conversely, if $Z(t)$ solves \eqref{eq:ode}, then $e^{Z(t)}$ has left
logarithmic derivative $\phi(\ad_Z)Z'=\phi(\ad_Z)\beta(\ad_Z)Y=Y$ and
initial value $e^X$. The equation $G'=GY$, $G(0)=e^X$, has a unique
solution: formally, comparing coefficients of $t^n$ determines $G$
recursively; analytically, a difference $\Delta G$ of two solutions
satisfies $\Delta G'=\Delta G\,Y$ and $\Delta G(0)=0$, so
$\Delta G(t)=\Delta G(0)e^{tY}=0$ (differentiate $\Delta G(t)e^{-tY}$ to
see it is constant). Hence $e^{Z(t)}=e^Xe^{tY}$ and $Z(t)=\BCH(X,tY)$ by
uniqueness of the formal, or of the local analytic, logarithm.
\end{proof}
This differential equation is a useful independent mechanism for
constructing the BCH series: its right-hand side involves only iterated
Lie brackets, and no appeal to the Hopf algebra structure is needed for the
recursions derived from it below.

\begin{theorem}[Poincar\'e integral formula]\label{thm:poincare}
With the same interpretations,
\begin{equation}\label{eq:poincare}
 \boxed{\displaystyle
 \BCH(X,Y)=X+\int_0^1\psi\bigl(e^{\ad_X}e^{t\ad_Y}\bigr)Y\,\dd t,\qquad
 \psi(x)=\frac{x\log x}{x-1},}
\end{equation}
where $\psi$ of an operator $E$ near the identity is defined by the series
\eqref{eq:psi} in $I-E$. In the formal setting the integral means
termwise integration of the polynomial coefficients in $t$.
\end{theorem}
\begin{proof}
By \cref{thm:campbell} and the multiplicativity of $\Ad$,
\[
 e^{\ad_{Z(t)}}=\Ad_{e^{Z(t)}}=\Ad_{e^Xe^{tY}}
 =\Ad_{e^X}\Ad_{e^{tY}}=e^{\ad_X}e^{t\ad_Y}.
\]
Formally, $\ad_{Z(t)}$ has positive degree, so substituting it into the
scalar identity $\psi(e^z)=\beta(z)$ of \eqref{eq:psibeta} gives
$\psi(e^{\ad_{Z(t)}})=\beta(\ad_{Z(t)})$ as operators on $\widehat T$; here
$\psi(e^{\ad_Z})$ is computed by substituting the positive-degree operator
$I-e^{\ad_Z}$ into \eqref{eq:psi}. Analytically, for small $X,Y$ the
operator $E(t)=e^{\ad_X}e^{t\ad_Y}$ satisfies $\norm{E(t)-I}<1$, its
logarithm series converges to $\ad_{Z(t)}$ (the unique small logarithm),
and again $\psi(E(t))=\beta(\ad_{Z(t)})$. Hence \eqref{eq:ode} reads
$Z'(t)=\psi\bigl(e^{\ad_X}e^{t\ad_Y}\bigr)Y$, and integrating from $0$ to
$1$ with $Z(0)=X$ and $Z(1)=\BCH(X,Y)$ gives \eqref{eq:poincare}. In the
formal setting each fixed-degree coefficient of the integrand is a
polynomial in $t$, so integration is termwise and no interchange of limits
is involved.
\end{proof}
No invertibility of $\ad_Z$ is needed in this argument: it uses an identity
of endomorphisms, not recovery of $Z$ from $\ad_Z$. Likewise $\psi(E)$ near
$E=I$ is defined by its power series in $E-I$, not by multiplication by a
nonexistent inverse of $E-I$.

\subsection{An explicit all-bidegree formula without an unevaluated integral}
Put $A=\ad_X$ and $B=\ad_Y$ and recall that $Z_{p,q}$ denotes the component
of $Z$ with $p$ letters $X$ and $q$ letters $Y$.
\begin{theorem}[Operator-block expansion of every bidegree]\label{thm:operatorblock}
The full BCH series is
\begin{equation}\label{eq:operatorblock}
 \boxed{\displaystyle
 Z=X+Y+\sum_{k\geq1}\frac{(-1)^{k+1}}{k(k+1)}
 \sum_{\substack{r_i,s_i\geq0\\r_i+s_i>0\ (1\leq i\leq k)}}
 \frac{A^{r_1}B^{s_1}\cdots A^{r_k}B^{s_k}Y}
 {\bigl(1+\sum_is_i\bigr)\prod_ir_i!s_i!}.}
\end{equation}
Consequently $Z_{1,0}=X$, $Z_{0,1}=Y$, $Z_{p,0}=Z_{0,p}=0$ for $p\geq2$, and
for $q\geq1$, $p+q\geq2$,
\begin{equation}\label{eq:bidegree}
 \boxed{\displaystyle
 Z_{p,q}=\frac1q\sum_{k=1}^{p+q-1}\frac{(-1)^{k+1}}{k(k+1)}
 \sum_{\substack{r_i,s_i\geq0,\ r_i+s_i>0\\
                   \sum_ir_i=p,\ \sum_is_i=q-1}}
 \frac{A^{r_1}B^{s_1}\cdots A^{r_k}B^{s_k}Y}{\prod_ir_i!s_i!}.}
\end{equation}
At fixed total degree all sums are finite; every summand with $s_k>0$
vanishes, and so does every summand with $p=0$ and $k\geq1$.
\end{theorem}
\begin{proof}
Set $E(t)=e^Ae^{tB}-I$. Expanding the two exponentials,
\[
 E(t)=\sum_{\substack{r,s\geq0\\r+s>0}}\frac{t^s}{r!s!}A^rB^s,
\]
so $E(t)^k$ is the sum over $k$-tuples of nonempty blocks of
$\frac{t^{\sum s_i}}{\prod r_i!s_i!}A^{r_1}B^{s_1}\cdots A^{r_k}B^{s_k}$.
By \eqref{eq:psi} with $1-x$ replaced by $-E(t)$,
\[
 \psi(I+E(t))=I-\sum_{k\geq1}\frac{(-E(t))^k}{k(k+1)}
 =I+\sum_{k\geq1}\frac{(-1)^{k+1}}{k(k+1)}E(t)^k.
\]
Insert this into \eqref{eq:poincare}. The term $I$ integrates to $Y$. A
block term carries $t^{\sum_is_i}$, whose integral over $[0,1]$ is
$1/(1+\sum_is_i)$; since the operator is applied to a final letter $Y$, the
number of letters $Y$ in the result is $q=1+\sum_is_i$ and the number of
letters $X$ is $p=\sum_ir_i$. This proves \eqref{eq:operatorblock} and, on
collecting the terms of bidegree $(p,q)$, \eqref{eq:bidegree}; the upper
limit $k\leq p+q-1$ holds because every block has positive degree and the
total degree of the blocks is $p+q-1$. A summand with $s_k>0$ ends with
$BY=[Y,Y]=0$ and vanishes; if $p=0$ every block has $r_i=0$, so $s_k>0$ and
the summand vanishes, which gives $Z_{0,q}=0$ for $q\geq2$. Finally
$Z_{p,0}=0$ for $p\geq2$ because $\BCH(X,0)=X$ by \eqref{eq:unitinverse}.
\end{proof}
For example, for $(p,q)=(1,2)$ the only block tuples with
$\sum r_i=1$, $\sum s_i=1$, and $s_k=0$ are $(B)(A)$ with $k=2$, giving
$\frac{1}{2}\cdot\frac{-1}{6}BAY=-\frac1{12}[Y,[X,Y]]$, in agreement with
\eqref{eq:z3}. Unlike a finite display with an unspecified remainder,
\eqref{eq:bidegree} gives every term with three, four, or arbitrarily many
occurrences of $Y$.

\subsection{The part linear in \texorpdfstring{$Y$}{Y} and a triangular recursion}
\begin{corollary}[All single-$Y$ terms]\label{cor:linearY}
The complete part of $Z(X,Y)$ that is linear in $Y$ is
\begin{equation}\label{eq:linearY}
 \beta(\ad_X)Y=\frac{\ad_X}{1-e^{-\ad_X}}Y
 =\frac{\ad_X}{2}\Bigl(1+\coth\frac{\ad_X}{2}\Bigr)Y
 =Y+\tfrac12[X,Y]+\tfrac1{12}[X,[X,Y]]-\tfrac1{720}\ad_X^4Y
   +\tfrac1{30240}\ad_X^6Y-\cdots,
\end{equation}
so that $Z(X,Y)=X+\beta(\ad_X)Y+O(Y^2)$, and the coefficient of $\ad_X^nY$
is $B_n^+/n!$.
\end{corollary}
\begin{proof}
Replace $Y$ by $\eps Y$ in \eqref{eq:ode}: $\frac{\dd}{\dd t}Z(X,t\eps Y)
=\eps\beta(\ad_{Z})Y$. The coefficient of $\eps^1$ in $Z(X,\eps Y)$ is
obtained by integrating the coefficient of $\eps^1$ on the right from
$t=0$ to $1$; since $Z=X+O(\eps)$, that coefficient is $\beta(\ad_X)Y$.
Alternatively, the terms with $q=1$ in \eqref{eq:bidegree} are those with
all $s_i=0$, that is $Z_{p,1}=\sum_{k}\frac{(-1)^{k+1}}{k(k+1)}
\sum_{r_1+\cdots+r_k=p,\ r_i\geq1}\frac{A^pY}{\prod r_i!}$, which is the
coefficient of $z^p$ in $\psi(e^z)=\beta(z)$ expanded by \eqref{eq:psi}.
The hyperbolic form is \eqref{eq:coth}.
\end{proof}
This proves both versions of the source page's linear-in-$Y$ formula. The
quotients and the product with $\coth$ denote the regular function $\beta$;
neither expression calls for an inverse of $\ad_X$ at zero. The remainder
$O(Y^2)$ means membership in the subspace of terms with at least two
occurrences of $Y$, not an unspecified bound on $\norm X$. The linear
formula by itself is not an exact BCH formula unless a further relation
annihilates all higher-$Y$ terms, as in \cref{thm:eigen}.

The remainder can be resolved at every order. Write
\begin{equation}\label{eq:Hqdef}
 Z(X,\eps Y)=\sum_{q\geq0}\eps^qH_q(X;Y),\qquad H_0=X,\qquad H_1=\beta(A)Y,
\end{equation}
where $H_q$ is a formal series in $X$ homogeneous of degree $q$ in $Y$.
\begin{proposition}[Triangular recursion in the number of $Y$'s]\label{prop:Hqrec}
For every $q\geq1$,
\begin{equation}\label{eq:Hqrec}
 \boxed{\displaystyle
 qH_q=\sum_{m\geq0}\bplus{m}
 \sum_{\substack{j_1,\ldots,j_m\geq0\\j_1+\cdots+j_m=q-1}}
 \ad_{H_{j_1}}\cdots\ad_{H_{j_m}}Y,}
\end{equation}
where for $m=0$ the inner term is $Y$ if $q=1$ and zero otherwise. Every
index $j_i$ is less than $q$, and in every fixed total degree the sum is
finite. In particular the complete quadratic part is
\begin{equation}\label{eq:H2}
 H_2=\frac12\sum_{m\geq1}\bplus{m}\sum_{r=0}^{m-1}A^r\bigl[H_1,A^{m-1-r}Y\bigr].
\end{equation}
The nonrecursive alternative is $H_q=\sum_{p\geq0}Z_{p,q}$ with
\eqref{eq:bidegree}.
\end{proposition}
\begin{proof}
Differentiate \eqref{eq:Hqdef} with respect to $\eps$. By
\cref{thm:ode} applied to $Z(X,\eps Y)$ as a function of $\eps$,
$\partial_\eps Z=\beta(\ad_Z)Y$. Expand
$\ad_Z=A+\sum_{k\geq1}\eps^k\ad_{H_k}$ and each power
$\ad_Z^m=\sum_{j_1,\ldots,j_m\geq0}\eps^{j_1+\cdots+j_m}
\ad_{H_{j_1}}\cdots\ad_{H_{j_m}}$ without changing the order of the
factors. The coefficient of $\eps^{q-1}$ on the right is the double sum
in \eqref{eq:Hqrec}, and on the left it is $qH_q$. An index $j_i=q-1$
forces all others to be zero, so only $H_j$ with $j\leq q-1$ occur. Although
there are arbitrarily many factors $\ad_{H_0}=A$, each raises the total
degree, so at fixed total degree only finitely many $m$ contribute. For
$q=2$ exactly one index equals $1$ and the rest are $0$; summing over its
position $r$ gives \eqref{eq:H2}. The last assertion is
\cref{thm:operatorblock} summed over $p$.
\end{proof}

\subsection{A bivariate generating function for the complete quadratic term}
\begin{theorem}[Quadratic kernel]\label{thm:kernel}
Let
\begin{equation}\label{eq:K}
 K(u,v)=\frac{\beta(u)}{2u}\bigl(\beta(u+v)-\beta(v)\bigr)
 =\sum_{p,q\geq0}k_{pq}u^pv^q,
\end{equation}
where the quotient is the regular divided difference
$(\beta(u+v)-\beta(v))/u=\sum_m\bplus{m}\sum_{r=0}^{m-1}(u+v)^rv^{m-1-r}$.
Then
\begin{equation}\label{eq:H2kernel}
 \boxed{\displaystyle
 H_2=\sum_{p,q\geq0}k_{pq}\,[A^pY,A^qY],\qquad
 k_{pq}=\frac12\sum_{a=0}^{p}\bplus{a}\,\bplus{p+q+1-a}
 \binom{p+q+1-a}{p+1-a}.}
\end{equation}
Thus every coefficient of the quadratic part of $Z$ is an explicit finite
sum of products of two Bernoulli coefficients.
\end{theorem}
\begin{proof}
Define a linear map $\Phi$ from formal power series in two commuting
variables $u,v$ to $\widehat T$ by $\Phi(u^av^b)=[A^aY,A^bY]$. By the
derivation rule \eqref{eq:derivation},
$A[A^aY,A^bY]=[A^{a+1}Y,A^bY]+[A^aY,A^{b+1}Y]$, that is,
$A\,\Phi(P)=\Phi\bigl((u+v)P\bigr)$ for every series $P$, and therefore
$A^r\Phi(P)=\Phi((u+v)^rP)$. Since $H_1=\beta(A)Y=\sum_a\bplus aA^aY$,
\[
 [H_1,A^{m-1-r}Y]=\sum_a\bplus a[A^aY,A^{m-1-r}Y]=\Phi\bigl(\beta(u)v^{m-1-r}\bigr).
\]
Hence \eqref{eq:H2} becomes
\[
 H_2=\frac12\sum_{m\geq1}\bplus m\sum_{r=0}^{m-1}
 \Phi\bigl((u+v)^r\beta(u)v^{m-1-r}\bigr)
 =\frac12\,\Phi\Bigl(\beta(u)\sum_{m\geq1}\bplus m\frac{(u+v)^m-v^m}{u}\Bigr),
\]
using the finite geometric identity
$\sum_{r=0}^{m-1}(u+v)^rv^{m-1-r}=\bigl((u+v)^m-v^m\bigr)/u$. The sum over
$m\geq1$ equals $(\beta(u+v)-\beta(v))/u$, since the $m=0$ term vanishes.
This proves $H_2=\Phi(K)$, which is the first part of \eqref{eq:H2kernel}.
For the coefficients, expand
$(u+v)^m-v^m=\sum_{j\geq1}\binom mju^jv^{m-j}$, so
\[
 K(u,v)=\frac12\sum_{a\geq0}\bplus au^a\sum_{m\geq1}\bplus m
 \sum_{j=1}^{m}\binom mju^{j-1}v^{m-j}.
\]
The coefficient of $u^pv^q$ collects the terms with $a+j-1=p$ and $m-j=q$,
that is $j=p+1-a$ and $m=p+q+1-a$, with $0\leq a\leq p$; this is the stated
finite sum.
\end{proof}
Because brackets are antisymmetric, $K$ may be replaced by its
antisymmetrization $\frac12(K(u,v)-K(v,u))$ without changing $\Phi(K)$. For
orientation, the first terms grouped by the number of $X$'s are
\begin{equation}\label{eq:H2first}
 H_2=-\frac1{12}[Y,AY]-\frac1{24}[Y,A^2Y]-\frac1{180}[Y,A^3Y]
 -\frac1{120}[AY,A^2Y]+O_X(4),
\end{equation}
where $O_X(4)$ means at least four $X$'s and exactly two $Y$'s. For
example, $k_{01}=\frac12\bplus0\bplus2\binom21=\frac1{12}$ and
$k_{10}=\frac12\bigl(\bplus0\bplus2\binom22+\bplus1\bplus1\binom11\bigr)
=\frac12\bigl(\frac1{12}+\frac14\bigr)=\frac16$, so the coefficient of
$[Y,AY]$ is $k_{01}-k_{10}=-\frac1{12}$, in agreement with the term
$\frac1{12}[Y,[Y,X]]$ of \eqref{eq:z3}. The kernel \eqref{eq:K}, not this
short display, is the full expansion.

\subsection{Analytic meaning when \texorpdfstring{$X$}{X} is not close to zero}\label{sec:regularX}
The recursion \eqref{eq:Hqrec} is a formal identity in the total-degree
completion. When $X$ is a fixed element of a finite-dimensional Lie algebra
$\g$ that is not small, the Bernoulli series $\beta(\ad_X)$ may diverge
even though the local expansion of $\log(e^Xe^{\eps Y})$ around $\eps=0$
exists. The following finite-at-each-order replacement gives the actual
local analytic expansion.

Let $A=\ad_X$, complexified if $\g$ is real, and assume the
\emph{regularity condition}
\begin{equation}\label{eq:regular}
 \spec(A)\cap\bigl(2\pi\ii\Z\setminus\{0\}\bigr)=\varnothing.
\end{equation}
By triangularization the eigenvalues of $\phi(A)$ are the values
$\phi(\lambda)$ at the eigenvalues $\lambda$ of $A$, and $\phi$ vanishes
exactly at $2\pi\ii\Z\setminus\{0\}$; so \eqref{eq:regular} holds if and
only if $\phi(A)$ is invertible (see \cref{prop:coframe} for the detailed
argument). Then $\beta(A)$ \emph{means} the inverse $\phi(A)^{-1}$, which
by holomorphic functional calculus equals
$\frac1{2\pi\ii}\oint_\Gamma\beta(z)(zI-A)^{-1}\dd z$ for any finite union
$\Gamma$ of positively oriented contours enclosing $\spec(A)$ and avoiding
the poles of $\beta$; it is not necessarily the Taylor series at zero.
Since $(\dd\exp)_X(H)=e^X\phi(A)H$ by \eqref{eq:leftdexp} (in a Lie group,
by the intrinsic version in \cref{sec:liegroups}), the differential of the
exponential at $X$ is invertible, and the inverse function theorem provides
an exponential chart centered at $X$; thus $\log(e^Xe^{\eps Y})$ has a
unique analytic branch $Z(\eps)$ with $Z(0)=X$ for $|\eps|$ small, and it
satisfies \eqref{eq:ode} in the form $Z'(\eps)=\phi(\ad_{Z(\eps)})^{-1}Y$.

For operators $E_1,\ldots,E_m$ on $\g_\C$ define
\begin{equation}\label{eq:frechet}
 T_m(A)[E_1,\ldots,E_m]
 =\frac1{2\pi\ii}\oint_\Gamma\beta(z)\,R(z)E_1R(z)E_2\cdots E_mR(z)\,\dd z,
 \qquad R(z)=(zI-A)^{-1}.
\end{equation}
The symmetrization $\sum_{\pi\in S_m}T_m(A)[E_{\pi(1)},\ldots,E_{\pi(m)}]$
is the $m$th Fr\'echet derivative $D^m\beta(A)[E_1,\ldots,E_m]$ of the
holomorphic matrix function $\beta$, but the unsymmetrized form is what the
recursion needs.
\begin{theorem}[Local expansion at a regular base point]\label{thm:Hqanalytic}
Under \eqref{eq:regular}, $Z(\eps)=\sum_{q\geq0}\eps^qH_q$ converges for
$|\eps|$ small, with $H_0=X$, $H_1=\beta(A)Y$, and for $q\geq2$
\begin{equation}\label{eq:Hqanalytic}
 \boxed{\displaystyle
 qH_q=\sum_{m=1}^{q-1}\ \sum_{\substack{j_1,\ldots,j_m\geq1\\j_1+\cdots+j_m=q-1}}
 T_m(A)\bigl[\ad_{H_{j_1}},\ldots,\ad_{H_{j_m}}\bigr]Y.}
\end{equation}
Both sums are finite, and only earlier $H_j$ occur. In particular
$2H_2=T_1(A)[\ad_{\beta(A)Y}]Y$. For a real Lie algebra all $H_q$ are real.
\end{theorem}
\begin{proof}
For an operator $E$ with $\norm E\,\sup_{z\in\Gamma}\norm{R(z)}<1$, the
resolvent of $A+E$ exists on $\Gamma$ and is given by the Neumann series
\[
 (zI-A-E)^{-1}=\bigl(R(z)^{-1}(I-R(z)E)\bigr)^{-1}
 =\sum_{k\geq0}\bigl(R(z)E\bigr)^kR(z),
\]
converging uniformly on $\Gamma$. For such $E$ the spectrum of $A+E$ lies
inside $\Gamma$ (a point $z\in\Gamma$ or outside is in the resolvent set by
the same Neumann series, and the interior of $\Gamma$ contains
$\spec(A)$ together with a neighborhood on which the estimate persists), so
by the holomorphic functional calculus
\[
 \beta(A+E)=\frac1{2\pi\ii}\oint_\Gamma\beta(z)(zI-A-E)^{-1}\dd z
 =\sum_{k\geq0}\frac1{2\pi\ii}\oint_\Gamma\beta(z)R(z)\bigl(ER(z)\bigr)^k\dd z,
\]
and $\beta(A+E)=\phi(A+E)^{-1}$ because $\phi\beta=1$ and functional
calculus is multiplicative. Now take $E=E(\eps)=\sum_{j\geq1}\eps^j\ad_{H_j}$,
where the $H_j$ are the Taylor coefficients of the analytic branch
$Z(\eps)$; for $|\eps|$ small the norm condition holds. Substituting into
the uniformly convergent double series and collecting the coefficient of
$\eps^{q-1}$ gives
\[
 \coeff{\eps^{q-1}}\beta(\ad_{Z(\eps)})
 =\sum_{m\geq0}\ \sum_{\substack{j_1,\ldots,j_m\geq1\\j_1+\cdots+j_m=q-1}}
 T_m(A)\bigl[\ad_{H_{j_1}},\ldots,\ad_{H_{j_m}}\bigr],
\]
where the $m=0$ term is $\beta(A)$ when $q=1$ and zero otherwise, and
$m\leq q-1$ because each $j_i\geq1$. The differential equation
$Z'(\eps)=\beta(\ad_{Z(\eps)})Y$ holds for the analytic branch by the
argument of \cref{thm:ode} (the left logarithmic derivative of
$e^Xe^{\eps Y}$ is $Y$, and $\phi(\ad_{Z(\eps)})$ is invertible for small
$\eps$ by continuity of the spectrum), so comparing Taylor coefficients
gives $qH_q=\coeff{\eps^{q-1}}\beta(\ad_{Z(\eps)})Y$, which is
\eqref{eq:Hqanalytic}. Convergence of the Taylor series of $Z(\eps)$
follows from analyticity of the local inverse of the exponential map. For a
real Lie algebra, complex conjugation on $\g_\C$ commutes with $\beta(\ad_Z)$
for real $Z$ (the contour may be chosen symmetric and $\beta$ has real
Taylor coefficients), so the uniquely determined coefficients $H_q$ are
real.
\end{proof}
This separates the all-orders local expansion around a regular base point
$X$ from the total-degree expansion around $(0,0)$: in the former one must
not evaluate the infinite Bernoulli sum in \eqref{eq:Hqrec} literally when
that series diverges.

\subsection{Two finite all-order recursions by total degree}\label{sec:homrec}
There is also a recursive proof of Lie-series existence that uses neither
the primitive-element theorem nor the Poincar\'e integral. Let
$G(t)=e^{tX}e^{tY}$, $Z(t)=\log G(t)=\sum_{n\geq1}t^nZ_n$, and
\begin{equation}\label{eq:Vj}
 V(t)=X+e^{t\ad_X}Y=\sum_{j\geq0}t^jV_j,\qquad
 V_0=X+Y,\qquad V_j=\frac{\ad_X^jY}{j!}\ (j\geq1).
\end{equation}
\begin{theorem}[Total-degree Bernoulli recursion]\label{thm:homrec}
For $n\geq1$,
\begin{equation}\label{eq:homrec}
 \boxed{\displaystyle
 nZ_n=\sum_{m=0}^{n-1}\bminus{m}
 \sum_{\substack{k_0\geq0,\ k_1,\ldots,k_m\geq1\\
                   k_0+k_1+\cdots+k_m=n-1}}
 \ad_{Z_{k_1}}\cdots\ad_{Z_{k_m}}V_{k_0},}
\end{equation}
where for $m=0$ the inner expression means $V_{n-1}$. Every $Z_{k_i}$ on
the right has index less than $n$; hence $Z_1=X+Y$ and, by induction, every
$Z_n$ is a Lie polynomial with rational coefficients. Equivalently, with
$S=X+Y$ and $D=X-Y$,
\begin{equation}\label{eq:compactode}
 Z'(t)=S+\tfrac12[D,Z]+\sum_{p\geq1}\bplus{2p}\ad_Z^{2p}S,
\end{equation}
so that for $n\geq2$
\begin{equation}\label{eq:compactrec}
 nZ_n=\tfrac12[D,Z_{n-1}]
 +\sum_{p=1}^{\lfloor(n-1)/2\rfloor}\bplus{2p}
 \sum_{\substack{i_1+\cdots+i_{2p}=n-1\\i_j\geq1}}
 \ad_{Z_{i_1}}\cdots\ad_{Z_{i_{2p}}}S.
\end{equation}
\end{theorem}
\begin{proof}
The right logarithmic derivative of $G$ is
$G'G^{-1}=Xe^{tX}e^{tY}e^{-tY}e^{-tX}+e^{tX}Ye^{tY}e^{-tY}e^{-tX}
=X+e^{tX}Ye^{-tX}=X+e^{t\ad_X}Y=V(t)$ by \cref{thm:campbell}. On the other
hand \eqref{eq:rightdexp} gives $G'G^{-1}=\phi(-\ad_Z)Z'$. Formally,
$\ad_{Z(t)}$ has positive degree in $t$, so $\gamma(\ad_Z)=\beta(-\ad_Z)
=\sum_m\bminus m\ad_Z^m$ is a well-defined operator inverting
$\phi(-\ad_Z)$, and
\begin{equation}\label{eq:homode}
 Z'(t)=\gamma(\ad_{Z(t)})V(t).
\end{equation}
Expand $\ad_Z^m=\sum t^{k_1+\cdots+k_m}\ad_{Z_{k_1}}\cdots\ad_{Z_{k_m}}$
over $k_i\geq1$, multiply by $V(t)$, and compare the coefficients of
$t^{n-1}$: the left side gives $nZ_n$ and the right side gives
\eqref{eq:homrec}. All indices $k_i$ are at least $1$ and sum to at most
$n-1$, so each is less than $n$. Starting from $Z_1=V_0=X+Y$, induction
shows that every $Z_n$ is a rational Lie polynomial. The series so
constructed satisfies \eqref{eq:homode}, hence $e^{Z(t)}$ has right
logarithmic derivative $V(t)$ and value $1$ at $t=0$, so it equals
$G(t)$ by uniqueness for $G'=VG$ (compare coefficients of $t$); thus it
coincides with the logarithmic construction \eqref{eq:Zdef} at $t=1$.

For the compact form, note that $e^{\ad_Z}=\Ad_{e^Z}=\Ad_{e^{tX}}\Ad_{e^{tY}}
=e^{t\ad_X}e^{t\ad_Y}$ and $e^{t\ad_Y}Y=Y$, so $e^{\ad_Z}Y=e^{t\ad_X}Y$ and
$V=X+e^{\ad_Z}Y$. Since $\gamma(z)e^z=ze^z/(e^z-1)=\beta(z)$,
\eqref{eq:homode} becomes $Z'=\gamma(\ad_Z)X+\beta(\ad_Z)Y$. The two scalar
functions satisfy $\gamma(z)+\beta(z)=2\sum_{p\geq0}\bplus{2p}z^{2p}$ (their
odd parts cancel because $\gamma(z)=\beta(-z)$) and
$\beta(z)-\gamma(z)=z$ (their odd part is $\bplus1z=z/2$ each, with
opposite signs, and $\bplus{2p+1}=0$ for $p\geq1$). Writing
$X=\frac12(S+D)$ and $Y=\frac12(S-D)$,
\[
 Z'=\tfrac12(\gamma+\beta)(\ad_Z)S+\tfrac12(\gamma-\beta)(\ad_Z)D
 =\sum_{p\geq0}\bplus{2p}\ad_Z^{2p}S-\tfrac12\ad_ZD,
\]
which is \eqref{eq:compactode}. Comparing coefficients of $t^{n-1}$ gives
\eqref{eq:compactrec}.
\end{proof}
The recursion \eqref{eq:homrec} gives $2Z_2=V_1+\bminus1[Z_1,V_0]
=[X,Y]-\tfrac12[X+Y,X+Y]=[X,Y]$ and
\[
 3Z_3=V_2+\bminus1\bigl([Z_1,V_1]+[Z_2,V_0]\bigr)+\bminus2[Z_1,[Z_1,V_0]]
 =\tfrac12[X,[X,Y]]-\tfrac12[X+Y,[X,Y]]+\tfrac14[X+Y,[X,Y]]
 =\tfrac14[X,[X,Y]]-\tfrac14[Y,[X,Y]],
\]
which is \eqref{eq:z3}. The compact form \eqref{eq:compactrec} with
$C=[X,Y]$ gives $2Z_2=\tfrac12[D,S]=C$, $3Z_3=\tfrac14[D,C]$ (since
$\ad_S^2S=0$), and
\[
 4Z_4=\tfrac12[D,Z_3]+\bplus2[S,[Z_2,S]]
 =\tfrac1{24}[D,[D,C]]-\tfrac1{24}[S,[S,C]]
 =-\tfrac16[Y,[X,C]],
\]
using $[X,[Y,C]]=[Y,[X,C]]$; this reproduces \eqref{eq:z4}. Eichler's
associativity-based argument \cite{eichler} proves the same structural
theorem by a different route; no claim that his recursion coincides with
\eqref{eq:homrec} is needed here.

\section{Convergence, remainder estimates, and logarithm obstructions}\label{sec:convergence}
The formal theorem is unconditional, but its numerical evaluation is not.
Four different questions must be kept apart: convergence of the Mercator
series $\sum_k(-1)^{k-1}(e^Xe^Y-I)^k/k$; absolute convergence of its
expansion into individual associative words or into the individual
commutator summands of Dynkin's formula; convergence of the homogeneous Lie
series $\sum_nZ_n$; and existence of some logarithm of $e^Xe^Y$. These are
different summation procedures, and outside a common domain of absolute
convergence they can behave differently; examples and further discussion
appear in \cite{biagi}. This section proves elementary quantitative
statements sufficient for every convergence claim on the source page,
without claiming optimal constants.

\subsection{Absolute convergence in a Banach algebra}
Let $\cA$ be a real or complex unital Banach algebra with a submultiplicative
norm, $\norm{AB}\leq\norm A\norm B$. Unless said otherwise we do not assume
$\norm I=1$: every estimate below involves only products of positive-length
words, so the norm of the identity never enters. In particular the
Hilbert--Schmidt norm on $d\times d$ matrices is admissible: by
Cauchy--Schwarz,
\[
 \norm{AB}_{\mathrm{HS}}^2=\sum_{i,j}\Bigl|\sum_kA_{ik}B_{kj}\Bigr|^2
 \leq\sum_{i,j}\Bigl(\sum_k|A_{ik}|^2\Bigr)\Bigl(\sum_k|B_{kj}|^2\Bigr)
 =\norm A_{\mathrm{HS}}^2\norm B_{\mathrm{HS}}^2,
\]
although $\norm I_{\mathrm{HS}}=\sqrt d$. Throughout, $s=\norm X+\norm Y$.

\begin{theorem}[Two convergence domains and explicit remainders]\label{thm:convergence}
\begin{enumerate}[label=(\roman*)]
\item If $s<\log2$, the associative block expansion \eqref{eq:blocks} is
absolutely convergent, the homogeneous BCH series $\sum_nZ_n(X,Y)$ converges
absolutely, and its sum $Z$ satisfies $e^Z=e^Xe^Y$. More precisely, for
every $r>0$ with $rs<\log2$,
\begin{equation}\label{eq:majorant}
\begin{gathered}
 \sum_{k\geq1}\frac1k\sum_{\substack{r_i+s_i>0}}
 \frac{\norm{X^{r_1}Y^{s_1}\cdots X^{r_k}Y^{s_k}}}{\prod_ir_i!s_i!}\,r^{\sum(r_i+s_i)}
 \leq M_s(r):=-\log\bigl(2-e^{rs}\bigr),\\
 \sum_{n\geq1}r^n\norm{Z_n(X,Y)}\leq M_s(r).
\end{gathered}
\end{equation}
\item If $1<r<\log2/s$ and $Z^{[N]}=\sum_{n=1}^NZ_n$, then
\begin{align}
 \norm{Z-Z^{[N]}}&\leq r^{-(N+1)}M_s(r),\label{eq:remainder}\\
 \norm{e^{Z^{[N]}}-e^Xe^Y}&\leq c_{\cA}\,e^{M_s(1)}\,r^{-(N+1)}M_s(r),
 \label{eq:expremainder}
\end{align}
where $c_{\cA}=\norm I^2$; thus $c_{\cA}=1$ for a norm with $\norm I=1$.
\item If $s<\tfrac12\log2$, the individual right-nested commutator summands
of Dynkin's formula \eqref{eq:dynkin} are absolutely summable, with
\begin{equation}\label{eq:dynkinradius}
 \sum_{\text{all Dynkin summands}}\norm{\text{summand}}
 \leq M_D(s):=-\tfrac12\log\bigl(2-e^{2s}\bigr),
\end{equation}
and if $Rs<\tfrac12\log2$ with $R>1$ the tail after degree $N$ of this
absolute sum is at most $R^{-(N+1)}M_D(Rs)$.
\item For a finite-dimensional normed Lie algebra $\g$ with
$\norm{[u,v]}\leq\kappa\norm u\norm v$, Dynkin's Lie series converges
absolutely whenever $\kappa(\norm X+\norm Y)<\log2$, with absolute sum at
most $\kappa^{-1}\bigl[-\log(2-e^{\kappa s})\bigr]$.
\end{enumerate}
\end{theorem}
\begin{proof}
(i) Expanding into nonempty words,
\[
 \norm{e^Xe^Y-I}\leq\sum_{r+s>0}\frac{\norm X^r\norm Y^s}{r!s!}=e^{s}-1,
\]
where a block with $r=0$ or $s=0$ is a pure power and contributes no factor
$\norm I$. Replacing each term of the block expansion of $(e^Xe^Y-I)^k$ by
the product of the norms of its blocks bounds the sum of the norms of all
terms of the $k$th power by $(e^s-1)^k$. If $s<\log2$ then $e^s-1<1$, and
\[
 \sum_{k\geq1}\frac{(e^s-1)^k}{k}=-\log\bigl(1-(e^s-1)\bigr)=-\log(2-e^s).
\]
Replacing $X,Y$ by $rX,rY$ multiplies a term of total degree $n$ by $r^n$
and gives the first inequality in \eqref{eq:majorant}. Absolute convergence
of the expanded series permits regrouping by total degree; the group of
degree $n$ is $Z_n(X,Y)$ by \cref{thm:wordcut} (the formal identity is a
statement about which word coefficients occur), so
$\sum_nr^n\norm{Z_n}\leq M_s(r)$, and with $r=1$ the homogeneous series
converges absolutely. Its sum is the Mercator logarithm $\log(I+U)$,
$U=e^Xe^Y-I$, $\norm U<1$. The scalar identity $\exp(\log(1+u))=1+u$ holds
for $|u|<1$ as an identity of absolutely convergent series, and substituting
$u\mapsto U$ is legitimate in the closed commutative subalgebra generated by
$U$, exactly as in \cref{lem:explog}. Hence $e^Z=I+U=e^Xe^Y$.

(ii) For $n>N$, $r^n\geq r^{N+1}$, so
$\sum_{n>N}\norm{Z_n}\leq r^{-(N+1)}\sum_{n>N}r^n\norm{Z_n}\leq r^{-(N+1)}M_s(r)$,
which is \eqref{eq:remainder}. For the exponentials, differentiating
$u\mapsto e^{(1-u)A}e^{uB}$ gives
\begin{equation}\label{eq:expdifference}
 e^B-e^A=\int_0^1e^{(1-u)A}(B-A)e^{uB}\,\dd u,
\end{equation}
since the derivative of the integrand's antiderivative
$e^{(1-u)A}e^{uB}$ is $e^{(1-u)A}(-A+B)e^{uB}$ (the factor $A$ commutes with
$e^{(1-u)A}$ and $B$ with $e^{uB}$). Submultiplicativity gives $\norm I=\norm{I^2}\leq\norm I^2$, so
$\norm I\geq1$, and therefore
$\norm{e^C}\leq\norm I+\sum_{n\geq1}\norm C^n/n!\leq\norm I\,e^{\norm C}$.
Consequently
\[
 \norm{e^B-e^A}\leq\norm I^2\int_0^1e^{(1-u)\norm A}e^{u\norm B}\dd u\,\norm{B-A}
 \leq\norm I^2e^{\max(\norm A,\norm B)}\norm{B-A}.
\]
Apply this with $A=Z^{[N]}$ and $B=Z$: both have norm at most
$\sum_n\norm{Z_n}\leq M_s(1)$, and $\norm{Z-Z^{[N]}}$ is bounded by
\eqref{eq:remainder}. Since $e^Z=e^Xe^Y$ by (i), this proves
\eqref{eq:expremainder}.

(iii) For a right-nested bracket of length $n$, iterating
$\norm{[U,V]}\leq2\norm U\norm V$ gives
$\norm{R(a_1\cdots a_n)}\leq2^{n-1}\prod_j\norm{a_j}$. In \eqref{eq:dynkin}
discard the factor $1/d\leq1$. The sum of the norms of the degree-$n$
summands is then at most $2^{n-1}$ times the sum of the norms of the
corresponding associative block terms with $X,Y$ replaced by $2X,2Y$, divided
by $2^n$; hence the total is at most
$\frac12\sum_{k\geq1}(e^{2s}-1)^k/k=-\frac12\log(2-e^{2s})$, finite when
$2s<\log2$. The tail bound follows exactly as in (ii) by applying the
estimate to $RX,RY$.

(iv) Replace $2^{n-1}$ by $\kappa^{n-1}$ in (iii); the norm bound on a
right-nested bracket of length $n$ becomes $\kappa^{n-1}\prod_j\norm{a_j}$,
and the same computation gives the bound $\kappa^{-1}[-\log(2-e^{\kappa s})]$
under $\kappa s<\log2$. Every finite-dimensional Lie algebra admits such a
$\kappa$, because a bilinear map is bounded on the product of the unit
spheres, so this proves local convergence without a matrix representation.
If the bracket is identically zero, BCH terminates at $X+Y$ and no
positive $\kappa$ is needed.
\end{proof}

The two radii $\log2$ and $\tfrac12\log2$ concern different ways of summing
expressions. The larger applies to the series of already-collected
homogeneous Lie polynomials and to the associative word expansion; the
smaller controls all the unreduced nested-commutator terms of Dynkin's
multiple sum separately, and it is the bound stated in the source page's
convergence discussion. Neither is claimed to be optimal.

\subsection{The Mercator tail and the role of \texorpdfstring{$\norm Z<\log2$}{||Z||<log 2}}
For a direct truncation of the Mercator series, put $q=e^s-1<1$. Bounding
$1/k\leq1/(K+1)$ in the tail gives
\begin{equation}\label{eq:mercatortail}
 \Bigl\|\log(e^Xe^Y)-\sum_{k=1}^K\frac{(-1)^{k-1}}{k}(e^Xe^Y-I)^k\Bigr\|
 \leq\sum_{k>K}\frac{q^k}{k}\leq\frac{q^{K+1}}{(K+1)(1-q)}.
\end{equation}
The source page displays the additional condition $\norm Z<\log2$ next to
its associative expansion. That condition is \emph{not} needed to define
$Z=\log(e^Xe^Y)$ by the convergent Mercator series: the bound actually
proved is $\norm Z\leq M_s(1)$, and one should not silently infer a
further norm bound on this constructed $Z$. The condition is relevant in
the converse direction. Suppose $Z_0$ is \emph{given} with $e^{Z_0}=e^Xe^Y$
and $\norm{Z_0}<\log2$. Then for $|t|\norm{Z_0}<\log2$ we have
$\norm{e^{tZ_0}-I}\leq e^{|t|\norm{Z_0}}-1<1$, so the Mercator logarithm
$L(t)=\log(e^{tZ_0})$ is a convergent power series in $t$ on that disk. Near
$t=0$ it equals $tZ_0$ by the formal identity $\log(e^{u})=u$ applied in the
commutative closed subalgebra generated by $Z_0$; both sides being analytic
in $t$ on the disk, they agree throughout it, and at $t=1$ we get
$\log(e^{Z_0})=Z_0$. Combined with $s<\log2$, this shows that the page's two
restrictions do suffice to conclude that its given $Z_0$ coincides with the
Mercator logarithm of $e^Xe^Y$, and hence with the BCH sum. Without a
restriction on $Z_0$ the conclusion fails: other logarithms exist, as
the trace discussion next shows.

\subsection{The trace identity and its branch restriction}
\begin{proposition}\label{prop:trace}
For the formal BCH series, and for finite complex or real matrices whenever
the homogeneous series converges,
\begin{equation}\label{eq:trace}
 \tr Z(X,Y)=\tr X+\tr Y.
\end{equation}
For an arbitrary complex matrix $L$ with $e^L=e^Xe^Y$, the branch-independent
conclusion is only
\begin{equation}\label{eq:tracebranch}
 \tr L-\tr X-\tr Y\in2\pi\ii\Z.
\end{equation}
\end{proposition}
\begin{proof}
The identity $\tr(AB)=\tr(BA)$ follows by interchanging the two summation
indices in $\sum_{i,j}A_{ij}B_{ji}$. Hence every commutator has trace zero,
and every nested bracket of length at least two, being a commutator, has
trace zero. By \cref{thm:formalbch} each $Z_n$ with $n\geq2$ is a sum of
such brackets, so $\tr Z_n=0$ for $n\geq2$ and $\tr Z_1=\tr X+\tr Y$. This
proves \eqref{eq:trace} degree by degree; when the series converges,
continuity of the trace (a linear functional on a finite-dimensional space)
gives it for the sum.

For \eqref{eq:tracebranch} we first show $\det e^A=e^{\tr A}$ for every
complex square matrix $A$. Every complex matrix is triangularizable: choose
an eigenvector $v_1$, complete it to a basis, and observe that in this basis
$A$ is block upper triangular with a $1\times1$ block; induct on the
dimension of the quotient. In an upper-triangular basis with diagonal
$\lambda_1,\ldots,\lambda_d$, every power $A^k$ is upper triangular with
diagonal $\lambda_j^k$, so $e^A$ is upper triangular with diagonal
$e^{\lambda_j}$, and $\det e^A=\prod_je^{\lambda_j}=e^{\sum\lambda_j}
=e^{\tr A}$. Applying this to $e^L=e^Xe^Y$ and using multiplicativity of
the determinant gives $e^{\tr L}=e^{\tr X}e^{\tr Y}$, that is
$e^{\tr L-\tr X-\tr Y}=1$, and the scalar exponential takes the value $1$
exactly on $2\pi\ii\Z$.
\end{proof}
A scalar counterexample to unrestricted use of the principal logarithm:
take $X=Y=3\pi\ii/4$ in dimension one. The BCH logarithm is $X+Y=3\pi\ii/2$,
whereas the principal logarithm of the product $e^Xe^Y=e^{3\pi\ii/2}=-\ii$
is $-\pi\ii/2$; they differ by $2\pi\ii$. Thus the word ``logarithm'' in
\eqref{eq:trace} cannot be left branch-independent. When $L,X,Y$ are all
real matrices the difference in \eqref{eq:tracebranch} is real, hence zero.

\subsection{The \texorpdfstring{$\mathfrak{sl}_2$}{sl2} example, completely resolved}
\begin{theorem}\label{thm:sl2}
Let
\begin{equation}\label{eq:sl2XY}
 X=\begin{pmatrix}0&\ii\pi\\\ii\pi&0\end{pmatrix},\qquad
 Y=\begin{pmatrix}0&1\\0&0\end{pmatrix}.
\end{equation}
Then
\begin{equation}\label{eq:sl2product}
 e^Xe^Y=\begin{pmatrix}-1&-1\\0&-1\end{pmatrix}=-I-Y.
\end{equation}
The complete set of complex $2\times2$ logarithms of this matrix is
\begin{equation}\label{eq:sl2logs}
 L_k=(2k+1)\pi\ii\,I+Y,\qquad k\in\Z.
\end{equation}
None of them is traceless, so none lies in $\mathfrak{sl}_2(\C)$, and the
homogeneous BCH series $\sum_nZ_n(X,Y)$ does not converge, even
conditionally, although the product has infinitely many logarithms in
$\mathfrak{gl}_2(\C)$.
\end{theorem}
\begin{proof}
Since $X^2=-\pi^2I$, the even and odd parts of the exponential series give
$e^X=\sum_m\frac{(-\pi^2)^m}{(2m)!}I+\sum_m\frac{(-\pi^2)^m}{(2m+1)!}X
=\cos\pi\,I+\frac{\sin\pi}{\pi}X=-I$. Since $Y^2=0$, $e^Y=I+Y$. Hence
$e^Xe^Y=-I-Y$, which is \eqref{eq:sl2product}.

Let $L$ be any logarithm, $e^L=-I-Y$. A matrix commutes with every power of
itself, hence with its exponential; so $L$ commutes with $-I-Y$ and
therefore with $Y$. Writing $L=\bigl(\begin{smallmatrix}a&b\\c&d\end{smallmatrix}\bigr)$,
the equation $LY=YL$ reads
$\bigl(\begin{smallmatrix}0&a\\0&c\end{smallmatrix}\bigr)
=\bigl(\begin{smallmatrix}c&d\\0&0\end{smallmatrix}\bigr)$, so $c=0$ and
$a=d$: the centralizer of $Y$ consists of the matrices $aI+bY$. For such a
matrix, $aI$ and $bY$ commute and $(bY)^2=0$, so
$e^{aI+bY}=e^a(I+bY)$. Equality with $-I-Y$ forces $e^a=-1$ and $e^ab=-1$,
that is $a=(2k+1)\pi\ii$ and $b=1$. This proves \eqref{eq:sl2logs}. The
trace of $L_k$ is $2(2k+1)\pi\ii\neq0$.

Suppose, to reach a contradiction, that $\sum_{n\geq1}Z_n(X,Y)$ converged to
some matrix $S$. Then $Z_n\to0$, so the coefficients are bounded and the
matrix power series $Z(t)=\sum_{n\geq1}t^nZ_n(X,Y)$ converges absolutely,
entry by entry, for $|t|<1$, defining an analytic matrix function there. For
$|t|(\norm X+\norm Y)<\log2$ we have $Z(t)=\BCH(tX,tY)$ by homogeneity, and
\cref{thm:convergence}(i) gives $e^{Z(t)}=e^{tX}e^{tY}$. Both sides are
analytic in $t$ on the unit disk (the exponential of an analytic matrix
function is analytic, being an absolutely and locally uniformly convergent
series of analytic functions), so by the identity theorem
$e^{Z(t)}=e^{tX}e^{tY}$ for all $|t|<1$.

Abel's limit theorem now gives $Z(t)\to S$ as $t\uparrow1$ along the real
axis. For completeness: let $S_m=\sum_{n=1}^mZ_n$ and $S_0=0$. For $0<t<1$,
summation by parts gives
\[
 Z(t)=\sum_{n\geq1}t^n(S_n-S_{n-1})=\sum_{n\geq1}t^nS_n-\sum_{n\geq1}t^{n+1}S_n
 =(1-t)\sum_{n\geq1}t^nS_n,
\]
all series converging absolutely because the $S_n$ are bounded. Since
$(1-t)\sum_{n\geq1}t^n=t$, we obtain
$Z(t)-tS=(1-t)\sum_{n\geq1}t^n(S_n-S)$. Given $\eps>0$ choose $N$ with
$\norm{S_n-S}<\eps$ for all $n>N$. The finite part
$(1-t)\sum_{n\leq N}t^n(S_n-S)$ tends to zero as $t\uparrow1$, and the tail
is bounded by $\eps(1-t)\sum_{n>N}t^n\leq\eps$. Hence
$\limsup_{t\uparrow1}\norm{Z(t)-tS}\leq\eps$ for every $\eps>0$, so
$Z(t)-tS\to0$ and $Z(t)\to S$. Continuity of the exponential then gives $e^S=\lim e^{Z(t)}
=\lim e^{tX}e^{tY}=e^Xe^Y=-I-Y$. But every $Z_n$ is traceless by
\cref{prop:trace}, so $\tr S=\lim\tr S_m=0$, contradicting the
classification \eqref{eq:sl2logs}. Hence the series does not converge.
\end{proof}
This example simultaneously distinguishes existence of a logarithm,
existence of a logarithm in the specified Lie algebra, and convergence of
the BCH series. They are genuinely different questions. The obstruction is
not that the invertible matrix $-I-Y$ lacks a logarithm; it has infinitely
many. The obstruction is that it has no logarithm in the traceless Lie
algebra generated by $X$ and $Y$.

\section{Lie groups: exponential coordinates, the local--global distinction, and nilpotent group laws}\label{sec:liegroups}
\subsection{Matrix groups}
For a matrix Lie group $G\subseteq\GL_d$, its Lie algebra $\g$ is the
tangent space at the identity. The left-invariant vector field associated
with $X\in\g$ is $X^L(g)=gX$, and the right-invariant field is $X^R(g)=Xg$. Regarding these as vector fields on the open
subset $\GL_d$ of the vector space of matrices, the bracket of two vector
fields $U,V$ is $[U,V](g)=DV(g)[U(g)]-DU(g)[V(g)]$, and since
$DX^L(g)[H]=HX$,
\[
 [X^L,Y^L](g)=(gX)Y-(gY)X=g(XY-YX)=[X,Y]^L(g).
\]
Thus the Lie bracket of $\g$, defined through left-invariant fields, is the
matrix commutator. The integral curve of $X^L$ through $I$ solves
$g'=gX$, $g(0)=I$; the matrix exponential $e^{tX}$ solves this equation, and
the solution is unique (differentiate $g(t)e^{-tX}$), so the Lie-group
exponential is the ordinary matrix exponential. This proves the matrix
statements of the source page.

For matrix Lie algebras, and more generally for closed Lie subalgebras of a
Banach algebra, the fact that the BCH logarithm lies in the Lie algebra can
be obtained directly from the differential equation of \cref{thm:ode},
without the formal theory of \cref{sec:hopf}.
\begin{proposition}[The BCH logarithm in a closed Lie subalgebra]\label{prop:closedlie}
Let $\cA$ be a real or complex unital Banach algebra. There is $\delta>0$
such that for every closed linear subspace $\g\subseteq\cA$ that is closed
under the commutator $[a,b]=ab-ba$, and all $X,Y\in\g$ with
$\norm X+\norm Y<\delta$, the logarithm $Z=\log(e^Xe^Y)$ of
\cref{cor:bchlog}, that is, the sum of the BCH series, lies in $\g$. In
particular, for a matrix Lie algebra $\g\subseteq\mathfrak{gl}_d$ and $X,Y\in\g$
sufficiently small, $e^Xe^Y=e^Z$ with $Z\in\g$.
\end{proposition}
\begin{proof}
Take $\delta$ as in \cref{thm:ode}, shrunk so that $Z(t)=\log(e^Xe^{tY})$
satisfies $\norm{Z(t)}\leq\tfrac18$ for $|t|\leq1$; then $Z(0)=X$ and
$Z'(t)=F(Z(t))$ with $F(W)=\beta(\ad_W)Y$. The map $F$ sends $\g$ into
itself: $\ad_W^nY$ is an iterated commutator of elements of $\g$, and $\g$
is closed, so the series $\sum_n\bplus n\ad_W^nY$ lies in $\g$. Moreover $F$
is Lipschitz on the ball $\norm W\leq\tfrac38$: with
$\norm{\ad_W}\leq2\norm W\leq\tfrac34$, one has
$\norm{\ad_W^{n}-\ad_{W'}^{n}}\leq n(\tfrac34)^{n-1}\norm{\ad_W-\ad_{W'}}
\leq2n(\tfrac34)^{n-1}\norm{W-W'}$, and $|\bplus n|\leq1$ by
\eqref{eq:bernrec}, so $\norm{F(W)-F(W')}\leq K\norm Y\,\norm{W-W'}$ with
$K=2\sum_{n\geq1}n(\tfrac34)^{n-1}$.

Let $d(t)=\operatorname{dist}(Z(t),\g)$, a continuous function on $[0,1]$
with $d(0)=0$. For $w\in\g$ with $\norm{Z(t)-w}$ close to $d(t)$ (so that
$\norm w\leq\tfrac38$) and $h>0$, the element $w+hF(w)$ lies in $\g$, hence
\begin{align*}
 d(t+h)&\leq\norm{Z(t+h)-w-hF(w)}\\
 &\leq\norm{Z(t+h)-Z(t)-hF(Z(t))}+\norm{Z(t)-w}+h\norm{F(Z(t))-F(w)}\\
 &\leq o(h)+(1+hK\norm Y)\norm{Z(t)-w}.
\end{align*}
Taking the infimum over $w$ gives $d(t+h)\leq(1+hK\norm Y)d(t)+o(h)$, so
the upper right Dini derivative of $d$ is at most $K\norm Y\,d(t)$.
Gronwall's inequality yields $d\equiv0$ on $[0,1]$, and since $\g$ is
closed, $Z=Z(1)\in\g$. A matrix Lie algebra is finite-dimensional, hence
closed, and closed under the commutator, so the last assertion follows.
\end{proof}

\subsection{The Maurer--Cartan form and an intrinsic differential of the exponential}
None of the assertions on the page requires a matrix embedding. For an
arbitrary finite-dimensional real or complex Lie group $G$ with Lie algebra
$\g$, define the left Maurer--Cartan form by
$\theta_g=(L_{g^{-1}})_*:T_gG\to T_eG=\g$; in a matrix group it is
$g^{-1}\dd g$. On a left-invariant field it takes the constant value
$\theta(U^L)=U$.

\begin{proposition}[Maurer--Cartan equation]\label{prop:MC}
The $\g$-valued one-form $\theta$ satisfies
\begin{equation}\label{eq:MC}
 \dd\theta+\tfrac12[\theta,\theta]=0,\qquad
 \tfrac12[\theta,\theta](U,V):=[\theta(U),\theta(V)].
\end{equation}
Consequently, for a smooth two-parameter map $g(s,t)$ into $G$, the
left-trivialized partial derivatives
$\alpha=\theta(\partial_sg)$ and $\eta=\theta(\partial_tg)$ satisfy the
mixed-variation identity
\begin{equation}\label{eq:mixedvariation}
 \partial_s\eta-\partial_t\alpha+[\alpha,\eta]=0.
\end{equation}
\end{proposition}
\begin{proof}
For a one-form $\omega$ and vector fields $U,V$ one has
$\dd\omega(U,V)=U(\omega(V))-V(\omega(U))-\omega([U,V])$. Evaluate on
left-invariant fields $U^L,V^L$: the functions $\theta(V^L)=V$ and
$\theta(U^L)=U$ are constant, so the first two terms vanish, and
$[U^L,V^L]=[U,V]^L$ by definition of the Lie algebra bracket; hence
$\dd\theta(U^L,V^L)=-\theta([U,V]^L)=-[U,V]$. On the other hand
$\frac12[\theta,\theta](U^L,V^L)=[\theta(U^L),\theta(V^L)]=[U,V]$. Both
sides of \eqref{eq:MC} are tensorial (bilinear over functions) and
left-invariant fields span every tangent space, so \eqref{eq:MC} holds
identically. In a matrix group the same result is the direct computation
$\dd(g^{-1})=-g^{-1}(\dd g)g^{-1}$, whence
$\dd(g^{-1}\dd g)=-g^{-1}\dd g\wedge g^{-1}\dd g$.

For \eqref{eq:mixedvariation}, pull $\theta$ back by $g$:
$g^*\theta=\alpha\,\dd s+\eta\,\dd t$. Pullback commutes with $\dd$ and
with the bracket of forms, so \eqref{eq:MC} gives
$\dd(g^*\theta)=-[\alpha,\eta]\,\dd s\wedge\dd t$, while directly
$\dd(\alpha\,\dd s+\eta\,\dd t)=(\partial_s\eta-\partial_t\alpha)\,\dd s\wedge\dd t$.
Comparing coefficients gives \eqref{eq:mixedvariation}. (In a matrix group:
$\partial_s(g^{-1}\partial_tg)=-\alpha\eta+g^{-1}\partial_s\partial_tg$ and
$\partial_t(g^{-1}\partial_sg)=-\eta\alpha+g^{-1}\partial_t\partial_sg$;
subtracting gives $-[\alpha,\eta]$.)
\end{proof}

\begin{proposition}[Intrinsic differential of the exponential]\label{prop:intrinsicdexp}
For $X,H\in\g$,
\begin{equation}\label{eq:intrinsicdexp}
 (L_{\exp(-X)})_*(\dd\exp)_X(H)=\phi(\ad_X)H=\int_0^1e^{-u\ad_X}H\,\dd u,
\end{equation}
and similarly $(R_{\exp(-X)})_*(\dd\exp)_X(H)=\phi(-\ad_X)H$. In
particular $(\dd\exp)_0=\id_\g$, and $\Ad_{\exp X}=e^{\ad_X}$ intrinsically.
\end{proposition}
\begin{proof}
Put $g(s,t)=\exp\bigl(s(X+tH)\bigr)$. For fixed $t$, $s\mapsto g(s,t)$ is
the one-parameter subgroup with velocity $X+tH$, so its left-trivialized
velocity is $\alpha(s,t)=\theta(\partial_sg)=X+tH$. Let
$\eta(s,t)=\theta(\partial_tg)$; then $\eta(0,t)=0$ because $g(0,t)=e$
for all $t$. By \eqref{eq:mixedvariation},
$\partial_s\eta=\partial_t\alpha-[\alpha,\eta]=H-[X+tH,\eta]$. At $t=0$
this is the linear equation $\partial_s\eta=H-\ad_X\eta$ with
$\eta(0)=0$, whose unique solution is
$\eta(s)=\int_0^se^{-(s-u)\ad_X}H\,\dd u$ (check by differentiation:
$\partial_s\eta=H-\ad_X\eta$). At $s=1$, substituting $v=1-u$,
$\eta(1,0)=\int_0^1e^{-v\ad_X}H\,\dd v=\phi(\ad_X)H$ by the termwise
integration of \cref{thm:dexp}. Since $\eta(1,0)=\theta(\partial_tg(1,0))
=(L_{\exp(-X)})_*(\dd\exp)_X(H)$, this is \eqref{eq:intrinsicdexp}. The
right-trivialized version follows by applying the same argument to the
right Maurer--Cartan form, or from
$(R_{\exp(-X)})_*=\Ad_{\exp X}\circ(L_{\exp(-X)})_*$ together with
$\Ad_{\exp X}\phi(\ad_X)=e^{\ad_X}\phi(\ad_X)=\phi(-\ad_X)$. At $X=0$ the
formula gives $(\dd\exp)_0(H)=H$. Finally, differentiating the
one-parameter group of automorphisms $t\mapsto\Ad_{\exp(tX)}$ gives
$\frac{\dd}{\dd t}\Ad_{\exp(tX)}=\ad_X\Ad_{\exp(tX)}$ with
$\Ad_{\exp(0)}=I$, whose unique solution is $e^{t\ad_X}$; this is the
intrinsic Campbell identity.
\end{proof}

\subsection{Local exponential coordinates and the local BCH law}
\begin{theorem}[Local BCH multiplication]\label{thm:local}
Let $G$ be a finite-dimensional real or complex Lie group. There is a
neighborhood $U$ of $0$ in $\g$ on which $\exp$ is a diffeomorphism onto a
neighborhood of the identity. After shrinking $U$, for all $X,Y\in U$,
\begin{equation}\label{eq:localproduct}
 \exp X\exp Y=\exp\bigl(\BCH(X,Y)\bigr),
\end{equation}
where the universal series converges absolutely in $\g$. Hence the local
multiplication law of $G$ is determined entirely by the Lie bracket.
\end{theorem}
\begin{proof}
By \cref{prop:intrinsicdexp}, $(\dd\exp)_0=\id$, so the inverse function
theorem provides the chart. Choose a norm on $\g$ and $\kappa$ with
$\norm{[u,v]}\leq\kappa\norm u\norm v$. By \cref{thm:convergence}(iv), for
$\kappa(\norm X+\norm Y)<\log2$ the series $Z(t)=\BCH(X,tY)=\sum_nZ_n(X,tY)$
converges absolutely, uniformly for $t\in[0,1]$, and so does the series
with $X,Y$ replaced by $rX,rY$ for some $r>1$, provided $X,Y$ are small
enough. Consequently $Z(t)$ is given by a power series in $t$ with radius
greater than $1$, and it may be differentiated termwise in $t$. The
differential equation \eqref{eq:ode}, $Z'(t)=\beta(\ad_{Z(t)})Y$, holds in
the formal setting coefficient by coefficient in $t$ (each coefficient
being a Lie polynomial identity in $X,Y$, valid in $\g$ by universality,
\cref{sec:universal}); for $\norm{Z(t)}$ small enough that
$\norm{\ad_{Z(t)}}<2\pi$, the right side is a convergent series in $t$ with
the same coefficients, so the equation holds analytically.

Now let $L(t)=\log(\exp X\exp(tY))$ in the exponential chart, defined for
small $X,Y$ and $t\in[0,1]$ after shrinking $U$ so that the relevant
products stay in the chart. Its left-trivialized derivative is
$\theta\bigl(\tfrac{\dd}{\dd t}\exp X\exp(tY)\bigr)=Y$ (the curve
$t\mapsto\exp X\exp(tY)$ is the left translate by $\exp X$ of a
one-parameter subgroup). By \cref{prop:intrinsicdexp} the same
left-trivialized derivative equals $\phi(\ad_{L(t)})L'(t)$. For $L$ small,
$\phi(\ad_L)$ is invertible with inverse $\beta(\ad_L)$, so
$L'=\beta(\ad_L)Y$ with $L(0)=X$. The map $L\mapsto\beta(\ad_L)Y$ is
analytic on a neighborhood of $0$, so the initial value problem has a
unique solution by the Picard--Lindel\"of theorem. Both $L(t)$ and $Z(t)$
solve it; hence $L(1)=Z(1)=\BCH(X,Y)$, which is \eqref{eq:localproduct}
since $\exp$ is injective on $U$.
\end{proof}
The local assertion is not a global classification. The additive group
$\R$ and the circle group $S^1=\{e^{\ii t}\}$ have the same one-dimensional
abelian Lie algebra, but one is noncompact and the other compact, so they
are not isomorphic as Lie groups; they differ in topology, not in local
structure. The identity map of their Lie algebras integrates from $\R$ to
$S^1$ as $t\mapsto e^{\ii t}$, but not from $S^1$ to $\R$: a continuous
homomorphism from the compact group $S^1$ to $\R$ has compact image, which is
a compact subgroup of $\R$, hence $\{0\}$, so its differential is zero.

\subsection{Homomorphisms, representations, and integration}
\begin{theorem}[Local Lie correspondence]\label{thm:localcorr}
Let $f:\g\to\h$ be a homomorphism of Lie algebras of Lie groups $G,H$. Then
\begin{equation}\label{eq:localhom}
 F(\exp_GX)=\exp_H(fX)
\end{equation}
defines a smooth map on an identity neighborhood of $G$ which is a local
homomorphism: $F(gh)=F(g)F(h)$ whenever $g,h,gh$ lie in a suitable smaller
neighborhood. Conversely, the differential at the identity of a smooth
local homomorphism preserves Lie brackets, and a local homomorphism is
determined near the identity by that differential.
\end{theorem}
\begin{proof}
On the exponential neighborhood of \cref{thm:local}, \eqref{eq:localhom}
defines a smooth map, being the composition $\exp_H\circ f\circ\log_G$. By
\cref{prop:symmetry}, $f$ preserves every $Z_n$, and by continuity it
preserves the convergent sum: $f(\BCH(X,Y))=\BCH(fX,fY)$. Hence, for $X,Y$
small enough that \eqref{eq:localproduct} applies in both groups,
\begin{align*}
 F(\exp X\exp Y)&=F\bigl(\exp\BCH(X,Y)\bigr)=\exp_H\bigl(f\,\BCH(X,Y)\bigr)
 =\exp_H\bigl(\BCH(fX,fY)\bigr)\\
 &=\exp_H(fX)\exp_H(fY)=F(\exp X)F(\exp Y).
\end{align*}
For the converse, let $F$ be a smooth local homomorphism with differential
$A=(\dd F)_e$. Differentiating $F(g\exp(tX))=F(g)F(\exp(tX))$ at $t=0$
shows that $F$ relates the left-invariant field $X^L$ on $G$ to the field
$(AX)^L$ on $H$, wherever both are defined. Related vector fields have
related brackets (a standard consequence of the chain rule applied to
$F_*$), so $[X^L,Y^L]=[X,Y]^L$ is related to $[(AX)^L,(AY)^L]=[AX,AY]^L$;
evaluating at $e$ gives $A[X,Y]=[AX,AY]$. Finally, since $F$ maps the
one-parameter subgroup $\exp(tX)$ to a local one-parameter subgroup of $H$
with initial velocity $AX$, uniqueness of integral curves gives
$F(\exp(tX))=\exp_H(tAX)$ for small $t$; as $\exp$ is onto an identity
neighborhood, $F$ is determined near $e$ by $A$.
\end{proof}

\begin{theorem}[Integration from a simply connected group]\label{thm:integration}
If $G$ is connected and simply connected and $H$ is a Lie group, every Lie
algebra homomorphism $f:\g\to\h$ integrates to a unique smooth homomorphism
$F:G\to H$ with $(\dd F)_e=f$. In particular every finite-dimensional
representation of $\g$ integrates to a representation of $G$.
\end{theorem}
\begin{proof}
Let $F_0$ be the local homomorphism of \cref{thm:localcorr}, defined and
multiplicative on an open identity neighborhood $W$ in the sense that
$F_0(uv)=F_0(u)F_0(v)$ whenever $u,v,uv\in W$. Choose a smaller symmetric
identity neighborhood $W_1$ with $W_1W_1\subseteq W$.

\emph{Step 1: paths.} Let $\gamma:[0,1]\to G$ be a continuous path with
$\gamma(0)=e$ and $\gamma(1)=g$. By uniform continuity there is a partition
$0=t_0<t_1<\cdots<t_m=1$ such that every increment
$u_j=\gamma(t_{j-1})^{-1}\gamma(t_j)$ lies in $W_1$. Define
\[
 F_\gamma(g)=F_0(u_1)F_0(u_2)\cdots F_0(u_m).
\]
This does not depend on the partition. Inserting one additional point $t'$
between $t_{j-1}$ and $t_j$ replaces $u_j$ by the product $u'u''$ of two
increments; if the partition is fine enough that $u',u''\in W_1$, then
$u'u''=u_j\in W$ and $F_0(u_j)=F_0(u')F_0(u'')$ by local multiplicativity,
so the product is unchanged. Any two sufficiently fine partitions have a
common refinement obtained by inserting points one at a time, so
$F_\gamma(g)$ is well defined.

\emph{Step 2: homotopy invariance.} Let $\Gamma:[0,1]^2\to G$ be a
continuous homotopy with $\Gamma(\cdot,0)=\gamma_0$, $\Gamma(\cdot,1)=\gamma_1$,
$\Gamma(0,\cdot)=e$, $\Gamma(1,\cdot)=g$. By uniform continuity there is a
rectangular grid so fine that for any two grid points $p,q$ of the same
closed cell, $\Gamma(p)^{-1}\Gamma(q)\in W_1$. Around a single cell with
corners $a,b,c,d$ (in cyclic order), the four increments
$u=\Gamma(a)^{-1}\Gamma(b)$, $v=\Gamma(b)^{-1}\Gamma(c)$,
$u'=\Gamma(a)^{-1}\Gamma(d)$, $v'=\Gamma(d)^{-1}\Gamma(c)$ satisfy
$uv=u'v'=\Gamma(a)^{-1}\Gamma(c)\in W_1$, so local multiplicativity gives
$F_0(u)F_0(v)=F_0(uv)=F_0(u')F_0(v')$. Hence the products of $F_0$ along the
two edge paths of the cell from $a$ to $c$ coincide. Sweeping across the
grid one cell at a time, replacing the lower-right edge pair of a cell by
its upper-left pair, converts the grid path along the bottom and right
boundary into the path along the left and top boundary without changing the
product. The left and right boundary edges are constant paths (at $e$ and
at $g$) and contribute factors $F_0(e)=e$. The bottom and top boundary
paths are $\gamma_0$ and $\gamma_1$ evaluated at the grid points, whose
increments form fine partitions in the sense of Step 1. Therefore
$F_{\gamma_0}(g)=F_{\gamma_1}(g)$.

\emph{Step 3: definition of $F$.} Since $G$ is simply connected, any two
paths from $e$ to $g$ are homotopic with fixed endpoints, and since $G$ is
connected such paths exist. Define $F(g)=F_\gamma(g)$ for any such path.

\emph{Step 4: multiplicativity, smoothness, differential.} Given paths
$\gamma$ from $e$ to $g$ and $\delta$ from $e$ to $h$, the concatenation of
$\gamma$ with the left translate $g\delta$ is a path from $e$ to $gh$. The
increments of $g\delta$ are $(g\delta(t_{j-1}))^{-1}g\delta(t_j)
=\delta(t_{j-1})^{-1}\delta(t_j)$, the original increments of $\delta$.
Hence $F(gh)=F(g)F(h)$. For $u\in W_1$, appending to a path to $g$ the
short path $t\mapsto g\exp(t\log u)$ from $g$ to $gu$ gives
$F(gu)=F(g)F_0(u)$; thus $F$ agrees near $g$ with the smooth map
$u\mapsto F(g)F_0(u)$, so $F$ is smooth, and at $g=e$ it agrees with $F_0$,
so $(\dd F)_e=f$.

\emph{Step 5: uniqueness.} Any homomorphism $F'$ with differential $f$
satisfies $F'(\exp X)=\exp_H(fX)$ (it maps one-parameter subgroups to
one-parameter subgroups with the given initial velocity), so $F'=F_0=F$ on
the exponential neighborhood. The set where $F'=F$ is a subgroup
containing an identity neighborhood, hence an open subgroup; its cosets
are open, so it is also closed; connectedness makes it all of $G$.

Taking $H=\GL(E)$ for a finite-dimensional vector space $E$ gives the
statement about representations.
\end{proof}
Without simple connectedness the path construction can yield a
nontrivial value around a loop; this is the precise global obstruction
absent from the local BCH law. Descent from a simply connected covering
group to a specified quotient requires the covering kernel to act trivially.

\subsection{Nilpotent Lie algebras and the polynomial group law}\label{sec:nilpotent}
Write $\g^1=\g$ and $\g^{r+1}=[\g,\g^r]$ for the lower central series. The
algebra is nilpotent of class at most $c$ if $\g^{c+1}=0$.
\begin{theorem}[Termination and the nilpotent group law]\label{thm:nilpotent}
In a nilpotent Lie algebra of class at most $c$, every $Z_n$ with $n>c$
vanishes identically, and
\begin{equation}\label{eq:nilpotentlaw}
 X*Y=\BCH(X,Y)=\sum_{n=1}^{c}Z_n(X,Y)
\end{equation}
is a polynomial group law on the vector space $\g$ with identity $0$ and
inverse $-X$. For finite-dimensional real $\g$ this makes $\g$ a connected,
simply connected Lie group whose Lie algebra is $\g$ with its original
bracket and whose exponential map is the identity of $\g$. Any connected
simply connected Lie group with Lie algebra $\g$ is isomorphic to it, and
in every Lie group with Lie algebra $\g$ the identity
$\exp X\exp Y=\exp\BCH(X,Y)$ holds for all $X,Y\in\g$.
\end{theorem}
\begin{proof}
By the Jacobi identity in the form \eqref{eq:derivation}, induction on $r$
gives $[\g^r,\g^s]\subseteq\g^{r+s}$: for $r=1$ this is the definition,
and if it holds for $r$ then for $x\in\g$, $y\in\g^r$, $z\in\g^s$,
$[[x,y],z]=[x,[y,z]]-[y,[x,z]]\in[\g,\g^{r+s}]+[\g^r,\g^{s+1}]\subseteq\g^{r+s+1}$.
Hence every bracket monomial of degree $n$ in elements of $\g$ lies in
$\g^n$, and vanishes for $n>c$. Every $Z_n$ is such a combination, so
$Z_n=0$ for $n>c$.

The formal identities \eqref{eq:assoc} and \eqref{eq:unitinverse} are
identities between Lie series in three free variables. Evaluating them in
$\g$ turns both sides into finite polynomial identities, because all Lie
monomials of degree greater than $c$ vanish, including those produced by
substituting one polynomial into another. Thus
$(X*Y)*T=X*(Y*T)$, $X*0=0*X=X$, and $X*(-X)=0$ hold for all $X,Y,T\in\g$,
not merely near zero. So $(\g,*)$ is a group, with polynomial multiplication
and inversion; it is a Lie group whose underlying manifold is a vector
space, hence connected and simply connected.

Its Lie algebra is computed from the group commutator. Write $a=sX$,
$b=tY$, and let $O(3)$ denote terms of total degree at least three in
$s,t$. Using $u*v=u+v+\frac12[u,v]+O(3)$ three times,
\begin{align*}
 (a*b)*(-a)&=\bigl(a+b+\tfrac12[a,b]\bigr)-a+\tfrac12[a+b,-a]+O(3)
 =b+[a,b]+O(3),\\
 \bigl((a*b)*(-a)\bigr)*(-b)&=\bigl(b+[a,b]\bigr)-b+\tfrac12[b,-b]+O(3)
 =[a,b]+O(3)=st[X,Y]+O(3).
\end{align*}
The mixed second derivative $\partial_s\partial_t$ at $s=t=0$ of the group
commutator $g_sh_tg_s^{-1}h_t^{-1}$ of two one-parameter subgroups is the
Lie bracket of their initial velocities in the Lie algebra of any Lie
group, so the bracket of $(\g,*)$ is the original bracket of $\g$.
(Alternatively, the left-invariant field of $X$ is $g\mapsto\frac{\dd}{\dd t}
g*(tX)|_{t=0}$, and computing brackets of these polynomial fields gives the
same result.) Since $Z(sX,tX)=(s+t)X$ (all brackets of multiples of $X$
vanish), the one-parameter subgroups are $t\mapsto tX$, so the exponential
map is the identity.

If $G$ is a connected simply connected Lie group with Lie algebra $\g$,
\cref{thm:integration} integrates the identity map of $\g$ in both
directions, giving homomorphisms $(\g,*)\to G$ and $G\to(\g,*)$ whose
composites have differential the identity, hence are the identity by the
uniqueness part; so they are inverse isomorphisms. In particular the
exponential map of $G$ is a global diffeomorphism. For an arbitrary Lie
group $G'$ with Lie algebra $\g$, the homomorphism $(\g,*)\to G'$
integrating the identity map sends $X$ to $\exp_{G'}X$ (it maps
one-parameter subgroups to one-parameter subgroups), and multiplicativity
reads $\exp_{G'}(X*Y)=\exp_{G'}X\exp_{G'}Y$.
\end{proof}
Termination is a statement about Lie brackets, not about matrix powers: a
nilpotent Lie algebra may be represented by matrices that are not
themselves nilpotent.

For a basis $P,Q,C$ with $[P,Q]=C$ and all other brackets zero (the
Heisenberg algebra, of class two), \eqref{eq:nilpotentlaw} reads
\begin{equation}\label{eq:heisenberglaw}
 (x,y,z)*(x',y',z')=\bigl(x+x',\,y+y',\,z+z'+\tfrac12(xy'-yx')\bigr)
\end{equation}
in coordinates $xP+yQ+zC$, since $\frac12[xP+yQ,x'P+y'Q]=\frac12(xy'-yx')C$.
A concrete matrix instance is $P=E_{12}$, $Q=E_{23}$, $C=E_{13}$ in the
strictly upper triangular $3\times3$ matrices, where both sides of
$e^Xe^Y=e^{X+Y+[X,Y]/2}$ are finite matrix polynomials. This group law
underlies the Weyl and displacement identities of \cref{sec:quantum}.

\section{Exactly summable commutator relations}\label{sec:special}
The universal series is infinite, but algebraic relations can remove almost
all of it. The proofs below establish \emph{global} exponential identities,
not merely small-neighborhood Taylor expansions. For Banach algebras all
exponentials exist; for finite-dimensional Lie groups the same proofs use
one-parameter subgroups and their right-trivialized logarithmic derivatives
(\cref{prop:intrinsicdexp}), together with uniqueness of solutions of the
ordinary differential equation $g'=V(t)^R(g)$ on the group. Complex
scalar parameters are understood in a complex Banach algebra or complex Lie
group; in a real Lie group all parameters and linear combinations must be
real.

\subsection{Commuting elements and central commutators}
\begin{proposition}[Commuting case]\label{prop:commuting}
If $[X,Y]=0$, then $Z(X,Y)=X+Y$ formally, and $e^Xe^Y=e^{X+Y}$ for all
bounded $X,Y$ without a smallness hypothesis.
\end{proposition}
\begin{proof}
For commuting elements the binomial theorem holds, and the Cauchy product
of the two absolutely convergent series gives
$e^Xe^Y=\sum_{r,s}X^rY^s/(r!s!)=\sum_n\frac1{n!}\sum_{r+s=n}\binom nrX^rY^s
=\sum_n(X+Y)^n/n!=e^{X+Y}$. Formally, the same computation is valid
coefficientwise, and then $Z(X,Y)=\log e^{X+Y}=X+Y$ by \cref{lem:explog};
this is also the content of \cref{thm:formalbch}, since every $Z_n$ with
$n\geq2$ lies in the ideal generated by $[X,Y]$.
\end{proof}

\begin{theorem}[Central-commutator identity]\label{thm:central}
Suppose $C=[X,Y]$ commutes with both $X$ and $Y$. Then, for every scalar
parameter $t$,
\begin{align}
 e^{tX}Ye^{-tX}&=Y+tC,\label{eq:centralconj}\\
 e^{tX}e^{tY}&=\exp\bigl(t(X+Y)+\tfrac12t^2C\bigr).\label{eq:centralt}
\end{align}
In particular,
\begin{equation}\label{eq:centralbch}
 e^Xe^Y=e^{X+Y+C/2},\qquad
 e^{X+Y}=e^Xe^Ye^{-C/2},\qquad
 e^{X/2}e^Ye^{X/2}=e^{X+Y},\qquad
 e^Xe^Ye^{-X}e^{-Y}=e^{C}.
\end{equation}
These are formal identities, global identities for bounded elements of a
Banach algebra, and global identities for elements of a finite-dimensional
Lie group whose Lie algebra elements satisfy the stated bracket relations.
\end{theorem}
\begin{proof}
In the Campbell series \eqref{eq:campbell}, $\ad_X^2Y=[X,C]=0$, so all
terms after the first commutator vanish, giving \eqref{eq:centralconj}.
Let $g(t)=e^{tX}e^{tY}$. Its right logarithmic derivative is
\[
 g'(t)g(t)^{-1}=X+e^{tX}Ye^{-tX}=X+Y+tC,
\]
so $g'=(X+Y+tC)g$ with $g(0)=I$. Let $K(t)=t(X+Y)+\frac12t^2C$. Since $C$
commutes with $X+Y$, $K(t)$ commutes with $K'(t)=X+Y+tC$, and therefore
$\ad_{K}K'=0$ and \eqref{eq:rightdexp} gives
$\frac{\dd}{\dd t}e^{K(t)}=\phi(-\ad_K)K'\,e^{K}=K'e^{K}=(X+Y+tC)e^{K(t)}$,
with $e^{K(0)}=I$. Both $g$ and $e^{K}$ solve $h'=(X+Y+tC)h$, $h(0)=I$;
the difference $\Delta$ of two solutions satisfies
$\Delta'=(X+Y+tC)\Delta$, $\Delta(0)=0$, hence $\Delta\equiv0$ by
uniqueness for linear equations with continuous bounded coefficients
(Gr\"onwall), or by comparing coefficients of $t$ formally. This proves
\eqref{eq:centralt}. In a Lie group the same argument uses the
right-trivialized derivative and uniqueness of integral curves of the
time-dependent right-invariant field $t\mapsto(X+Y+tC)^R$.

Setting $t=1$ gives the first identity of \eqref{eq:centralbch}; since
$C/2$ commutes with $X+Y+C/2$, \cref{prop:commuting} gives
$e^{X+Y+C/2}e^{-C/2}=e^{X+Y}$, the second identity. For the symmetric
identity apply the first identity to the pair $(X/2,Y)$, whose commutator
is $C/2$, and then to the pair $(X/2+Y+C/4,\,X/2)$, whose commutator is
$\frac12[Y,X]=-C/2$ (all other brackets vanish and $C$ is central):
\[
 e^{X/2}e^Ye^{X/2}=e^{X/2+Y+C/4}e^{X/2}=e^{X/2+Y+C/4+X/2-C/4}=e^{X+Y}.
\]
Finally, by \eqref{eq:centralconj} and multiplicativity of conjugation,
$e^Xe^Ye^{-X}=e^{Y+C}$, and since $Y$ and $C$ commute,
$e^{Y+C}e^{-Y}=e^{C}$.
\end{proof}
Centrality is needed only relative to the Lie subalgebra generated by $X$
and $Y$, not relative to a larger ambient algebra. This proof is global and
does not infer convergence from the formal truncation alone; for unbounded
operators a common invariant domain and integrability hypotheses are still
necessary, as \cref{sec:quantum} shows.

\begin{corollary}[Several factors with central pairwise commutators]\label{cor:centralmany}
If all commutators $[X_i,X_j]$ commute with every $X_k$, then
\begin{equation}\label{eq:centralmany}
 \prod_{j=1}^{m}e^{X_j}
 =\exp\Bigl(\sum_{j=1}^{m}X_j+\frac12\sum_{1\leq i<j\leq m}[X_i,X_j]\Bigr),
\end{equation}
the product being ordered by increasing $j$.
\end{corollary}
\begin{proof}
Induct on $m$; the case $m=1$ is trivial. Let $E_{m-1}=\sum_{j<m}X_j
+\frac12\sum_{i<j<m}[X_i,X_j]$, so that $\prod_{j<m}e^{X_j}=e^{E_{m-1}}$ by
the induction hypothesis. The commutator $[E_{m-1},X_m]=\sum_{j<m}[X_j,X_m]$,
because the central part of $E_{m-1}$ commutes with $X_m$; it is a sum of
central elements, hence commutes with $E_{m-1}$ and $X_m$. \Cref{thm:central}
applied to $(E_{m-1},X_m)$ gives
$e^{E_{m-1}}e^{X_m}=\exp\bigl(E_{m-1}+X_m+\frac12\sum_{j<m}[X_j,X_m]\bigr)$,
which is \eqref{eq:centralmany} for $m$ factors.
\end{proof}
This is the all-factors form behind the Weyl and displacement
multiplication laws of \cref{sec:quantum}. In the Heisenberg algebra all
higher BCH terms vanish because every Lie monomial of degree at least three
contains a bracket with the central element $C$, in agreement with
\cref{thm:nilpotent}.

\subsection{An eigenvector of the adjoint action}
\begin{theorem}[The relation $[X,Y]=sY$]\label{thm:eigen}
Suppose $[X,Y]=sY$ with $s$ a scalar. Put
\[
 q_s(t)=\int_0^te^{-su}\,\dd u=\begin{cases}(1-e^{-st})/s,&s\neq0,\\ t,&s=0,\end{cases}
 \qquad q_s(1)=\phi(s).
\]
Then for every scalar $c$ and every $t$,
\begin{equation}\label{eq:eigenfactor}
 e^{t(X+cY)}=e^{tX}e^{c\,q_s(t)Y},\qquad\text{in particular}\qquad
 e^{X+cY}=e^Xe^{c\phi(s)Y}.
\end{equation}
Consequently, if $s\notin2\pi\ii\Z\setminus\{0\}$,
\begin{equation}\label{eq:eigenbch}
 \boxed{\displaystyle e^Xe^Y=\exp\Bigl(X+\frac{s}{1-e^{-s}}\,Y\Bigr)
 =\exp\bigl(X+\beta(s)Y\bigr),}
\end{equation}
with the coefficient interpreted as $\beta(0)=1$ at $s=0$. For every $s$,
including the resonant values,
\begin{equation}\label{eq:eigenbraid}
 e^Xe^Ye^{-X}=e^{e^sY},\qquad e^Xe^Y=e^{e^sY}e^X.
\end{equation}
The formal BCH series is
\begin{equation}\label{eq:eigenseries}
 Z(X,Y)=X+\sum_{n\geq0}\bplus{n}s^nY=X+\beta(s)Y,
\end{equation}
where the scalar series converges for $|s|<2\pi$ and diverges for
$|s|\geq2\pi$ if $Y\neq0$. All exponential identities are global in the
same settings as \cref{thm:central}.
\end{theorem}
\begin{proof}
The relation gives $\ad_X^nY=s^nY$ by induction, so the Campbell series
\eqref{eq:campbell} yields $e^{tX}Ye^{-tX}=e^{st}Y$; conjugating the
exponential series of $Y$ gives $e^{tX}e^{Y}e^{-tX}=e^{e^{st}Y}$, and at
$t=1$ this is \eqref{eq:eigenbraid}, whose second form follows by
multiplying on the right by $e^X$. No division has occurred, so these hold
for every $s$.

For \eqref{eq:eigenfactor}, let $F(t)=e^{tX}e^{c\,q_s(t)Y}$. Its right
logarithmic derivative is
\[
 F'F^{-1}=X+e^{tX}\bigl(c\,q_s'(t)Y\bigr)e^{-tX}=X+c\,e^{-st}e^{st}Y=X+cY,
\]
using $q_s'(t)=e^{-st}$ and the fact that $c\,q_s(t)Y$ commutes with its
own derivative, so that \eqref{eq:rightdexp} for the second factor reduces
to plain differentiation. Hence $F'=(X+cY)F$ with $F(0)=I$. The function
$e^{t(X+cY)}$ solves the same initial value problem, and uniqueness (as in
the proof of \cref{thm:central}) gives \eqref{eq:eigenfactor}. At $t=1$,
$q_s(1)=\phi(s)$, with the removable value $1$ at $s=0$.

If $s\notin2\pi\ii\Z\setminus\{0\}$, then $\phi(s)\neq0$ (the zeros of
$\phi$ are exactly the nonzero points of $2\pi\ii\Z$), and choosing
$c=1/\phi(s)=\beta(s)$ in \eqref{eq:eigenfactor} gives \eqref{eq:eigenbch}.

For the formal assertion \eqref{eq:eigenseries}, we show that every Lie
monomial in $X,Y$ containing at least two occurrences of $Y$ vanishes under
the relation. Consider a bracket tree evaluating such a monomial. A branch
containing no $Y$ is a bracket of copies of $X$ alone, which is $X$ if it
has length one and $0$ otherwise. A branch containing exactly one $Y$ is,
by induction on its size, a scalar multiple of $Y$: at each node one
combines either $X$ with a multiple of $Y$ (giving a multiple of $sY$) or a
zero branch with something (giving $0$). At the first node where two
branches each containing at least one $Y$ meet, both are multiples of $Y$
and their bracket is zero. Hence only the terms linear in $Y$ survive, and
by \cref{cor:linearY} these are $\beta(\ad_X)Y=\sum_n\bplus n\ad_X^nY
=\sum_n\bplus ns^nY$. The scalar series is the Taylor series of $\beta$ at
$s$, with radius exactly $2\pi$ by \cref{prop:bernoulli}; by
\eqref{eq:bernzeta} its even terms have magnitude
$2\zeta(2m)(|s|/2\pi)^{2m}\geq2(|s|/2\pi)^{2m}$, which does not tend to
zero when $|s|\geq2\pi$, so the series diverges there.
\end{proof}
The differential-equation proof of \eqref{eq:eigenbch} establishes the
exponential identity even where the Taylor series \eqref{eq:eigenseries}
diverges: resummation and Taylor convergence are distinct assertions. In
the same vein,
\begin{equation}\label{eq:eigenscaled}
 \BCH(tX,tY)=tX+t\beta(ts)Y
\end{equation}
as a Taylor series in $t$ at $0$, whose radius is exactly $2\pi/|s|$ when
$s\neq0$ and $Y\neq0$; real $s>2\pi$ therefore makes this series divergent
at $t=1$ although the resummed identity remains valid. The braiding law
\eqref{eq:eigenbraid}, being an entire exponential identity, has no
resonance exclusion at all.

\begin{example}[Resonance]\label{ex:resonance}
Let
\begin{equation}\label{eq:eigenmatrices}
 X=\begin{pmatrix}s&0\\0&0\end{pmatrix},\qquad
 Y=\begin{pmatrix}0&1\\0&0\end{pmatrix},\qquad [X,Y]=sY.
\end{equation}
Direct exponentiation gives
\[
 e^X=\begin{pmatrix}e^s&0\\0&1\end{pmatrix},\quad
 e^Xe^Y=\begin{pmatrix}e^s&e^s\\0&1\end{pmatrix},\quad
 e^{X+cY}=\begin{pmatrix}e^s&c\,(e^s-1)/s\\0&1\end{pmatrix},
\]
the last by the triangular structure ($X+cY$ has eigenvalues $s,0$ and its
exponential's off-diagonal entry is $c$ times the divided difference of
$e^z$ at $s,0$), with the continuous value $c$ at $s=0$. If $s=3\pi$, the
choice $c=\beta(3\pi)$ works although the BCH series \eqref{eq:eigenseries}
diverges. If $s=2\pi\ii k\neq0$, then $e^s=1$, every $e^{X+cY}$ equals
$I$, while $e^Xe^Y=I+Y\neq I$; thus no logarithm of the form $X+cY$ exists
at a resonance, although the logarithm $Y$ does exist, and the braiding
identity $e^Xe^Y=e^{e^sY}e^X=e^Ye^X$ remains true. What fails at resonance
is the shape $X+cY$ of the logarithm, not the existence of a logarithm.
\end{example}

\section{Zassenhaus factorization: every exponent, three recursions, a tree formula, and convergence}\label{sec:zassenhaus}
BCH combines a product of exponentials into one exponential. Zassenhaus
factorization performs the ordered operation in the opposite direction:
it writes the exponential of a sum as an ordered product of exponentials of
homogeneous Lie corrections. The order of the factors is part of the
statement. Recursive constructions are central in the computational
literature \cite{casasmuruanadinic}; explicit reorganizations are
considered in \cite{kimura}. We prove existence and uniqueness, three
recursive constructions, a finite nonrecursive formula for every exponent,
a globally convergent ordered-integral expansion, and a quantitative
convergence theorem for the product.

\subsection{Existence, uniqueness, and the residual algorithm}
\begin{theorem}[Formal Zassenhaus factorization]\label{thm:zassenhaus}
There is a unique sequence of homogeneous Lie polynomials $C_n(X,Y)$ of
degree $n$, $n\geq2$, such that
\begin{equation}\label{eq:zassenhaus}
 \boxed{\displaystyle
 e^{t(X+Y)}=e^{tX}e^{tY}\prod_{n=2}^{\infty}e^{t^nC_n(X,Y)}
 =e^{tX}e^{tY}e^{t^2C_2}e^{t^3C_3}e^{t^4C_4}\cdots.}
\end{equation}
The infinite product is ordered by increasing $n$ from left to right and is
well defined coefficientwise in $t$: modulo $t^{N+1}$ only the factors
through $C_N$ matter. Define
\begin{equation}\label{eq:R1}
 R_1(t)=e^{-tY}e^{-tX}e^{t(X+Y)},
\end{equation}
and recursively, for $n\geq2$,
\begin{equation}\label{eq:residual}
 C_n=\coeff{t^n}R_{n-1}(t)=\coeff{t^n}\log R_{n-1}(t),\qquad
 R_n(t)=e^{-t^nC_n}R_{n-1}(t).
\end{equation}
Then $R_n(t)=I+O(t^{n+1})$ for every $n\geq1$, and \eqref{eq:residual}
computes all the exponents.
\end{theorem}
\begin{proof}
Work in $\widehat T[[t]]$, or equivalently in $\widehat T$ with $t$ tracking
total degree. The coefficient of $t^1$ in $R_1$ is $-Y-X+(X+Y)=0$, so
$R_1=I+O(t^2)$. Suppose inductively that $R_{n-1}=I+t^nC_n+O(t^{n+1})$.
Then $\log R_{n-1}=t^nC_n+O(t^{n+1})$, which gives the second expression
for $C_n$ in \eqref{eq:residual}, and
$R_n=e^{-t^nC_n}R_{n-1}=(I-t^nC_n+O(t^{2n}))(I+t^nC_n+O(t^{n+1}))
=I+O(t^{n+1})$. Thus the recursion removes one degree at a time.

Each residual is group-like in the completed Hopf algebra of
\cref{sec:hopf} (extended by the central scalar $t$): $R_1$ is a product of
exponentials of the primitive elements $-tY$, $-tX$, $t(X+Y)$, hence
group-like by \cref{lem:grouplike}; and if $R_{n-1}$ is group-like and
$C_n$ is a Lie polynomial, then $e^{-t^nC_n}$ is group-like and so is the
product $R_n$. To see that $C_n$ is a Lie polynomial, compare the
coefficients of $t^n$ in $\Delta R_{n-1}=R_{n-1}\otimes R_{n-1}$: on the
left it is $\Delta C_n$; on the right, since
$R_{n-1}=I+t^nC_n+O(t^{n+1})$ has no terms of $t$-degree between $1$ and
$n-1$, it is $C_n\otimes1+1\otimes C_n$. Hence $C_n$ is primitive, and by
\cref{thm:friedrichs} it is a Lie polynomial; it is homogeneous of degree
$n$ in the letters because every input has matching degree in $t$ and in
$X,Y$. This completes the induction.

Undoing the residuals, $R_1=e^{t^2C_2}e^{t^3C_3}\cdots e^{t^NC_N}R_N$ with
$R_N=I+O(t^{N+1})$, so modulo $t^{N+1}$ the product of the first $N-1$
factors equals $R_1$; multiplying by $e^{tX}e^{tY}$ on the left proves
\eqref{eq:zassenhaus} degree by degree. For uniqueness, suppose
$e^{t(X+Y)}=e^{tX}e^{tY}\prod_ne^{t^nC_n'}$ with homogeneous $C_n'$ of
degree $n$. Then $\prod_{n\geq2}e^{t^nC_n'}=R_1$. If $C_j'=C_j$ for
$j<n$, then multiplying on the left by the inverses of the first factors
gives $\prod_{j\geq n}e^{t^jC_j'}=R_{n-1}$, and comparing coefficients of
$t^n$ gives $C_n'=C_n$. Induction proves uniqueness.
\end{proof}
When the identity is written as $e^{-tX}e^{t(X+Y)}=\prod_{n\geq1}e^{t^nC_n}$,
as on the source page, the convention \emph{must} be $C_1=Y$; the notation
would otherwise omit a necessary first-degree factor. Extracting $t^n$
from $R_{n-1}=e^{-t^{n-1}C_{n-1}}\cdots e^{-t^2C_2}R_1$ gives the finite
coefficient recurrence
\begin{equation}\label{eq:zassfinitecoeff}
 C_n=\sum_{\substack{k_2,\ldots,k_{n-1},a,b,c\geq0\\
                 \sum_{j=2}^{n-1}jk_j+a+b+c=n}}
 \Bigl(\prod_{j=n-1}^{2}\frac{(-C_j)^{k_j}}{k_j!}\Bigr)
 \frac{(-Y)^a}{a!}\frac{(-X)^b}{b!}\frac{(X+Y)^c}{c!},
\end{equation}
in which the product over $j$ is taken in \emph{descending} order of $j$
from left to right. It involves associative products, but the result is a
Lie polynomial by the theorem, and \eqref{eq:dsw} converts it into nested
commutators.

\subsection{A recursion consisting entirely of commutators}
Define the right logarithmic derivatives $F_n(t)=R_n'(t)R_n(t)^{-1}$.
\begin{proposition}[Differential Zassenhaus algorithm]\label{prop:zassdiff}
\begin{align}
 F_1(t)&=e^{-t\ad_Y}\bigl(e^{-t\ad_X}-I\bigr)Y
 =\sum_{k\geq1}t^kf_{1,k},\qquad
 f_{1,k}=(-1)^k\sum_{j=1}^{k}\frac{\ad_Y^{k-j}\ad_X^{j}Y}{(k-j)!\,j!},
 \label{eq:F1}\\
 C_n&=\frac1n\coeff{t^{n-1}}F_{n-1}(t),\qquad
 F_n(t)=e^{-t^n\ad_{C_n}}\bigl(F_{n-1}(t)-nt^{n-1}C_n\bigr)\quad(n\geq2).
 \label{eq:Frec}
\end{align}
Writing $F_n=\sum_{k\geq0}t^kf_{n,k}$, with $f_{n,k}=0$ for $k<0$, the
coefficients satisfy the fully finite update
\begin{equation}\label{eq:frec}
 \boxed{\displaystyle
 f_{n,k}=\sum_{j=0}^{\lfloor k/n\rfloor}\frac{(-1)^j}{j!}
 \ad_{C_n}^{\,j}f_{n-1,k-nj}-nC_n\,\delta_{k,n-1},\qquad
 C_n=\frac{f_{n-1,n-1}}{n}.}
\end{equation}
Only previously computed homogeneous Lie polynomials occur; this is a
Lie-only algorithm at every degree.
\end{proposition}
\begin{proof}
For a product $ABC$ of differentiable invertible curves,
$(ABC)'(ABC)^{-1}=A'A^{-1}+A(B'B^{-1})A^{-1}+AB(C'C^{-1})B^{-1}A^{-1}$.
With $A=e^{-tY}$, $B=e^{-tX}$, $C=e^{t(X+Y)}$ we have $A'A^{-1}=-Y$,
$B'B^{-1}=-X$, $C'C^{-1}=X+Y$, and conjugations are adjoint exponentials
by \cref{thm:campbell}:
\begin{align*}
 F_1&=-Y-e^{-t\ad_Y}X+e^{-t\ad_Y}e^{-t\ad_X}(X+Y)\\
 &=-Y-e^{-t\ad_Y}X+e^{-t\ad_Y}X+e^{-t\ad_Y}e^{-t\ad_X}Y
 =e^{-t\ad_Y}\bigl(e^{-t\ad_X}-I\bigr)Y,
\end{align*}
using $e^{-t\ad_X}X=X$ and $e^{-t\ad_Y}Y=Y$. Expanding both adjoint
exponentials gives the double series
$\sum_{i\geq0,j\geq1}(-1)^{i+j}t^{i+j}\ad_Y^i\ad_X^jY/(i!j!)$, and
collecting $t^k$ gives $f_{1,k}$.

Since $R_{n-1}=I+t^nC_n+O(t^{n+1})$, we have
$R_{n-1}'=nt^{n-1}C_n+O(t^n)$ and $R_{n-1}^{-1}=I+O(t^n)$, so
$F_{n-1}=nt^{n-1}C_n+O(t^n)$; in particular $f_{n-1,k}=0$ for $k<n-1$
and $f_{n-1,n-1}=nC_n$, the first part of \eqref{eq:Frec}. Differentiating
$R_n=e^{-t^nC_n}R_{n-1}$,
\[
 F_n=\Bigl(\frac{\dd}{\dd t}e^{-t^nC_n}\Bigr)e^{t^nC_n}
 +e^{-t^nC_n}F_{n-1}e^{t^nC_n}
 =-nt^{n-1}C_n+e^{-t^n\ad_{C_n}}F_{n-1},
\]
because $t^nC_n$ commutes with its derivative and conjugation by
$e^{-t^nC_n}$ is $e^{-t^n\ad_{C_n}}$. As $e^{-t^n\ad_{C_n}}C_n=C_n$, the
correction may be moved inside the operator exponential, giving the second
part of \eqref{eq:Frec}. Expanding
$e^{-t^n\ad_{C_n}}=\sum_j(-1)^jt^{nj}\ad_{C_n}^j/j!$ and extracting the
coefficient of $t^k$ gives \eqref{eq:frec}; the term $-nt^{n-1}C_n$ is
annihilated by $\ad_{C_n}^j$ for $j\geq1$ and contributes only through
$j=0$.
\end{proof}

\subsection{The exponents displayed on the source page}
\begin{proposition}\label{prop:zassdisplayed}
\begin{align}
 C_2&=-\tfrac12[X,Y],\label{eq:c2}\\
 C_3&=\tfrac16[X,[X,Y]]+\tfrac13[Y,[X,Y]]
     =\tfrac16[X,[X,Y]]-\tfrac13[[X,Y],Y],\label{eq:c3}\\
 C_4&=-\tfrac1{24}[X,[X,[X,Y]]]-\tfrac18[Y,[X,[X,Y]]]-\tfrac18[Y,[Y,[X,Y]]]
 \label{eq:c4}\\
    &=-\tfrac1{24}\bigl([[[X,Y],X],X]+3[[[X,Y],X],Y]+3[[[X,Y],Y],Y]\bigr).
 \label{eq:c4left}
\end{align}
\end{proposition}
\begin{proof}
From \eqref{eq:F1},
\begin{align*}
 f_{1,1}&=-[X,Y],\\
 f_{1,2}&=\ad_Y\ad_XY+\tfrac12\ad_X^2Y=[Y,[X,Y]]+\tfrac12[X,[X,Y]],\\
 f_{1,3}&=-\tfrac12\ad_Y^2\ad_XY-\tfrac12\ad_Y\ad_X^2Y-\tfrac16\ad_X^3Y.
\end{align*}
By \eqref{eq:frec}, $C_2=f_{1,1}/2=-\tfrac12[X,Y]$. Next,
$f_{2,2}=f_{1,2}-\ad_{C_2}f_{1,0}=f_{1,2}$ since $f_{1,0}=0$, so
$C_3=f_{2,2}/3$, which is \eqref{eq:c3}; the second form uses
$[Y,[X,Y]]=-[[X,Y],Y]$. Then
$f_{2,3}=f_{1,3}-\ad_{C_2}f_{1,1}=f_{1,3}+\tfrac12[[X,Y],[X,Y]]=f_{1,3}$,
and $f_{3,3}=f_{2,3}-\ad_{C_3}f_{2,0}=f_{2,3}$, so $C_4=f_{1,3}/4$, which is
\eqref{eq:c4}. For \eqref{eq:c4left}, antisymmetry gives
$[[X,Y],X]=-[X,[X,Y]]$, hence $[[[X,Y],X],X]=[X,[X,[X,Y]]]$,
$[[[X,Y],X],Y]=[Y,[X,[X,Y]]]$, and $[[[X,Y],Y],Y]=[Y,[Y,[X,Y]]]$.
Substituting these into \eqref{eq:c4} gives \eqref{eq:c4left}, exactly the
left-nested form on the source page. Since the corrections at each step
were computed exactly, no higher-order approximation has been left implicit.
\end{proof}
The complete fifth and sixth exponents are printed in \cref{app:tables}
(\cref{tab:zass56}), and all exponents through degree twelve accompany the
article as data. There is also a left-oriented companion factorization:
replacing $t$ by $-t$ in \eqref{eq:zassenhaus} and inverting both sides
reverses the factor order and gives
\begin{equation}\label{eq:zassleft}
 e^{t(X+Y)}=\cdots e^{t^4\widehat C_4}e^{t^3\widehat C_3}e^{t^2\widehat C_2}e^{tY}e^{tX},
 \qquad \widehat C_n=(-1)^{n+1}C_n,
\end{equation}
since $(e^{(-t)^nC_n})^{-1}=e^{(-1)^{n+1}t^nC_n}$.

\subsection{A finite nonrecursive formula using weighted trees}
The previous algorithms determine the exponents recursively. For a finite
nonrecursive expression at every fixed order, first define the explicit
associative polynomials
\begin{equation}\label{eq:An}
 A_n=\coeff{t^n}R_1(t)=\sum_{a+b+c=n}\frac{(-1)^{a+b}}{a!\,b!\,c!}\,Y^aX^b(X+Y)^c
 \qquad(n\geq0),
\end{equation}
where $(X+Y)^c$ is the sum of all $2^c$ words of length $c$; $A_0=I$ and
$A_1=0$. Coefficient extraction from the ordered product
$R_1=\prod_{j\geq2}e^{t^jC_j}$ gives the forward partition recurrence
\begin{equation}\label{eq:Cpartitions}
 C_n=A_n-\sum_{\substack{k_2,\ldots,k_{n-1}\geq0\\\sum_{j=2}^{n-1}jk_j=n}}
 \frac{C_2^{k_2}C_3^{k_3}\cdots C_{n-1}^{k_{n-1}}}{k_2!\,k_3!\cdots k_{n-1}!},
\end{equation}
with the factors in \emph{increasing} order of degree. Indeed, choosing
from each factor $e^{t^jC_j}$ the term $t^{jk_j}C_j^{k_j}/k_j!$ with
$\sum jk_j=n$ gives all contributions to $t^n$; the unique choice with
$k_n=1$ contributes $C_n$, and every other choice has $k_n=0$ and at
least two factors. The sum is empty for $n=2,3$.

\begin{definition}\label{def:trees}
Let $\cT_n$ be the set of finite plane rooted trees with a positive integer
weight at each vertex, subject to: the root has weight $n$; every weight is
at least $2$; and each internal vertex $v$ of weight $n_v$ has $r\geq2$
ordered children whose weights $2\leq j_1\leq j_2\leq\cdots\leq j_r$ are
weakly increasing from left to right and sum to $n_v$. Leaves carry no
further restriction. Let $I(T)$ be the number of internal vertices, let
$m_{v,j}$ be the number of children of weight $j$ at the internal vertex
$v$, and let $\ell_1,\ldots,\ell_q$ be the weights of the leaves read from
left to right in planar order. Define
\begin{equation}\label{eq:treeweight}
 \omega(T)=\frac{(-1)^{I(T)}}{\prod_{v\ \mathrm{internal}}\prod_{j\geq2}m_{v,j}!},
 \qquad
 A(T)=A_{\ell_1}A_{\ell_2}\cdots A_{\ell_q}.
\end{equation}
Children of equal weight are distinct ordered slots whose subtrees may
differ; interchanging two unequal subtrees changes $A(T)$ in general,
because the $A_j$ do not commute.
\end{definition}

\begin{theorem}[All-terms Zassenhaus tree formula]\label{thm:tree}
For every $n\geq2$, $\cT_n$ is finite and
\begin{equation}\label{eq:treeformula}
 \boxed{\displaystyle C_n=\sum_{T\in\cT_n}\omega(T)\,A(T).}
\end{equation}
This is a finite formula involving only factorials, sums, and products of
$X$ and $Y$. An explicitly nested-commutator version is
\begin{equation}\label{eq:treelie}
 C_n=\frac1n\sum_{T\in\cT_n}\omega(T)\sum_{w\in\{X,Y\}^n}\coeff{w}A(T)\,R(w).
\end{equation}
\end{theorem}
\begin{proof}
Finiteness: the leaf weights are at least $2$ and sum to $n$ (every
internal vertex's weight is the sum of its children's weights, so by
induction the root weight is the sum of the leaf weights), so there are at
most $\lfloor n/2\rfloor$ leaves; since every internal vertex has at least
two children, a plane tree with $q$ leaves has fewer than $q$ internal
vertices; and each vertex weight is an integer in $[2,n]$. Hence only
finitely many shapes and weight assignments occur.

We prove \eqref{eq:treeformula} by strong induction on $n$. For $n=2,3$
the sum in \eqref{eq:Cpartitions} is empty, $C_n=A_n$, and $\cT_n$ consists
of the single-vertex tree (a root of weight $n$ with no children, since
no weakly increasing list of at least two integers $\geq2$ sums to $2$ or
$3$), for which $\omega=1$ and $A(T)=A_n$. For $n\geq4$, assume the formula
for all degrees below $n$. In \eqref{eq:Cpartitions}, a choice of exponents
$(k_2,\ldots,k_{n-1})$ with $\sum jk_j=n$ corresponds bijectively to the
weakly increasing list of child weights of a root of weight $n$, namely
$k_2$ copies of $2$, then $k_3$ copies of $3$, and so on; the product
$C_2^{k_2}\cdots C_{n-1}^{k_{n-1}}$ is the ordered product of $C_{j_i}$
over this list, and the coefficient $-1/\prod_jk_j!$ is exactly the root's
contribution $(-1)\prod_j1/m_{\mathrm{root},j}!$ to $\omega$. Substituting
the induction hypothesis $C_{j_i}=\sum_{T_i\in\cT_{j_i}}\omega(T_i)A(T_i)$
for each factor and expanding the product gives a sum over all tuples
$(T_1,\ldots,T_r)$ of subtrees, one in each child slot, of
$(-1)\prod_j\frac1{m_{\mathrm{root},j}!}\prod_i\omega(T_i)\prod_iA(T_i)$.
Attaching $T_1,\ldots,T_r$ in order to the root produces exactly one tree
$T\in\cT_n$ with at least one internal vertex; conversely every such tree
arises exactly once (remove the root). For that tree,
$\omega(T)=(-1)\prod_j\frac1{m_{\mathrm{root},j}!}\prod_i\omega(T_i)$
(the internal vertices of $T$ are the root and those of the $T_i$), and
$A(T)=\prod_iA(T_i)$ because the leaves of $T$ in planar order are the
leaves of $T_1$, then of $T_2$, and so on. The remaining term $A_n$ is the
contribution of the single-vertex tree. This proves \eqref{eq:treeformula}.
Finally, $C_n$ is a homogeneous Lie polynomial of degree $n$ by
\cref{thm:zassenhaus}, hence primitive, so \cref{lem:dsw} gives
$C_n=\frac1nD(C_n)=\frac1n\sum_w\coeff{w}C_n\,R(w)$; substituting the word
expansion of \eqref{eq:treeformula} gives \eqref{eq:treelie}.
\end{proof}
For orientation, the first instances are
\begin{equation}\label{eq:CinA}
\begin{aligned}
 C_2&=A_2,\qquad C_3=A_3,\qquad C_4=A_4-\tfrac12A_2^2,\qquad
 C_5=A_5-A_2A_3,\\
 C_6&=A_6-A_2A_4-\tfrac12A_3^2+\tfrac13A_2^3,\qquad
 C_7=A_7-A_2A_5-A_3A_4+\tfrac12A_2^2A_3+\tfrac12A_3A_2^2.
\end{aligned}
\end{equation}
The coefficient $\tfrac13$ of $A_2^3$ in $C_6$ is
$\tfrac12-\tfrac16$: the tree with root children $(2,4)$ and the child of
weight $4$ split into $(2,2)$ contributes $(-1)^2\cdot\frac1{2!}=\frac12$,
and the tree with root children $(2,2,2)$ contributes $-\frac1{3!}$. In
$C_7$ the two terms $A_2^2A_3$ and $A_3A_2^2$ are different polynomials:
the first collects the tree with root list $(2,5)$ and $5\to(2,3)$
(weight $+1$) and the tree with root list $(2,2,3)$ (weight $-\tfrac1{2!}$),
while the second comes from the root list $(3,4)$ with $4\to(2,2)$
(weight $+\tfrac1{2!}$). This illustrates why an unordered product would
lose information. The data file \texttt{zassenhaus\_tree\_polynomials.json}
lists these polynomials in the $A_j$ through degree ten, and the
verification program of \cref{sec:computation} confirms that
\eqref{eq:treeformula} agrees with the two recursive constructions through
that degree.

\subsection{A globally convergent ordered-simplex expansion}
There is another useful all-terms expansion, a sum of ordered products
rather than a product of exponentials; its convergence is therefore a
separate and easier question than Zassenhaus convergence.
\begin{theorem}[Ordered-simplex expansion]\label{thm:simplex}
For bounded $X,Y$ in a unital Banach algebra and every scalar $t$,
\begin{equation}\label{eq:simplex}
\begin{split}
 e^{t(X+Y)}=e^{tX}\Biggl[I+\sum_{k\geq1}\ \sum_{n_1,\ldots,n_k\geq0}
 &\frac{(-1)^{n_1+\cdots+n_k}\,t^{k+n_1+\cdots+n_k}}
 {\displaystyle\prod_{j=1}^{k}n_j!\ \prod_{j=1}^{k}\bigl(n_j+n_{j+1}+\cdots+n_k+k-j+1\bigr)}\\[-4pt]
 &\qquad\cdot\bigl(\ad_X^{n_1}Y\bigr)\bigl(\ad_X^{n_2}Y\bigr)\cdots\bigl(\ad_X^{n_k}Y\bigr)\Biggr],
\end{split}
\end{equation}
and the multiple series is absolutely convergent, with the sum of the norms
of its terms bounded by $\exp\bigl(|t|\norm Ye^{2|t|\norm X}\bigr)$.
\end{theorem}
\begin{proof}
Put $H(t)=e^{-tX}e^{t(X+Y)}$. Differentiating,
$H'(t)=-Xe^{-tX}e^{t(X+Y)}+e^{-tX}(X+Y)e^{t(X+Y)}=e^{-tX}Ye^{t(X+Y)}
=\bigl(e^{-tX}Ye^{tX}\bigr)H(t)=B(t)H(t)$ with $B(t)=e^{-t\ad_X}Y$ by
\cref{thm:campbell}, and $H(0)=I$. First take real $t\geq0$. Integrating,
$H(t)=I+\int_0^tB(s)H(s)\dd s$, and substituting this expression for $H(s)$
repeatedly $N$ times gives
\begin{align*}
 H(t)=I&+\sum_{k=1}^{N}\int_{0<t_k<\cdots<t_1<t}B(t_1)\cdots B(t_k)\,\dd t_k\cdots\dd t_1\\
 &+\int_{0<t_{N+1}<\cdots<t_1<t}B(t_1)\cdots B(t_{N+1})H(t_{N+1})\,\dd t_{N+1}\cdots\dd t_1.
\end{align*}
With $b=\sup_{0\leq s\leq t}\norm{B(s)}\leq\norm Ye^{2t\norm X}$ and
$\norm{H(s)}\leq\norm I^2e^{s\norm X+s\norm{X+Y}}$ bounded on $[0,t]$, the
remainder has norm at most $b^{N+1}t^{N+1}/(N+1)!$ times that bound, since
the ordered simplex has volume $t^{N+1}/(N+1)!$. Hence the remainder tends
to zero and $H(t)=I+\sum_{k\geq1}\int_{\text{simplex}}B(t_1)\cdots B(t_k)$.
Now expand $B(s)=\sum_{n\geq0}(-s)^n\ad_X^nY/n!$ in each factor. The
resulting multiple series of monomials is absolutely convergent: the sum of
the norms of the terms of $B(s)$ is at most
$\sum_n s^n(2\norm X)^n\norm Y/n!=\norm Ye^{2s\norm X}\leq b$, so the sum
of the norms of all monomials in the $k$th ordered integral is at most
$b^kt^k/k!$, and summing over $k$ gives the bound $e^{bt}$. Absolute
convergence justifies integrating term by term. The iterated monomial
integral is evaluated from the inside out: integrating $t_k^{n_k}$ from $0$
to $t_{k-1}$ gives $t_{k-1}^{n_k+1}/(n_k+1)$; then
$t_{k-1}^{n_{k-1}+n_k+1}$ integrates to $t_{k-2}^{n_{k-1}+n_k+2}/(n_{k-1}+n_k+2)$;
continuing,
\[
 \int_{0<t_k<\cdots<t_1<t}\prod_{j=1}^kt_j^{n_j}\,\dd t_k\cdots\dd t_1
 =\frac{t^{k+n_1+\cdots+n_k}}{\prod_{j=1}^{k}(n_j+n_{j+1}+\cdots+n_k+k-j+1)}.
\]
This proves the coefficients in \eqref{eq:simplex} for real $t\geq0$.
For complex $t$, parametrize the segment from $0$ to $t$ by $s=tu$,
$0\leq u\leq1$; the same Picard iteration and the same integrals, now with
the factor $t^{k+\sum n_j}$ pulled out, give the identical formula and the
bound with $|t|$. Multiplying by $e^{tX}$ gives \eqref{eq:simplex}.
\end{proof}

\subsection{An elementary sufficient convergence condition for the product}
Formal stabilization of a product does not automatically imply analytic
convergence. The following bound supplies a concrete domain; it is
deliberately elementary and is not asserted to be optimal.
\begin{theorem}[Convergence of the ordered Zassenhaus product]\label{thm:zconv}
Let $X,Y$ be elements of a unital Banach algebra with $\norm I=1$ and put
$a=\norm X+\norm Y$. Let $r>0$ satisfy
\begin{equation}\label{eq:zconv}
 h(r):=e^{2ar}-1-2ar<2\log2-1.
\end{equation}
Then $\sum_{n\geq2}r^n\norm{C_n}<\infty$, the ordered product in
\eqref{eq:zassenhaus} converges in norm uniformly for $|t|\leq r$, and the
identity \eqref{eq:zassenhaus} holds there. Condition \eqref{eq:zconv} is
satisfied whenever $ar<u_*/2=0.38349\ldots$, where $u_*=0.76699\ldots$ is
the positive solution of $e^{u}-1-u=2\log2-1$; a particularly simple
sufficient condition is $ar\leq\tfrac14$.
\end{theorem}
\begin{proof}
\emph{Majorant.} Expanding the three exponentials in \eqref{eq:R1}, the
sum of the norms of all terms of degree $n$ in $t$ is at most the
coefficient of $t^n$ in $e^{t\norm Y}e^{t\norm X}e^{t(\norm X+\norm Y)}
=e^{2at}$, so $\norm{A_n}\leq(2a)^n/n!$ for $n\geq2$, and
$h(z)=\sum_{n\geq2}(2a)^nz^n/n!$ is a nonnegative-coefficient majorant of
$R_1-I$. Define a scalar formal series $c(z)=\sum_{n\geq2}c_nz^n$ by
\begin{equation}\label{eq:zmajorant}
 c=h+\bigl(e^{c}-1-c\bigr).
\end{equation}
Since $c$ has order $2$, the expression $e^c-1-c=\sum_{m\geq2}c^m/m!$ has
its degree-$n$ coefficient determined by $c_2,\ldots,c_{n-2}$; so
\eqref{eq:zmajorant} determines the $c_n$ recursively, and they are
nonnegative because all ingredients are. In the commutative scalar algebra,
$\prod_{j\geq2}e^{c_jz^j}=e^{c(z)}$, whose coefficient of $z^n$ is
$\sum_{\sum jk_j=n}\prod_jc_j^{k_j}/k_j!$; the choice $k_n=1$ contributes
$c_n$, and the remaining choices, all with at least two factors, contribute
exactly $\coeff{z^n}(e^c-1-c)$. Comparing with \eqref{eq:Cpartitions} and
using the triangle inequality and submultiplicativity,
\[
 \norm{C_n}\leq\norm{A_n}+\sum_{\substack{\sum jk_j=n\\k_n=0}}
 \prod_j\frac{\norm{C_j}^{k_j}}{k_j!}
 \leq h_n+\sum_{\substack{\sum jk_j=n\\k_n=0}}\prod_j\frac{c_j^{k_j}}{k_j!}=c_n,
\]
where the middle inequality uses the induction hypothesis
$\norm{C_j}\leq c_j$ for $j<n$ and monotonicity in the $c_j\geq0$.

\emph{Convergence of the majorant.} The function $q(d)=2d-e^d+1$ is
strictly increasing on $[0,\log2)$, since $q'(d)=2-e^d>0$ there, from
$q(0)=0$ to $q(\log2)=2\log2-1$. Since $h$ is increasing and continuous and
\eqref{eq:zconv} is strict, choose $R>r$ with $h(R)<2\log2-1$, and then
$d\in[0,\log2)$ with $q(d)=h(R)$. Let $\mathcal B$ be the Banach space of
functions holomorphic on $|z|<R$ and continuous on $|z|\leq R$ with the
supremum norm, and let $\mathcal B_d$ be its closed ball of radius $d$.
For $f\in\mathcal B_d$ define $\mathcal F(f)(z)=h(z)+e^{f(z)}-1-f(z)$.
Since $h$ has nonnegative coefficients, $|h(z)|\leq h(R)$ on the closed
disk, and $|e^{f}-1-f|\leq\sum_{m\geq2}|f|^m/m!\leq e^d-1-d$; hence
$\norm{\mathcal F(f)}\leq h(R)+e^d-1-d=q(d)+e^d-1-d=d$, so $\mathcal F$
maps $\mathcal B_d$ into itself. For $f,g\in\mathcal B_d$,
$e^{f}-f-(e^{g}-g)=\int_0^1(e^{g+\tau(f-g)}-1)\,\dd\tau\,(f-g)$, so
$\norm{\mathcal F(f)-\mathcal F(g)}\leq(e^d-1)\norm{f-g}$ with
$e^d-1<1$. By the contraction mapping theorem $\mathcal F$ has a unique
fixed point $f_*\in\mathcal B_d$. Its Taylor coefficients satisfy
\eqref{eq:zmajorant} (coefficient extraction is continuous and
$e^{f_*}-1-f_*$ has the same coefficient recursion), $f_*$ has order at
least $2$ because $h$ does and $e^{f}-1-f$ has order at least twice that of
$f$, and the formal solution of \eqref{eq:zmajorant} is unique; hence
$f_*(z)=\sum_nc_nz^n$ on $|z|<R$. As $c_n\geq0$ and $r<R$,
\[
 \sum_{n\geq2}r^n\norm{C_n}\leq\sum_{n\geq2}c_nr^n=f_*(r)\leq d<\log2.
\]

\emph{Convergence of the product.} Let $P_N(t)=\prod_{n=2}^Ne^{t^nC_n}$.
For $|t|\leq r$ and $M>N$, $P_M-P_N=P_N\bigl(\prod_{n=N+1}^Me^{t^nC_n}-I\bigr)$.
Expanding each factor as $I+(e^{t^nC_n}-I)$ and multiplying out, every term
other than $I$ contains at least one factor $e^{t^nC_n}-I$ of norm at most
$e^{r^n\norm{C_n}}-1$, so
$\norm{\prod_{n=N+1}^Me^{t^nC_n}-I}\leq\exp\bigl(\sum_{n>N}r^n\norm{C_n}\bigr)-1$,
which tends to zero as $N\to\infty$; and $\norm{P_N}\leq\exp(\sum_nr^n\norm{C_n})\leq e^d$.
Hence $(P_N)$ is uniformly Cauchy on $|t|\leq r$ and converges uniformly to
a limit $P(t)$, continuous on the closed disk and analytic on $|t|<r$ as a
uniform limit of entire functions. For $n\leq N$ the coefficient of $t^n$ in
$P_N$ does not depend on $N$ (the omitted factors contribute only degrees
above $N$), and by the Cauchy coefficient formula uniform convergence gives
$\coeff{t^n}P=\coeff{t^n}P_N$. The formal identity of \cref{thm:zassenhaus}
says that $\coeff{t^n}\bigl(e^{tX}e^{tY}P_N(t)\bigr)=\coeff{t^n}e^{t(X+Y)}$
for $n\leq N$. Thus the analytic functions $e^{tX}e^{tY}P(t)$ and
$e^{t(X+Y)}$ have the same Taylor coefficients at $0$, so they agree on
$|t|<r$, and by continuity on $|t|\leq r$.

\emph{The numerical conditions.} The function $u\mapsto e^u-1-u$ is
strictly increasing on $[0,\infty)$ from $0$, so $e^{u}-1-u=2\log2-1
=0.386294\ldots$ has a unique positive solution $u_*$; evaluating,
$e^{0.767}-1-0.767=0.38629\ldots$ and $e^{0.766}-1-0.766=0.38514\ldots$, so
$u_*=0.7669980\ldots$ and \eqref{eq:zconv} holds exactly when
$2ar<u_*$. In particular $ar\leq\tfrac14$ gives
$h(r)\leq e^{1/2}-\tfrac32<1.6488-1.5=0.1488<0.386$.
\end{proof}
The assumption $\norm I=1$ was used only in the bound $\norm{P_N}\leq e^d$;
for a general submultiplicative norm a fixed factor $\norm I$ appears and
nothing else changes.

\subsection{Entire resummation in the eigenvector case}
When $[X,Y]=sY$, the factorization \eqref{eq:eigenfactor} with $c=1$ gives
$e^{t(X+Y)}=e^{tX}e^{q_s(t)Y}$, hence
\[
 R_1(t)=e^{-tY}e^{-tX}e^{tX}e^{q_s(t)Y}=e^{(q_s(t)-t)Y},\qquad
 q_s(t)-t=\sum_{n\geq2}\frac{(-1)^{n-1}s^{n-1}}{n!}t^n.
\]
All exponents are multiples of the same $Y$ and commute, so by
\cref{lem:explog}
\[
 e^{(q_s(t)-t)Y}=\prod_{n\geq2}\exp\Bigl(\frac{(-1)^{n-1}s^{n-1}}{n!}\,t^nY\Bigr),
\]
and uniqueness in \cref{thm:zassenhaus} yields the all-orders closed form
\begin{equation}\label{eq:eigenzass}
 \boxed{\displaystyle C_n=\frac{(-1)^{n-1}s^{n-1}}{n!}\,Y\qquad(n\geq2).}
\end{equation}
This product converges for every finite $s$ and $t$, a useful contrast with
the Bernoulli BCH series \eqref{eq:eigenseries}, whose radius in $s$ is only
$2\pi$. If instead the commutator $C=[X,Y]$ is central, \eqref{eq:centralt}
gives $e^{t(X+Y)}=e^{tX}e^{tY}e^{-t^2C/2}$, so $C_2=-\tfrac12[X,Y]$ and
$C_n=0$ for every $n\geq3$.

\section{Trotter, symmetric, and higher-order product formulas with complete error series}\label{sec:splitting}
The source page mentions the Suzuki--Trotter decomposition as a consequence
of the Zassenhaus formula without displaying a particular version. We prove
the basic product limit, the symmetric second-order formula, and a
recursive construction of arbitrarily high even order, each with its
complete logarithmic error series. All analytic statements in this section
concern bounded elements $X,Y$ of a unital Banach algebra with $\norm I=1$;
they do not assert a product theorem for unbounded semigroup generators,
for which the classical sources \cite{trotter} require additional
hypotheses. Throughout $a=\norm X+\norm Y$ and $S=X+Y$.

We first record the multifactor version of \cref{thm:convergence}(i):
for bounded $A_1,\ldots,A_m$ with $\sum_i\norm{A_i}=\sigma<\log2$, the
Mercator logarithm of $\prod_ie^{A_i}$ converges absolutely, equals the
sum of the homogeneous parts of \eqref{eq:manyexp}, exponentiates to the
product, and the sum of the norms of its terms of degree at least $N+1$ is
at most $\rho^{-(N+1)}\bigl[-\log(2-e^{\rho\sigma})\bigr]$ for any
$\rho>1$ with $\rho\sigma<\log2$. The proof is identical: the sum of the
norms of the nonempty words of $\prod e^{A_i}-I$ is at most $e^{\sigma}-1$.

\begin{lemma}\label{lem:tail}
If $\sigma<\tfrac12\log2$, the terms of degree at least $2$ in the
logarithm of $\prod_ie^{A_i}$ have total norm at most $c_0\sigma^2$, and
those of degree at least $3$ have total norm at most $c_0'\sigma^3$, with
absolute constants $c_0=4(\log2)^{-2}\log\frac1{2-\sqrt2}$ and
$c_0'=8(\log2)^{-3}\log\frac1{2-\sqrt2}$.
\end{lemma}
\begin{proof}
Take $\rho=\log2/(2\sigma)>1$, so $e^{\rho\sigma}=\sqrt2$ and
$-\log(2-e^{\rho\sigma})=\log\frac1{2-\sqrt2}$. The tail bound with $N=1$
gives $\rho^{-2}\log\frac1{2-\sqrt2}=\frac{4\sigma^2}{(\log2)^2}\log\frac1{2-\sqrt2}$,
and with $N=2$ it gives the cubic bound.
\end{proof}

\subsection{The Lie--Trotter product formula}
\begin{theorem}[Lie--Trotter formula and its complete logarithmic error]\label{thm:trotter}
For every fixed scalar $t$,
\begin{equation}\label{eq:trotter}
 \lim_{n\to\infty}\bigl(e^{tX/n}e^{tY/n}\bigr)^n=e^{t(X+Y)}
\end{equation}
in norm, with error $O(n^{-1})$ uniformly for $t$ in compact sets. For all
$n>|t|a/\log2$ one has the exact convergent expansion
\begin{equation}\label{eq:trotterall}
 \bigl(e^{tX/n}e^{tY/n}\bigr)^n
 =\exp\Bigl(t(X+Y)+\sum_{k\geq2}\frac{t^k}{n^{k-1}}Z_k(X,Y)\Bigr),
\end{equation}
so every term of the logarithmic error is explicitly available from
\eqref{eq:wordcut} or Dynkin's formula.
\end{theorem}
\begin{proof}
For $n>|t|a/\log2$ the pair $(tX/n,tY/n)$ satisfies the hypothesis of
\cref{thm:convergence}(i), so $L_n:=\log(e^{tX/n}e^{tY/n})=\sum_kZ_k(tX/n,tY/n)
=\sum_k(t/n)^kZ_k(X,Y)$ converges absolutely, using homogeneity of $Z_k$, and
$e^{L_n}=e^{tX/n}e^{tY/n}$. Powers of a single exponential satisfy
$(e^{L})^n=e^{nL}$ by \cref{lem:explog} (in its analytic form), so the
$n$th power equals $\exp(nL_n)=\exp\bigl(tS+\sum_{k\geq2}t^kn^{1-k}Z_k\bigr)$,
which is \eqref{eq:trotterall}. Let $E_n=\sum_{k\geq2}t^kn^{1-k}Z_k
=n\sum_{k\geq2}Z_k(tX/n,tY/n)$. For $n>2|t|a/\log2$, \cref{lem:tail}
with $\sigma=|t|a/n$ gives $\norm{E_n}\leq nc_0(|t|a/n)^2=c_0t^2a^2/n$.
By the exponential difference estimate \eqref{eq:expdifference} with
$\norm I=1$,
\[
 \bigl\|(e^{tX/n}e^{tY/n})^n-e^{tS}\bigr\|=\norm{e^{tS+E_n}-e^{tS}}
 \leq e^{|t|a+\norm{E_n}}\norm{E_n}\leq e^{|t|a+c_0t^2a^2/n}\,\frac{c_0t^2a^2}{n},
\]
which is $O(n^{-1})$ uniformly for $|t|$ bounded.
\end{proof}
Thus the Zassenhaus viewpoint explains which correction factors are
discarded, while BCH provides the most direct proof of the limit and of
all corrections to its effective exponent.

\subsection{Symmetry removes every even degree}
Put
\begin{equation}\label{eq:S2}
 S_2(t)=e^{tX/2}e^{tY}e^{tX/2}.
\end{equation}
\begin{proposition}[Symmetric splitting]\label{prop:symmetric}
The formal logarithm of $S_2(t)$ is odd in $t$:
\begin{equation}\label{eq:S2log}
 \log S_2(t)=t(X+Y)+\sum_{j\geq1}t^{2j+1}E_{2j+1},\qquad
 E_3=-\frac1{24}[X,[X,Y]]-\frac1{12}[Y,[X,Y]].
\end{equation}
Every $E_{2j+1}$ is a Lie polynomial given by the finite multifactor
word-cut formula \eqref{eq:manyexp} with $(A_1,A_2,A_3)=(X/2,Y,X/2)$,
followed by the Dynkin projection. For bounded $X,Y$ and $n>|t|a/\log2$,
\begin{equation}\label{eq:S2all}
 S_2(t/n)^n=\exp\Bigl(t(X+Y)+\sum_{j\geq1}\frac{t^{2j+1}}{n^{2j}}E_{2j+1}\Bigr)
 =e^{t(X+Y)}+O(n^{-2}),
\end{equation}
uniformly for $t$ in compact sets.
\end{proposition}
\begin{proof}
Since $(e^{A})^{-1}=e^{-A}$ and inversion reverses products,
$S_2(-t)=e^{-tX/2}e^{-tY}e^{-tX/2}=S_2(t)^{-1}$. Let $L(t)=\log S_2(t)$ be
the formal logarithm, a series in $\widehat T$ whose degree-$n$ part is
$t^n$ times a polynomial. Then $e^{L(-t)}=S_2(-t)=S_2(t)^{-1}=e^{-L(t)}$,
and uniqueness of the formal logarithm gives $L(-t)=-L(t)$; hence only odd
powers of $t$ occur. By \cref{cor:manyexp} and \cref{thm:friedrichs}
applied to the three-letter alphabet (the product of exponentials of
primitives is group-like), each homogeneous part is a Lie polynomial given
by the finite cut sum.

For $E_3$ we compute by two applications of BCH through degree three. Put
$A=tX/2$ and $B=tY$, so that $S_2=e^Ae^Be^A$. By
\eqref{eq:z1}--\eqref{eq:z3},
\[
 W:=\BCH(A,B)=A+B+\tfrac12[A,B]+\tfrac1{12}[A,[A,B]]+\tfrac1{12}[B,[B,A]]+O(t^4).
\]
Then $\log S_2=\BCH(W,A)=W+A+\tfrac12[W,A]+\tfrac1{12}[W,[W,A]]+\tfrac1{12}[A,[A,W]]+O(t^4)$.
Writing $W=W_1+W_2+W_3+O(t^4)$ with $W_1=A+B$, $W_2=\tfrac12[A,B]$,
$W_3=\tfrac1{12}[A,[A,B]]+\tfrac1{12}[B,[B,A]]$, the degree-two part is
$W_2+\tfrac12[W_1,A]=\tfrac12[A,B]+\tfrac12[B,A]=0$, as oddness predicts.
The degree-three part is
\[
 W_3+\tfrac12[W_2,A]+\tfrac1{12}[W_1,[W_1,A]]+\tfrac1{12}[A,[A,W_1]],
\]
and the four pieces are, respectively,
$\tfrac1{12}[A,[A,B]]+\tfrac1{12}[B,[B,A]]$;
$\tfrac14[[A,B],A]=-\tfrac14[A,[A,B]]$;
$\tfrac1{12}[A+B,[B,A]]=-\tfrac1{12}[A,[A,B]]+\tfrac1{12}[B,[B,A]]$; and
$\tfrac1{12}[A,[A,B]]$. Collecting, the coefficient of $[A,[A,B]]$ is
$\tfrac1{12}-\tfrac14-\tfrac1{12}+\tfrac1{12}=-\tfrac16$ and that of
$[B,[B,A]]$ is $\tfrac1{12}+\tfrac1{12}=\tfrac16$. Substituting
$A=tX/2$, $B=tY$: $-\tfrac16\cdot\tfrac{t^3}4[X,[X,Y]]+\tfrac16\cdot\tfrac{t^3}2[Y,[Y,X]]
=-\tfrac{t^3}{24}[X,[X,Y]]-\tfrac{t^3}{12}[Y,[X,Y]]$, which is $E_3$.

For \eqref{eq:S2all}, apply the multifactor convergence statement to the
three factors of $S_2(t/n)$, whose norms sum to $|t|a/n$; the logarithm
converges for $n>|t|a/\log2$, and raising to the $n$th power multiplies
it by $n$, giving the displayed exponent by homogeneity. The error
exponent $E_n'=n\sum_{j\geq1}(t/n)^{2j+1}E_{2j+1}$ consists of the terms of
degree at least $3$ of $n\log S_2(t/n)$, so \cref{lem:tail} gives
$\norm{E_n'}\leq nc_0'(|t|a/n)^3=O(n^{-2})$ for $n>2|t|a/\log2$, and the
exponential difference estimate finishes the proof exactly as in
\cref{thm:trotter}.
\end{proof}
If $[X,Y]$ is central, \eqref{eq:centralbch} shows that already
$S_2(t)=e^{t(X+Y)}$ exactly.

\subsection{Arbitrarily high even order by a five-factor recursion}
The recursive symmetric construction below is Suzuki's fractal
decomposition \cite{suzuki1990,suzuki1991}; its order conditions follow
directly from BCH.
\begin{theorem}[Symmetric product formulas of every even order]\label{thm:suzuki}
For $k\geq2$ set
\begin{equation}\label{eq:suzukip}
 p_k=\frac1{4-4^{1/(2k-1)}},
\end{equation}
and define recursively, starting from \eqref{eq:S2},
\begin{equation}\label{eq:suzuki}
 S_{2k}(t)=S_{2k-2}(p_kt)^2\,S_{2k-2}\bigl((1-4p_k)t\bigr)\,S_{2k-2}(p_kt)^2.
\end{equation}
Then $S_{2k}(-t)=S_{2k}(t)^{-1}$ and
\begin{equation}\label{eq:suzukiorder}
 \log S_{2k}(t)=t(X+Y)+O(t^{2k+1}),\qquad
 S_{2k}(t/n)^n=e^{t(X+Y)}+O(n^{-2k})
\end{equation}
for fixed $k$ and $t$; the first statement is formal and locally analytic,
the second is a norm statement for bounded $X,Y$. Every term of the local
logarithm of $S_{2k}$ is a Lie polynomial with an explicit finite
multifactor word-cut expansion.
\end{theorem}
\begin{proof}
We prove by induction on $k\geq1$ that $S_{2k}(-t)=S_{2k}(t)^{-1}$ and
$\log S_{2k}(t)=tS+t^{2k+1}E^{(k)}+O(t^{2k+3})$ for some Lie polynomial
$E^{(k)}$, where the logarithm is the formal one. The case $k=1$ is
\cref{prop:symmetric}. Assume the statement for $k-1\geq1$ and abbreviate
$p=p_k$, $c=(p,p,1-4p,p,p)$. Symmetry: since inversion reverses the order
of a product and $S_{2k-2}(-ct)=S_{2k-2}(ct)^{-1}$,
\begin{align*}
 S_{2k}(-t)&=S_{2k-2}(pt)^{-2}\,S_{2k-2}((1-4p)t)^{-1}\,S_{2k-2}(pt)^{-2}\\
 &=\bigl(S_{2k-2}(pt)^2\,S_{2k-2}((1-4p)t)\,S_{2k-2}(pt)^2\bigr)^{-1}=S_{2k}(t)^{-1},
\end{align*}
using that the list $c$ is palindromic. Hence $\log S_{2k}$ is odd in $t$,
by the uniqueness argument of \cref{prop:symmetric}.

Order: the five factors are exponentials of
$W_i=c_itS+c_i^{2k-1}t^{2k-1}E^{(k-1)}+O(t^{2k+1})$, $i=1,\ldots,5$. By
the formal BCH theorem for several factors (\cref{cor:manyexp} together
with \cref{thm:friedrichs}), $\log\prod_ie^{W_i}=\sum_iW_i+(\text{Lie
polynomials of degree at least two in the }W_i)$. A bracket in which every
argument is a multiple of $tS$ vanishes, since $[S,S]=0$; a bracket
containing at least one factor of order $t^{2k-1}$ and at least one other
factor has order at least $t^{2k}$. Therefore
\[
 \log S_{2k}(t)=tS\sum_ic_i+t^{2k-1}E^{(k-1)}\sum_ic_i^{2k-1}+O(t^{2k}).
\]
Now $\sum_ic_i=4p+(1-4p)=1$. From \eqref{eq:suzukip},
$4-4^{1/(2k-1)}=1/p$, so $1-4p=-4^{1/(2k-1)}p$ and
$(1-4p)^{2k-1}=-4p^{2k-1}$; hence $\sum_ic_i^{2k-1}=4p^{2k-1}+(1-4p)^{2k-1}=0$.
Thus $\log S_{2k}(t)=tS+O(t^{2k})$, and oddness removes the degree-$2k$
term, leaving $tS+t^{2k+1}E^{(k)}+O(t^{2k+3})$ with $E^{(k)}$ a Lie
polynomial. This completes the induction and proves the first part of
\eqref{eq:suzukiorder}; the finite word-cut description of every term
follows from \cref{cor:manyexp} applied to the $5^{k-1}\cdot3$ exponential
factors of $S_{2k}$ with their scalar weights.

For the $n$-step estimate, $S_{2k}(t/n)$ is a finite product of
exponentials whose norms sum to $\sigma_n=\lambda_k|t|a/n$, where
$\lambda_k=\prod_{j=2}^k(4p_j+|1-4p_j|)$ accounts for the (possibly
negative) substeps. For $n$ large, $\sigma_n<\tfrac12\log2$ and the
multifactor logarithm converges; $n\log S_{2k}(t/n)=tS+E_n''$ where
$E_n''$ consists of the terms of degree at least $2k+1$ of
$n\log S_{2k}(t/n)$, hence has norm at most
$n\rho^{-(2k+1)}\log\frac1{2-\sqrt2}$ with $\rho=\log2/(2\sigma_n)$, that
is $O(n\sigma_n^{2k+1})=O(n^{-2k})$. The exponential difference estimate
then gives the second part of \eqref{eq:suzukiorder}.
\end{proof}
The middle coefficient $1-4p_k$ is negative (for $k=2$,
$p_2=1/(4-4^{1/3})\approx0.4145$ and $1-4p_2\approx-0.658$). Negative
substeps are harmless for bounded operators and for unitary groups but can
be a genuine obstruction for one-sided unbounded semigroups; the
boundedness hypothesis of this section is therefore part of the result, not
an optional technicality. The proof asserts nothing about convergence as
$k\to\infty$ at fixed $t$.

\section{Exponential coordinates, coframes, and metrics to all orders}\label{sec:geometry}
The differential identity \eqref{eq:leftdexp} has an intrinsic geometric
meaning. This section develops it into complete coordinate expansions and
isolates the corrections needed in the geometric paragraph of the source
page.

\subsection{The infinitesimal formula and its full component expansion}
Let $G$ be a finite-dimensional Lie group with Lie algebra $\g$, basis
$e_1,\ldots,e_d$, and structure constants
\begin{equation}\label{eq:structure}
 [e_i,e_j]=f_{ij}{}^{k}e_k,\qquad X=X^ie_i,
\end{equation}
with repeated upper--lower indices summed. For a differentiable
$\g$-valued function $X$ on a manifold, applying \eqref{eq:leftdexp}
(intrinsically, \cref{prop:intrinsicdexp}) to every tangent vector gives
the complete infinitesimal identity
\begin{equation}\label{eq:infinitesimal}
\begin{aligned}
 e^{-X}\dd e^X:=\exp(X)^*\theta&=\phi(\ad_X)\,\dd X
 =\sum_{n\geq0}\frac{(-1)^n}{(n+1)!}\ad_X^n(\dd X)\\
 &=\dd X-\tfrac12[X,\dd X]+\tfrac16[X,[X,\dd X]]-\cdots,
\end{aligned}
\end{equation}
an equality of $\g$-valued one-forms; in a matrix group the left side is
literally $g^{-1}\dd g$ with $g=e^X$. Inserting \eqref{eq:structure} into
the successive commutators gives the component terms displayed on the
source page,
\begin{equation}\label{eq:components}
 e^{-X}\dd e^X=e_k\,\dd X^k-\tfrac12X^i\dd X^jf_{ij}{}^ke_k
 +\tfrac16X^iX^j\dd X^kf_{jk}{}^{\ell}f_{i\ell}{}^me_m-\cdots,
\end{equation}
whose indices and signs are fixed by the bracket convention rather than
guessed. Define the matrix of $\ad_X$ and the coframe matrix by
\begin{equation}\label{eq:MW}
 M^a{}_b=X^if_{ib}{}^a,\qquad
 W(X)=\phi(M)=\sum_{n\geq0}\frac{(-1)^n}{(n+1)!}M^n=\int_0^1e^{-uM}\,\dd u,
\end{equation}
so that $\ad_Xe_b=M^a{}_be_a$. Then every component at every order is
specified by
\begin{equation}\label{eq:thetaW}
\begin{gathered}
 e^{-X}\dd e^X=e_a\,W^a{}_b(X)\,\dd X^b,\\
 (M^n)^a{}_b=\sum_{i_1,\ldots,i_n}\ \sum_{c_1,\ldots,c_{n-1}}
 X^{i_1}\cdots X^{i_n}\,f_{i_1c_{n-1}}{}^{a}f_{i_2c_{n-2}}{}^{c_{n-1}}\cdots f_{i_nb}{}^{c_1},
\end{gathered}
\end{equation}
which is ordinary matrix multiplication written out; for $n=1$ it reads
$X^if_{ib}{}^a$, and for $n=2$ it reproduces the third term of
\eqref{eq:components} after renaming dummy indices. The factorial series
in \eqref{eq:MW} converges for every finite matrix $M$, and the integral
form follows from \eqref{eq:leftdexp}.

\subsection{Which matrix is the vielbein?}
\begin{proposition}[The coframe and its regularity]\label{prop:coframe}
In exponential coordinates $X^1,\ldots,X^d$ near the identity, the one-forms
$\theta^a=W^a{}_b\,\dd X^b$ form a coframe (the components of the left
Maurer--Cartan form), and the dual frame has coefficient matrix $W^{-1}$.
The matrix $M=\ad_X$ is not this coframe: in positive dimension it is
\emph{always singular}, whereas $W(0)=I$. More precisely,
\begin{equation}\label{eq:Wcriterion}
 W(X)\text{ is invertible}\quad\Longleftrightarrow\quad
 \spec(\ad_X)\cap\bigl(2\pi\ii\Z\setminus\{0\}\bigr)=\varnothing,
\end{equation}
the spectrum being computed after complexification for a real Lie algebra,
and on this set $W^{-1}=\beta(M)$ by holomorphic functional calculus (not
necessarily by a convergent Taylor series about $M=0$).
\end{proposition}
\begin{proof}
Equation \eqref{eq:thetaW} identifies $\theta^a$ with the components of
the pullback of $\theta$ under $\exp$, and since $\exp$ is a local
diffeomorphism at $0$ (\cref{thm:local}) and $W(0)=I$, the $\theta^a$ are
linearly independent near $0$: they form a coframe, and the dual frame is
given by the inverse matrix. If $X\neq0$ then $MX=[X,X]=0$ with $X\neq0$,
so $M$ has a nontrivial kernel; if $X=0$ then $M=0$. Hence $M$ is singular
in every positive-dimensional Lie algebra and cannot be a frame or coframe
matrix. Over $\C$ triangularize $M$; a power series in a triangular matrix
is triangular with diagonal entries equal to the scalar function evaluated
at the diagonal entries, so $\det W=\prod_{\lambda\in\spec M}\phi(\lambda)$
with multiplicities. The zeros of the entire function $\phi(z)=(1-e^{-z})/z$
are exactly the nonzero points of $2\pi\ii\Z$ ($1-e^{-z}=0$ iff
$z\in2\pi\ii\Z$, and at $z=0$ the value is $1$). This proves
\eqref{eq:Wcriterion}; in particular the eigenvalue $0$ forced by
$MX=0$ is harmless for $W$. Where no zero is present, the analytic
identity $\phi\beta=1$ and multiplicativity of functional calculus give
$W^{-1}=\beta(M)$.
\end{proof}
Some authors reserve ``vielbein'' for the dual vector frame rather than
the coframe; with that convention the frame coefficients are $W^{-1}$,
still not $M$. The quotient $(I-e^{-M})M^{-1}$ used on the source page is a
shorthand for the entire function $\phi(M)$; a literal inverse of $M$ does
not exist in this setting. This is a correction of a mathematical
misidentification on the page, not an alternative sign convention. The
defining differential relation of the coframe is the Maurer--Cartan
equation of \cref{prop:MC} in components,
\begin{equation}\label{eq:structureequation}
 \dd\theta^a=-\tfrac12f_{bc}{}^a\,\theta^b\wedge\theta^c,
\end{equation}
obtained by writing $\theta=e_a\theta^a$ and
$\frac12[\theta,\theta]=\frac12\theta^b\wedge\theta^c[e_b,e_c]$.

\subsection{Pullback tensors: the Jacobian and the nondegeneracy hypotheses}
Let $B$ be a symmetric bilinear form on $\g$ with matrix
$B_{ab}=B(e_a,e_b)$. Left translation defines a left-invariant symmetric
tensor $g_G$ on $G$ by $g_G(u,v)=B(\theta u,\theta v)$; it is a
pseudo-Riemannian metric exactly when $B$ is nondegenerate, and Riemannian
when $B$ is positive definite. No invariance of $B$ under the adjoint
action is needed for left invariance.

\begin{theorem}[Pullback formula with its hypotheses]\label{thm:pullback}
Let $F:N\to G$ be smooth. Near a point $q_0\in N$, after a fixed left
translation (which does not change the pullback of a left-invariant
tensor), write $F(q)=\exp X(q)$ with $X(q)=X^b(q)e_b$ in coordinates
$q^\mu$ on $N$, and set
\begin{equation}\label{eq:E}
 E^a{}_\mu(q)=W^a{}_b\bigl(X(q)\bigr)\,\partial_\mu X^b(q).
\end{equation}
Then
\begin{equation}\label{eq:pullback}
 (F^*\theta)^a=E^a{}_\mu\,\dd q^\mu,\qquad
 (F^*g_G)_{\mu\nu}=B_{ab}\,E^a{}_\mu E^b{}_\nu.
\end{equation}
In exponential coordinates themselves, $X^b=q^b$, the Jacobian is the
identity and
\begin{equation}\label{eq:metricW}
 g_{ij}=B_{ab}W^a{}_iW^b{}_j,\qquad g=W^{\mathsf T}BW.
\end{equation}
If $B$ is positive definite, $F^*g_G$ is a Riemannian metric on $N$ if and
only if $F$ is an immersion. If $B$ is indefinite, $F^*g_G$ is a
pseudo-Riemannian metric if and only if $F$ is an immersion and $B$
restricted to $\theta(\dd F(T_qN))$ is nondegenerate at every $q$.
\end{theorem}
\begin{proof}
The chain rule applied to \eqref{eq:thetaW} gives
$F^*\theta=(\exp\circ X)^*\theta=e_aW^a{}_b(X)\,\dd X^b
=e_aW^a{}_b(X)\partial_\mu X^b\,\dd q^\mu$, the first identity. By
definition of the pullback of a bilinear form,
$(F^*g_G)(\partial_\mu,\partial_\nu)=B(\theta F_*\partial_\mu,\theta F_*\partial_\nu)
=B(e_aE^a{}_\mu,e_bE^b{}_\nu)=B_{ab}E^a{}_\mu E^b{}_\nu$. If $B$ is
positive definite, then $(F^*g_G)(v,v)=B(\theta F_*v,\theta F_*v)\geq0$
with equality iff $F_*v=0$, so the pullback is positive definite iff
$\dd F$ is injective at every point. If $B$ is indefinite, the pullback is
the restriction of the nondegenerate form $B$ to the subspace
$\theta(\dd F(T_qN))$ transported by the linear map $\theta\circ\dd F$,
which is an isomorphism onto that subspace iff $\dd F$ is injective; the
transported form is nondegenerate iff the restriction is.
\end{proof}
Three distinct qualifications follow. An arbitrary smooth map does not
induce a metric: a constant map has zero pullback. For an indefinite $B$
even an immersion can have a degenerate pullback: the immersion
$q\mapsto(q,q)$ into $(\R^2,\dd x^2-\dd y^2)$ pulls the metric back to
$(1-1)\dd q^2=0$. And an arbitrary previously specified metric on $N$
cannot be asserted to be such a pullback without additional data. The
Jacobian factors in \eqref{eq:E} are indispensable unless the coordinates
on $N$ are the exponential coordinates $X^i$ themselves.

The Killing form is
\begin{equation}\label{eq:killing}
 B_K(u,v)=\tr(\ad_u\ad_v),\qquad (B_K)_{ab}=f_{ac}{}^df_{bd}{}^c,
\end{equation}
the second expression being the trace of the product of the two adjoint
matrices. It is symmetric by $\tr(AB)=\tr(BA)$, and it is invariant:
\begin{align*}
 B_K([x,u],v)+B_K(u,[x,v])
 &=\tr\bigl([\ad_x,\ad_u]\ad_v+\ad_u[\ad_x,\ad_v]\bigr)\\
 &=\tr\bigl(\ad_x\ad_u\ad_v-\ad_u\ad_x\ad_v+\ad_u\ad_x\ad_v-\ad_u\ad_v\ad_x\bigr)\\
 &=\tr[\ad_x,\ad_u\ad_v]=0,
\end{align*}
using $\ad_{[x,u]}=[\ad_x,\ad_u]$ from \eqref{eq:derivation}. But the
Killing form need not be nondegenerate: in an abelian Lie algebra every
adjoint map is zero and $B_K\equiv0$, and more generally it vanishes on
the center of any Lie algebra. (By Cartan's criterion, which is not needed
here, $B_K$ is nondegenerate exactly for semisimple Lie algebras; and for
compact semisimple algebras it is negative definite, so even then it is
$-B_K$ that is Riemannian.) Consequently calling its pullback a ``metric'' without a
nondegeneracy hypothesis is not a valid general assertion, and it is not a
universal choice of Riemannian metric; an arbitrary positive definite inner
product at the identity always gives a left-invariant Riemannian metric
instead. When $B_K$ is nondegenerate it may be used as $B$ in
\cref{thm:pullback}, with its signature respected. Formula
\eqref{eq:metricW} is valid regardless of whether it defines a metric or
only a possibly degenerate tensor.

\subsection{All-orders metric coefficients, their even resummation, and the volume density}
Without assuming invariance, substituting \eqref{eq:MW} into
\eqref{eq:metricW} gives the absolutely convergent double series
\begin{equation}\label{eq:metricdouble}
 g_{ij}=\sum_{p,q\geq0}\frac{(-1)^{p+q}}{(p+1)!\,(q+1)!}\,B_{ab}(M^p)^a{}_i(M^q)^b{}_j,
\end{equation}
valid for every $M$. If $B$ is invariant the answer is considerably
simpler.
\begin{theorem}[Complete invariant-metric expansion]\label{thm:metric}
Suppose $B([X,u],v)+B(u,[X,v])=0$ for all $u,v$ (for example $B=B_K$).
Then
\begin{align}
 g=W^{\mathsf T}BW&=B\,\phi(-M)\phi(M)
 =B\,\frac{e^M+e^{-M}-2I}{M^2}
 =B\,\frac{2(\cosh M-I)}{M^2}
 =B\Bigl(\frac{\sinh(M/2)}{M/2}\Bigr)^2\label{eq:metricclosed}\\
 &=B\sum_{n\geq0}\frac{2M^{2n}}{(2n+2)!}
 =B\Bigl(I+\frac{M^2}{12}+\frac{M^4}{360}+\frac{M^6}{20160}+\cdots\Bigr).
 \label{eq:metriceven}
\end{align}
All quotients denote entire functions with their removable values at
zero; no inverse of $M$ or $M^2$ is assumed. If moreover $B$ and $W$ are
invertible, $g^{-1}=\beta(M)B^{-1}\beta(M)^{\mathsf T}$.
\end{theorem}
\begin{proof}
In components, invariance reads $B_{ab}M^a{}_cu^cv^b=-B_{ab}u^aM^b{}_cv^c$
for all $u,v$, that is $M^{\mathsf T}B=-BM$. Induction gives
$(M^{\mathsf T})^nB=(-1)^nBM^n$: if it holds for $n$, then
$(M^{\mathsf T})^{n+1}B=M^{\mathsf T}(-1)^nBM^n=(-1)^{n+1}BM^{n+1}$. Hence
\[
 W^{\mathsf T}B=\sum_n\frac{(-1)^n}{(n+1)!}(M^{\mathsf T})^nB
 =B\sum_n\frac{M^n}{(n+1)!}=B\,\phi(-M),
\]
and $W^{\mathsf T}BW=B\,\phi(-M)\phi(M)$. The scalar identity
\[
 \phi(-z)\phi(z)=\frac{(e^z-1)(1-e^{-z})}{z^2}=\frac{e^z+e^{-z}-2}{z^2}
 =\frac{2(\cosh z-1)}{z^2}=\frac{4\sinh^2(z/2)}{z^2}
 =\sum_{n\geq0}\frac{2z^{2n}}{(2n+2)!}
\]
holds between entire functions (expand $\cosh z-1=\sum_{m\geq1}z^{2m}/(2m)!$),
and since all the matrix functions involved are power series in the single
matrix $M$, they commute and the scalar identity transfers to $M$. The
coefficients are $2/2!=1$, $2/4!=\tfrac1{12}$, $2/6!=\tfrac1{360}$,
$2/8!=\tfrac1{20160}$. The resulting matrix is symmetric because it was
obtained as $W^{\mathsf T}BW$. Finally
$g^{-1}=W^{-1}B^{-1}(W^{\mathsf T})^{-1}=\beta(M)B^{-1}\beta(M)^{\mathsf T}$
when the inverses exist.
\end{proof}
Note that all odd orders disappear only under the invariance hypothesis;
no such cancellation occurs in \eqref{eq:metricdouble}.

When $B$ is nondegenerate and the $X^i$ are exponential coordinates, the
metric matrix is $W^{\mathsf T}BW$, so its volume density is
\begin{equation}\label{eq:volume}
 \sqrt{\abs{\det g}}=\abs{\det W}\sqrt{\abs{\det B}},\qquad
 \det W=\prod_{\lambda\in\spec M}\phi(\lambda),
\end{equation}
with the factor $1$ at $\lambda=0$ and multiplicities included, by the
triangularization used in \cref{prop:coframe}. If the spectral radius of
$M$ is less than $2\pi$, then
\begin{equation}\label{eq:logvolume}
 \det W=\exp\Bigl(-\tfrac12\tr M+\sum_{m\geq1}\frac{\bplus{2m}}{2m}\tr\bigl(M^{2m}\bigr)\Bigr)
 =\exp\Bigl(-\tfrac12\tr M+\frac{\tr M^2}{24}-\frac{\tr M^4}{2880}+\cdots\Bigr).
\end{equation}
\begin{proof}[Proof of \eqref{eq:logvolume}]
For $|z|<2\pi$, $\phi(z)\neq0$ and
\[
 \frac{\dd}{\dd z}\log\phi(z)=\frac{e^{-z}}{1-e^{-z}}-\frac1z
 =\frac1{e^z-1}-\frac1z=\frac{\gamma(z)-1}{z}
 =\sum_{m\geq1}\bminus{m}z^{m-1}=-\frac12+\sum_{m\geq1}\bplus{2m}z^{2m-1},
\]
using $\bminus1=-\tfrac12$ and $\bminus{2m}=\bplus{2m}$, $\bminus{2m+1}=0$
for $m\geq1$. Integrating with $\log\phi(0)=0$ gives
$\log\phi(z)=-\tfrac z2+\sum_{m\geq1}\tfrac{\bplus{2m}}{2m}z^{2m}$, a
series with radius $2\pi$ (the nearest zeros of $\phi$). Hence
$\phi(\lambda_j)=\exp(-\tfrac{\lambda_j}2+\sum_m\tfrac{\bplus{2m}}{2m}\lambda_j^{2m})$
for every eigenvalue $\lambda_j$ with $|\lambda_j|<2\pi$; multiplying over
$j$ and using $\sum_j\lambda_j^{k}=\tr M^k$ (absolute convergence allows
the interchange of the sums over $j$ and $m$) gives \eqref{eq:logvolume}.
The first coefficients are $\bplus2/2=\tfrac1{24}$ and
$\bplus4/4=-\tfrac1{2880}$.
\end{proof}
The determinant formula \eqref{eq:volume} itself has no radius restriction.

\subsection{A concrete three-dimensional resummation}
Identify the rotation Lie algebra $\mathfrak{so}(3)$ with $\R^3$ and the
bracket with the cross product, $[x,v]=x\times v$. Then $Mv=x\times v$ and,
with $r=|x|$, the vector triple-product identity gives
$M^2v=x\times(x\times v)=x(x\cdot v)-r^2v$ and hence
$M^3v=x\times\bigl(x(x\cdot v)-r^2v\bigr)=-r^2\,x\times v$, that is
$M^3=-r^2M$. Consequently $M^{2k+1}=(-r^2)^kM$ and $M^{2k+2}=(-r^2)^kM^2$
for $k\geq0$, and summing the exponential series,
\[
 e^{-sM}=I+\sum_{k\geq0}\frac{(-s)^{2k+1}(-r^2)^k}{(2k+1)!}M
 +\sum_{k\geq0}\frac{(-s)^{2k+2}(-r^2)^k}{(2k+2)!}M^2
 =I-\frac{\sin(sr)}{r}M+\frac{1-\cos(sr)}{r^2}M^2,
\]
which is Rodrigues' formula. Integrating over $s\in[0,1]$ by
\eqref{eq:MW} gives the complete coframe
\begin{equation}\label{eq:so3W}
 W=I-\frac{1-\cos r}{r^2}M+\frac{r-\sin r}{r^3}M^2,
\end{equation}
with the Taylor values $\tfrac12$ and $\tfrac16$ of the two coefficients at
$r=0$. The Euclidean dot product is invariant, since
$(x\times u)\cdot v+u\cdot(x\times v)=0$. For $r>0$ let
$P_\parallel=xx^{\mathsf T}/r^2$ and $P_\perp=I-P_\parallel$; then $M=0$ on
the range of $P_\parallel$ and $M^2=-r^2I$ on the range of $P_\perp$ (the
plane orthogonal to $x$, where $x\cdot v=0$). Formula
\eqref{eq:metricclosed} with $B=I$ therefore gives
\begin{equation}\label{eq:so3metric}
 W^{\mathsf T}W=P_\parallel+\frac{2(1-\cos r)}{r^2}P_\perp
 =P_\parallel+\Bigl(\frac{\sin(r/2)}{r/2}\Bigr)^2P_\perp,
 \qquad
 \det W=\Bigl(\frac{\sin(r/2)}{r/2}\Bigr)^2,
\end{equation}
the determinant following either from the eigenvalues $0,\ii r,-\ii r$ of
$M$ and $\phi(\ii r)\phi(-\ii r)=|1-e^{-\ii r}|^2/r^2=(2-2\cos r)/r^2$, or
directly from \eqref{eq:so3W}. Thus the coordinate coframe of the rotation
group in exponential coordinates is singular exactly at the radii
$r=2\pi,4\pi,\ldots$, precisely as \eqref{eq:Wcriterion} predicts; at
$r=0$ the whole expression has the unambiguous limit $I$.

\section{Quantum-mechanical identities with explicit operator domains}\label{sec:quantum}
The formal central-commutator calculation of \cref{thm:central} predicts
the exponential formulas on the source page. For operators on a Hilbert
space, however, unboundedness and domains cannot be suppressed: we first
prove the identities directly in the Schr\"odinger and Bargmann--Fock
realizations, and then exhibit why a commutator identity on a dense domain
alone is insufficient. This avoids the Stone--von Neumann classification
theorem, which is mentioned but not stated on the page; background on these
realizations may be found in \cite{hall,glauber,wilcox}. Inner products are
linear in the second argument, and $\ket n,\bra n$ denote basis vectors and
functionals.

\subsection{Position, momentum, and the Weyl relations}
On $L^2(\R,\dd x)$ fix $\hbar>0$ and define, on the Schwartz space
$\cS(\R)$,
\begin{equation}\label{eq:QP}
 (Qf)(x)=xf(x),\qquad (Pf)(x)=-\ii\hbar f'(x).
\end{equation}
Both map $\cS$ into itself, and direct differentiation gives
$(QPf)(x)-(PQf)(x)=-\ii\hbar xf'(x)+\ii\hbar(xf)'(x)=\ii\hbar f(x)$, that is,
\begin{equation}\label{eq:CCR}
 [Q,P]f=\ii\hbar f\qquad(f\in\cS(\R)).
\end{equation}
Both operators are symmetric on $\cS$ (integration by parts for $P$, with
boundary terms vanishing by decay). Their self-adjoint realizations are the
maximal multiplication operator $\overline Q$ on
$\{f:xf\in L^2\}$ and $\overline P=\mathcal F^{-1}(\hbar k)\mathcal F$, the
Fourier conjugate of multiplication by $\hbar k$, on $\{f:k\hat f\in L^2\}$;
these agree with \eqref{eq:QP} on $\cS$, since the Fourier transform of
$-\ii\hbar f'$ is $\hbar k\hat f$ for $f\in\cS$. Schwartz space is a core
for each. For multiplication by a real function $m$ (here $m(x)=x$ or, after
Fourier transform, $m(k)=\hbar k$, and $\mathcal F$ maps $\cS$ onto $\cS$)
this follows from the basic criterion: a symmetric operator $A$ is
essentially self-adjoint iff $(A\pm\ii)\Dom A$ are dense; if $g\perp(m\pm\ii)f$
for all $f\in\cS$, then $\int\overline{(m\mp\ii)g}\,f=0$ for all $f\in\cS$
(as $m$ is real), so the locally integrable function $(m\pm\ii)\bar g$
vanishes almost everywhere, and since $m\pm\ii\neq0$ we get $g=0$.

The unitary groups are concrete:
\begin{equation}\label{eq:UV}
 (U(a)f)(x)=(e^{\ii aQ}f)(x)=e^{\ii ax}f(x),\qquad
 (V(b)f)(x)=(e^{\ii bP}f)(x)=f(x+b\hbar),\qquad a,b\in\R.
\end{equation}
The first is the spectral calculus of the multiplication operator. For the
second, $e^{\ii b\overline P}=\mathcal F^{-1}e^{\ii b\hbar k}\mathcal F$, and
multiplication of $\hat f$ by $e^{\ii b\hbar k}$ is the Fourier transform
of the translate $f(\cdot+b\hbar)$. Both families are unitary by a change
of variables and strongly continuous: for $f\in C_c^\infty$,
$\norm{V(b)f-f}\to0$ by uniform continuity, and this extends to $L^2$ by
density and the uniform bound $\norm{V(b)}=1$; similarly for $U$.

\begin{theorem}[Weyl relations and exact symmetric splitting]\label{thm:weyl}
For real $a,b$, the operator $aQ+bP$ on $\cS(\R)$ is essentially
self-adjoint, and the unitary group of its closure acts by
\begin{equation}\label{eq:sumaction}
 \bigl(e^{\ii(aQ+bP)}f\bigr)(x)=e^{\ii a(x+b\hbar/2)}f(x+b\hbar)=:(\mathcal W(a,b)f)(x).
\end{equation}
Consequently, as identities between bounded operators on all of $L^2(\R)$,
\begin{align}
 e^{\ii aQ}e^{\ii bP}&=e^{-\ii ab\hbar/2}\,e^{\ii(aQ+bP)}
 =\exp\bigl(\ii(aQ+bP-\tfrac12ab\hbar I)\bigr),\label{eq:weylbch}\\
 e^{\ii(aQ+bP)}&=e^{\ii aQ/2}e^{\ii bP}e^{\ii aQ/2},\label{eq:weylsym}\\
 e^{\ii aQ}e^{\ii bP}&=e^{-\ii ab\hbar}\,e^{\ii bP}e^{\ii aQ},\label{eq:weylcomm}\\
 \mathcal W(a,b)\mathcal W(c,d)&=e^{\ii\hbar(bc-ad)/2}\,\mathcal W(a+c,b+d).\label{eq:weylgroup}
\end{align}
\end{theorem}
\begin{proof}
Let $b\neq0$ and let $T$ be multiplication by $\tau(x)=e^{\ii ax^2/(2b\hbar)}$,
a unitary operator preserving $\cS$. For $f\in\cS$, since
$\tau'=(\ii ax/(b\hbar))\tau$,
\[
 T^{-1}(bP)Tf=\tau^{-1}(-\ii\hbar b)(\tau f)'
 =(-\ii\hbar b)\Bigl(\frac{\ii ax}{b\hbar}f+f'\Bigr)=axf-\ii\hbar bf'=(aQ+bP)f.
\]
Thus $aQ+bP$ on $\cS$ is unitarily equivalent to $bP$ on $\cS$, which is
essentially self-adjoint with closure $b\overline P$; hence $aQ+bP$ is
essentially self-adjoint with closure $T^{-1}b\overline PT$, and its unitary
group is $T^{-1}e^{\ii tbP}T=T^{-1}V(tb)T$. Evaluating,
\begin{align*}
 (T^{-1}V(b)Tf)(x)&=e^{-\ii ax^2/(2b\hbar)}e^{\ii a(x+b\hbar)^2/(2b\hbar)}f(x+b\hbar)\\
 &=e^{\ii a(2xb\hbar+b^2\hbar^2)/(2b\hbar)}f(x+b\hbar)=e^{\ii a(x+b\hbar/2)}f(x+b\hbar),
\end{align*}
which is \eqref{eq:sumaction}; for $b=0$ the statement is the multiplication
case $U(a)$, also of the form \eqref{eq:sumaction}. (One can also verify
directly that $\mathcal W(ta,tb)$ is a strongly continuous one-parameter
unitary group whose derivative at $t=0$ on $\cS$ is $\ii(aQ+bP)$.)

Now compute with \eqref{eq:UV}: $(U(a)V(b)f)(x)=e^{\ii ax}f(x+b\hbar)
=e^{-\ii ab\hbar/2}(\mathcal W(a,b)f)(x)$, which is \eqref{eq:weylbch}; the
second form holds because the scalar $-\tfrac12ab\hbar I$ commutes with
everything. Next, $(U(a/2)V(b)U(a/2)f)(x)=e^{\ii ax/2}e^{\ii a(x+b\hbar)/2}f(x+b\hbar)
=e^{\ii a(x+b\hbar/2)}f(x+b\hbar)$, which is \eqref{eq:weylsym}. Further,
$(V(b)U(a)f)(x)=e^{\ii a(x+b\hbar)}f(x+b\hbar)=e^{\ii ab\hbar}(U(a)V(b)f)(x)$,
which is \eqref{eq:weylcomm}. Finally,
\[
 (\mathcal W(a,b)\mathcal W(c,d)f)(x)
 =e^{\ii a(x+b\hbar/2)}e^{\ii c(x+b\hbar+d\hbar/2)}f(x+(b+d)\hbar),
\]
and the phase of $\mathcal W(a+c,b+d)f$ is $e^{\ii(a+c)(x+(b+d)\hbar/2)}$;
the difference of exponents is
$\ii\hbar\bigl(\tfrac{ab}2+cb+\tfrac{cd}2-\tfrac{(a+c)(b+d)}2\bigr)=\ii\hbar\tfrac{bc-ad}2$,
which proves \eqref{eq:weylgroup}. All four identities involve only
bounded unitary operators and hold on all of $L^2(\R)$.
\end{proof}
The formal central-commutator calculation agrees: $[\ii aQ,\ii bP]=-ab\cdot\ii\hbar I
=-\ii ab\hbar I$ is central, so \eqref{eq:centralbch} predicts exactly the
correction $-\tfrac12\ii ab\hbar I$ in \eqref{eq:weylbch}, and the symmetric
splitting of \cref{thm:central} becomes exact. The direct proof establishes
that the formal calculation is legitimate for the chosen self-adjoint
realizations; \eqref{eq:weylbch} has not been justified by substituting
unbounded operators into a norm-convergent series. Induction in
\eqref{eq:weylgroup} gives the complete finite-factor formula
\begin{equation}\label{eq:weylmulti}
 \prod_{j=1}^m\mathcal W(a_j,b_j)
 =\exp\Bigl(\frac{\ii\hbar}{2}\sum_{j<k}(b_ja_k-a_jb_k)\Bigr)\,
 \mathcal W\Bigl(\sum_ja_j,\sum_jb_j\Bigr),
\end{equation}
which is \eqref{eq:centralmany} in its unitary Heisenberg representation,
with no omitted higher-order terms.

\subsection{Why the infinitesimal commutator is not enough}
\begin{proposition}\label{prop:ccrcounter}
A relation $[Q,P]=\ii I$ on a dense common invariant domain, even for
self-adjoint $Q$ and $P$, does not by itself imply the Weyl relations.
\end{proposition}
\begin{proof}
On $L^2(0,1)$ let $Q$ be multiplication by $x$, a bounded self-adjoint
operator, and let $P=-\ii\,\dd/\dd x$ with the periodic domain
$\{f\in H^1(0,1):f(0)=f(1)\}$. Under the Fourier series isomorphism
$L^2(0,1)\cong\ell^2(\Z)$ given by the orthonormal basis $e^{2\pi\ii nx}$,
$P$ becomes multiplication by the real sequence $2\pi n$ on its maximal
domain; so $P$ is self-adjoint, with $e^{\ii P}$ acting as multiplication by
$e^{2\pi\ii n}=1$, that is, $e^{\ii P}=I$. The dense subspace
$C_c^\infty(0,1)$ lies in the domains of $Q$ and $P$, is invariant under
both, and satisfies $[Q,P]=\ii I$ there by differentiation. If the Weyl
relation \eqref{eq:weylcomm} (with $\hbar=1$) held for all real $a$ and
$b=1$, it would give $e^{\ii aQ}=e^{-\ii a}e^{\ii aQ}$, hence $e^{-\ii a}=1$,
which is false unless $a\in2\pi\Z$.
\end{proof}
The subspace $C_c^\infty(0,1)$ is not a core for the periodic derivative,
and the periodic domain is not preserved by multiplication by $e^{\ii ax}$:
that is exactly where a naive exponentiation argument breaks down. The
failure is a domain and integrability issue, not an error in the formal
central-commutator algebra.

There cannot even be bounded operators $A,B$ on a Banach space with
$[A,B]=I$. Indeed, induction gives $[A,B^n]=nB^{n-1}$: if it holds for
$n$, then $[A,B^{n+1}]=[A,B^n]B+B^n[A,B]=nB^{n-1}B+B^n=(n+1)B^n$. No power
of $B$ vanishes: if $B^m=0$ with $m\geq1$ minimal, then
$0=[A,B^m]=mB^{m-1}$ forces $B^{m-1}=0$, contradicting minimality unless
$m=1$, and $B=0$ contradicts $[A,B]=I$. Hence
$n\norm{B^{n-1}}\leq2\norm A\norm{B^n}\leq2\norm A\norm B\norm{B^{n-1}}$
with $\norm{B^{n-1}}\neq0$, so $n\leq2\norm A\norm B$ for every $n$, which
is impossible. (For finite matrices, taking traces gives the contradiction
$0=\tr[A,B]=\tr I=d$ immediately.) Some unboundedness is therefore
intrinsic to the canonical relation, and the bounded-operator convergence
theorems of \cref{sec:convergence} cannot be applied directly to $Q,P$.

\subsection{Creation and annihilation operators in Bargmann--Fock space}
Let $\cF$ be the space of entire functions $F$ on $\C$ with
\begin{equation}\label{eq:focknorm}
 \norm F_{\cF}^2=\frac1\pi\int_{\C}|F(z)|^2e^{-|z|^2}\,\dd^2z<\infty.
\end{equation}
Polar integration gives $\langle z^m,z^n\rangle=\frac1\pi\int_0^\infty\!\!\int_0^{2\pi}
\rho^{m+n}e^{\ii(n-m)\varphi}e^{-\rho^2}\rho\,\dd\varphi\,\dd\rho=\delta_{mn}\,2\int_0^\infty\rho^{2n+1}e^{-\rho^2}\dd\rho=\delta_{mn}n!$,
so
\begin{equation}\label{eq:fockbasis}
 \ket n=e_n(z)=\frac{z^n}{\sqrt{n!}}\qquad(n\geq0)
\end{equation}
is an orthonormal family. It is complete: expanding an entire $F$ in its
Taylor series $\sum c_nz^n$, which converges uniformly on each circle,
integrating first over the angle (where the cross terms vanish) and then
over the radius by monotone convergence gives
$\norm F^2=\sum_nn!|c_n|^2$, so $F=\sum_nc_n\sqrt{n!}\,\ket n$ in norm and
polynomials are dense. Conversely $\sum n!|c_n|^2<\infty$ implies that
$\sum c_nz^n$ is entire, by Cauchy--Schwarz with
$\sum|z|^{2n}/n!=e^{|z|^2}$; the same estimate gives
\begin{equation}\label{eq:pointeval}
 |F(z)|\leq\Bigl(\sum_nn!|c_n|^2\Bigr)^{1/2}\Bigl(\sum_n\frac{|z|^{2n}}{n!}\Bigr)^{1/2}
 =\norm F_{\cF}\,e^{|z|^2/2},
\end{equation}
so norm convergence in $\cF$ implies locally uniform pointwise convergence.
On the dense polynomial domain $\mathcal P$ define
\begin{equation}\label{eq:creation}
\begin{gathered}
 aF=F',\qquad a^\dagger F=zF,\qquad
 [a,a^\dagger]=I,\qquad[a^\dagger,a]=-I,\\
 a\ket n=\sqrt n\,\ket{n-1},\qquad a^\dagger\ket n=\sqrt{n+1}\,\ket{n+1}.
\end{gathered}
\end{equation}
The ladder formulas are differentiation and multiplication of $z^n/\sqrt{n!}$;
the commutator is $(zF)'-zF'=F$. The ladder formulas show
$\langle a^\dagger F,G\rangle=\langle F,aG\rangle$ for $F,G\in\mathcal P$,
which justifies the notation $\dagger$ (the maximal closed realizations are
Hilbert-space adjoints of each other).

\begin{theorem}[Displacement operators, generator, and normal ordering]\label{thm:displacement}
For $v\in\C$ the formula
\begin{equation}\label{eq:Daction}
 (D(v)F)(z)=e^{-|v|^2/2+vz}\,F(z-\bar v)
\end{equation}
defines a unitary operator on $\cF$. For fixed $v$, $t\mapsto D(tv)$
$(t\in\R)$ is a strongly continuous one-parameter unitary group whose
skew-adjoint generator is the closure of the operator
$T_v=va^\dagger-\bar va$ defined on $\mathcal P$; in this precise sense
\begin{equation}\label{eq:normalorder}
 D(v)=e^{va^\dagger-\bar va}=e^{va^\dagger}e^{-\bar va}e^{-|v|^2/2},
\end{equation}
where the right-hand product is defined on $\mathcal P$ (the middle factor
by a finite sum, the left factor by a norm-convergent series) and extends by
continuity to the unitary operator \eqref{eq:Daction}.
\end{theorem}
\begin{proof}
\emph{Unitarity.} Substitute $z=w+\bar v$ in \eqref{eq:focknorm}: using
$\operatorname{Re}(v\bar v)=|v|^2$ and
$|w+\bar v|^2=|w|^2+2\operatorname{Re}(vw)+|v|^2$,
\[
 \bigl|e^{-|v|^2/2+vz}\bigr|^2e^{-|z|^2}
 =e^{-|v|^2+2\operatorname{Re}(vw)+2|v|^2-|w|^2-2\operatorname{Re}(vw)-|v|^2}
 =e^{-|w|^2}.
\]
Hence
$\norm{D(v)F}=\norm F$, and $D(v)F$ is entire, so $D(v)$ is an isometry of
$\cF$. Direct substitution shows $D(-v)D(v)=I=D(v)D(-v)$, so $D(v)$ is
unitary.

\emph{Group law and continuity.} For real $s,t$, applying
\eqref{eq:Daction} twice,
$(D(tv)D(sv)F)(z)=e^{-t^2|v|^2/2+tvz}e^{-s^2|v|^2/2+sv(z-t\bar v)}F(z-(t+s)\bar v)
=e^{-(t+s)^2|v|^2/2+(t+s)vz}F(z-(t+s)\bar v)$, since
$-t^2|v|^2/2-s^2|v|^2/2-ts|v|^2=-(t+s)^2|v|^2/2$; thus $D(tv)D(sv)=D((t+s)v)$.
For a polynomial $F$ of degree $m$ and $|t|\leq1$, the integrand of
$\norm{D(tv)F-F}^2$ is bounded by $C(1+|z|)^{2m}e^{2|v||z|-|z|^2}$, which is
integrable, and converges pointwise to $0$ as $t\to0$; dominated convergence
gives $\norm{D(tv)F-F}\to0$. For general $F$, approximate by polynomials and
use $\norm{D(tv)}=1$. By Stone's theorem the group has a skew-adjoint
generator $G$.

\emph{The exponential series on polynomials.} On the span of
$\ket0,\ldots,\ket j$, the ladder formulas give $\norm a\leq\sqrt j$ and
$\norm{a^\dagger}\leq\sqrt{j+1}$, and $T_v$ maps this span into the span of
$\ket0,\ldots,\ket{j+1}$ with norm at most $|v|(\sqrt j+\sqrt{j+1})\leq2|v|\sqrt{j+1}$.
Iterating from a polynomial $F$ of degree at most $m$,
\begin{equation}\label{eq:analyticvectors}
 \norm{T_v^kF}\leq(2|v|)^k\sqrt{\frac{(m+k)!}{m!}}\,\norm F,
\end{equation}
so $\sum_k|t|^k\norm{T_v^kF}/k!$ converges for every $t$ (the ratio of
consecutive terms is $2|v||t|\sqrt{m+k+1}/(k+1)\to0$). Define
$E(t)F=\sum_kt^kT_v^kF/k!$, a norm-convergent series. We claim
$E(t)F=D(tv)F$. As a differential operator on entire functions,
$T_v=vz-\bar v\,\partial_z$. For a polynomial $F$ and fixed $z$, direct
differentiation of \eqref{eq:Daction} gives
\begin{align*}
 \frac{\dd}{\dd t}\bigl(D(tv)F\bigr)(z)
 &=e^{-t^2|v|^2/2+tvz}\bigl[(vz-t|v|^2)F(z-t\bar v)-\bar vF'(z-t\bar v)\bigr]\\
 &=e^{-t^2|v|^2/2+tvz}\bigl[v(z-t\bar v)F(z-t\bar v)-\bar vF'(z-t\bar v)\bigr]
 =\bigl(D(tv)(T_vF)\bigr)(z).
\end{align*}
Since $T_vF$ is again a polynomial, induction gives
$\frac{\dd^k}{\dd t^k}(D(tv)F)(z)=(D(tv)T_v^kF)(z)$, and at $t=0$ this is
$(T_v^kF)(z)$. The function $t\mapsto(D(tv)F)(z)$ is entire in $t$ (a
product of entire functions), so it equals its Taylor series:
$(D(tv)F)(z)=\sum_kt^k(T_v^kF)(z)/k!$. By \eqref{eq:pointeval}, the
norm-convergent series $E(t)F$ converges pointwise to the same sum. Two
elements of $\cF$ that agree pointwise are equal; hence $E(t)F=D(tv)F$.

\emph{Identification of the generator.} From the norm-convergent power
series, $t\mapsto D(tv)F=E(t)F$ is norm-differentiable at $t=0$ with
derivative $T_vF$; thus $\mathcal P\subseteq\Dom G$ and $G|_{\mathcal P}=T_v$,
that is $T_v\subseteq G$. As $G$ is closed, $T_v$ is closable with
$\overline{T_v}\subseteq G$. We show that $(I-T_v)\mathcal P$ is dense. If
$h\perp(I-T_v)\mathcal P$, then $\langle h,T_vF\rangle=\langle h,F\rangle$
for every polynomial $F$, and applying this to $T_v^{k-1}F$ gives
$\langle h,T_v^kF\rangle=\langle h,F\rangle$ for all $k$. Therefore
\[
 \langle h,D(tv)F\rangle=\langle h,E(t)F\rangle
 =\sum_k\frac{t^k}{k!}\langle h,T_v^kF\rangle=e^{t}\langle h,F\rangle
 \qquad(t\in\R),
\]
and the left side is bounded by $\norm h\norm F$ for all $t$, so letting
$t\to+\infty$ forces $\langle h,F\rangle=0$; density of $\mathcal P$ gives
$h=0$. The same argument with $e^{-t}$ and $t\to-\infty$ shows that
$(I+T_v)\mathcal P$ is dense. Now $\overline{T_v}$ is skew-symmetric, so
$\norm{(I\mp\overline{T_v})x}^2=\norm x^2+\norm{\overline{T_v}x}^2$; hence
$I-\overline{T_v}$ is bounded below and, being closed, has closed range,
which contains the dense set $(I-T_v)\mathcal P$ and is therefore all of
$\cF$. Given $y\in\Dom G$, choose $x\in\Dom\overline{T_v}$ with
$(I-\overline{T_v})x=(I-G)y$; since $\overline{T_v}\subseteq G$ this reads
$(I-G)x=(I-G)y$, and $I-G$ is injective because $G$ is skew-adjoint
($\norm{(I-G)u}^2=\norm u^2+\norm{Gu}^2$). Thus $y=x\in\Dom\overline{T_v}$,
proving $\overline{T_v}=G$: the closure of $T_v$ is exactly the generator,
with $\mathcal P$ as a core, and $D(v)=e^{\overline{T_v}}$.

\emph{Normal ordering.} For a polynomial $F$, the series
$e^{-\bar va}F=\sum_k(-\bar v)^kF^{(k)}/k!$ is a finite sum equal to
$F(z-\bar v)$ by Taylor's formula for polynomials. The series
$e^{va^\dagger}G=\sum_kv^kz^kG/k!$ converges in $\cF$ for every polynomial
$G$: for $G=\sum_jg_jz^j$, $\norm{z^kz^j}=\sqrt{(k+j)!}$ and
$\sum_k|v|^k\sqrt{(k+j)!}/k!<\infty$, and its sum is the entire function
$e^{vz}G(z)$, again by \eqref{eq:pointeval}. Hence
$e^{va^\dagger}e^{-\bar va}e^{-|v|^2/2}F=e^{-|v|^2/2}e^{vz}F(z-\bar v)=D(v)F$
on $\mathcal P$, and the bounded extension is $D(v)$.
\end{proof}
Although the ordered product in \eqref{eq:normalorder} has a unitary
extension, its two nonscalar factors are separately unbounded; the identity
does not assert that each is everywhere defined.

\begin{theorem}[Composition law and all matrix elements]\label{thm:Dcomposition}
For $u,v\in\C$,
\begin{equation}\label{eq:Dproduct}
 D(v)D(u)=\exp\Bigl(\frac{v\bar u-u\bar v}{2}\Bigr)D(v+u),
\end{equation}
and for any finite list $v_1,\ldots,v_m$,
\begin{equation}\label{eq:Dmany}
 \prod_{j=1}^mD(v_j)=\exp\Bigl(\frac12\sum_{j<k}(v_j\bar v_k-v_k\bar v_j)\Bigr)
 D\Bigl(\sum_jv_j\Bigr).
\end{equation}
The full normally ordered expansion, its number-state matrix elements, and
the coherent states are
\begin{align}
 D(v)&=e^{-|v|^2/2}\sum_{r,s\geq0}\frac{v^r(-\bar v)^s}{r!\,s!}(a^\dagger)^ra^s
 \qquad\text{on }\mathcal P,\label{eq:Dall}\\
 \bra mD(v)\ket n&=e^{-|v|^2/2}\sqrt{m!\,n!}
 \sum_{k=0}^{\min(m,n)}\frac{v^{m-k}(-\bar v)^{n-k}}{k!\,(m-k)!\,(n-k)!},\label{eq:Dmatrix}\\
 \ket v:=D(v)\ket0&=e^{-|v|^2/2}\sum_{n\geq0}\frac{v^n}{\sqrt{n!}}\ket n,
 \qquad\braket uv=\exp\Bigl(-\frac{|u|^2+|v|^2}2+\bar uv\Bigr),\qquad\norm{\ket v}=1.
 \label{eq:coherent}
\end{align}
For $m\geq n$, equivalently,
\begin{equation}\label{eq:laguerre}
 \bra mD(v)\ket n=e^{-|v|^2/2}\sqrt{\frac{n!}{m!}}\,v^{m-n}L_n^{(m-n)}(|v|^2),
 \qquad L_n^{(\alpha)}(x):=\sum_{j=0}^n\binom{n+\alpha}{n-j}\frac{(-x)^j}{j!},
\end{equation}
and the case $n\geq m$ follows from $D(v)^*=D(-v)$.
\end{theorem}
\begin{proof}
Applying \eqref{eq:Daction} twice,
\begin{align*}
 (D(v)D(u)F)(z)&=e^{-|v|^2/2+vz}e^{-|u|^2/2+u(z-\bar v)}F(z-\bar v-\bar u),\\
 (D(v+u)F)(z)&=e^{-|v+u|^2/2+(v+u)z}F(z-\bar v-\bar u).
\end{align*}
The difference of the exponents is
$-\tfrac{|v|^2}2-\tfrac{|u|^2}2-u\bar v+\tfrac{|v|^2+|u|^2+v\bar u+u\bar v}2
=\tfrac{v\bar u-u\bar v}2$, which is purely imaginary (it equals
$\ii\operatorname{Im}(v\bar u)$), proving \eqref{eq:Dproduct}; the factor is
a phase. Induction on $m$, exactly as in \cref{cor:centralmany}, gives
\eqref{eq:Dmany}: the new phase from $D(v_1+\cdots+v_{m-1})D(v_m)$ is
$\frac12\sum_{j<m}(v_j\bar v_m-v_m\bar v_j)$.

Expanding the two exponentials in \eqref{eq:normalorder} on a polynomial
(the $a$-series is finite and the $a^\dagger$-series converges in norm)
gives \eqref{eq:Dall}. Acting on $\ket n$: the term $a^s$ with $s=n-k$
lowers $\ket n$ to $\sqrt{n!/k!}\,\ket k$, and $(a^\dagger)^r$ with $r=m-k$
raises $\ket k$ to $\sqrt{m!/k!}\,\ket m$; the product of the ladder factors
is $\sqrt{m!n!}/k!$, and summing over the admissible intermediate levels
$0\leq k\leq\min(m,n)$ gives \eqref{eq:Dmatrix}. Setting $n=0$ (only $k=0$)
gives $\bra mD(v)\ket0=e^{-|v|^2/2}v^m/\sqrt{m!}$, which is the first
formula in \eqref{eq:coherent}; equivalently
$D(v)\ket0=e^{-|v|^2/2}e^{vz}$ by \eqref{eq:Daction}. Its norm squared is
$e^{-|v|^2}\sum|v|^{2n}/n!=1$, and
$\braket uv=e^{-(|u|^2+|v|^2)/2}\sum_n\bar u^nv^n/n!$ gives the overlap.
For \eqref{eq:laguerre} with $m\geq n$, substitute $j=n-k$ in
\eqref{eq:Dmatrix}: $v^{m-k}(-\bar v)^{n-k}=v^{m-n}(-|v|^2)^j$ and
\begin{align*}
 \sqrt{\frac{n!}{m!}}\,L_n^{(m-n)}(|v|^2)
 &=\sqrt{\frac{n!}{m!}}\sum_{j=0}^{n}\frac{m!}{(n-j)!\,(m-n+j)!}\frac{(-|v|^2)^j}{j!}\\
 &=\sqrt{m!\,n!}\sum_{j=0}^{n}\frac{(-|v|^2)^j}{(n-j)!\,(m-n+j)!\,j!},
\end{align*}
which matches \eqref{eq:Dmatrix} term by term with $k=n-j$. Since
$D(v)^*=D(v)^{-1}=D(-v)$, $\bra mD(v)\ket n=\overline{\bra nD(-v)\ket m}$
handles $n\geq m$.
\end{proof}
Algebraically, $[T_v,T_u]=(v\bar u-u\bar v)I$ on $\mathcal P$, which is
central and purely imaginary, and all triple commutators vanish: this is
the class-two nilpotent Heisenberg structure of \eqref{eq:heisenberglaw}
behind \eqref{eq:Dproduct}. The representation-level proof above
establishes the identity globally, without an unjustified interchange of
unbounded-operator series. Finitely many independent modes are handled by
tensoring these representations: all cross-mode commutators vanish, every
scalar product $v\bar u$ in the phase formulas is replaced by
$\sum_\ell v_\ell\bar u_\ell$, and the vacuum and matrix-element formulas
factor by mode; no further BCH terms appear. These formulas provide all
terms of the oscillator applications on the source page, with no asymptotic
truncation remaining.

\section{Exact computation, additional terms, and reproducibility}\label{sec:computation}
The all-order assertions of this article are proved algebraically or
analytically; a finite computation is not their substitute. Computation is
nevertheless useful for certifying transcription, signs, bracket
conventions, and the additional explicit terms, and for detecting errors in
the finite displays. This section describes the exact computations that
accompany the article.

\subsection{Representation and independently organized calculations}
An associative polynomial is stored as a dictionary from words in $\{X,Y\}$
to exact rational numbers (Python's \texttt{fractions.Fraction}); the empty
word represents $1$. Multiplication concatenates words and discards
degrees beyond a requested cutoff. Exponential and logarithm are evaluated
by their defining series, which become finite under this truncation.
Because no relation is imposed on the two generators, equality of the
resulting dictionaries is equality in the truncated free associative
algebra, not a test on selected numerical matrices in which independent
commutators might collapse.

The three source treatments each supplied a standard-library Python
verifier, and all three are retained unchanged in the repository
(\texttt{paper-1/code}, \texttt{paper-2/code}, \texttt{paper-3/code}). Between
them they perform the following independently organized calculations, each
compared coefficient by coefficient with the reference expansion
$\log(e^Xe^Y)$:
\begin{itemize}
\item the finite word-cut formula \eqref{eq:wordcut} evaluated by a dynamic
program over all contiguous cuts of every word, including words with zero
coefficient (through degree $10$; $2046$ words);
\item the Poincar\'e operator-block expansion \eqref{eq:operatorblock},
assembled by repeated application of $e^{\ad_X}e^{t\ad_Y}-I$ with exact
integration of powers of $t$ (through degree $10$);
\item the right Dynkin projection \eqref{eq:dynkinword}, verified to fix
every homogeneous component (through degree $12$);
\item the total-degree Bernoulli recursion \eqref{eq:homrec} (through
degree $10$), the single-$Y$ Bernoulli series \eqref{eq:linearY} (through
degree $12$), and the quadratic kernel \eqref{eq:H2kernel} (through degree
$10$);
\item the primitive coproduct condition, by explicit computation of the
reduced unshuffle coproduct (through degree $10$);
\item the exchange parity \eqref{eq:parity} and word-reversal parity
\eqref{eq:reversal} (through degree $12$), and three-letter associativity
\eqref{eq:assoc} (through degree $6$);
\item the displayed terms \eqref{eq:z1}--\eqref{eq:z6}, entered
independently as nested brackets and expanded by \eqref{eq:bracketexpand};
\item the Zassenhaus exponents computed three ways, by residual extraction
\eqref{eq:residual}, by the Lie recurrence \eqref{eq:frec}, and by the
forward partition recurrence \eqref{eq:Cpartitions}, compared with each
other and with the explicit enumeration of the weighted-tree formula
\eqref{eq:treeformula} (through degree $10$), followed by reconstruction of
the ordered product and comparison with $e^{t(X+Y)}$ (through degree
$12$), and by the displayed $C_2,C_3,C_4$;
\item the ordered-simplex expansion \eqref{eq:simplex} (through degree
$10$), the symmetric splitting's oddness and cubic coefficient
\eqref{eq:S2log} (through degree $12$), Campbell's conjugation series
\eqref{eq:campbell}, the Bernoulli inverse identity $\phi\beta=1$, and the
invariant-metric even-series identity \eqref{eq:metriceven} (through degree
$10$).
\end{itemize}
The supplied runs report $20$, $30$, and $26$ passed named checks
respectively; the script \texttt{code/run\_all\_verifiers.py} in the
article's directory reruns all three and records their reports. These are
exact finite algebraic verifications, not floating-point tolerance tests
and not a claim of formal proof-assistant certification. Finite-degree
agreement does not prove an all-orders identity, an analytic convergence
theorem, or a domain statement for an unbounded operator; those are
established by the proofs above.

\subsection{Two commutator encodings}
The right-nested data use exactly $R(w)$ of \eqref{eq:rightbracket}, so
that $Z_n=\sum_{|w|=n}\frac{c(w)}nR(w)$. This encoding is redundant: every
word ending in two equal letters represents zero, and the remaining brackets
are not claimed to be linearly independent.

For compact tables we use the standard Lyndon bracketing. Order the
alphabet by $X<Y$ and words lexicographically, a proper prefix preceding
its extensions. A nonempty word is \emph{Lyndon} if it is strictly smaller
than each of its nonempty proper suffixes. Define $L(X)=X$, $L(Y)=Y$, and
for a Lyndon word $w$ of length at least two write $w=uv$ with $v$ its
longest proper suffix that is Lyndon, and set
\begin{equation}\label{eq:lyndon}
 L(w)=[L(u),L(v)].
\end{equation}
(The standard factorization theorem guarantees that $u$ is again Lyndon;
the table generator applies the recursion to $u$ and $v$ in any case, so
every entry is an explicitly specified bracket polynomial.) For example
$L(\mathrm{XXY})=[X,[X,Y]]$ and $L(\mathrm{XYY})=[[X,Y],Y]$. Thus
$L(\mathrm{XYY})$ is \emph{not} $R(\mathrm{XYY})$, which is zero; confusing
the two conventions would corrupt many coefficients. No theorem about
Lyndon bases is needed to verify the tables: the generator expands each
tabulated bracket polynomial into words and requires the result to equal
the associative polynomial exactly. A table is therefore a certificate for
an explicit finite commutator identity whether or not one invokes the
general basis theorem. The counts in \cref{tab:counts} are numbers of
nonzero entries in these encodings, not dimensions of homogeneous free Lie
algebras and not counts of terms in a minimal basis.

\subsection{Data files and regeneration}
The principal machine-readable files in the article's \texttt{data/}
directory are listed in \cref{tab:files}. In an associative file a row
gives the coefficient of the literal word; in a \texttt{lyndon} file it
multiplies $L(w)$, and in the right-commutator file it multiplies $R(w)$.
Numerators and denominators are separate integers, and missing rows have
coefficient zero. The tree-polynomial file gives each $C_n$ as a polynomial
in the $A_j$ of \eqref{eq:An}; a key such as \texttt{2 3} denotes the
ordered product $A_2A_3$.
\begin{table}[htbp]
\centering\small
\begin{tabular}{@{}p{0.42\textwidth}p{0.52\textwidth}@{}}
\toprule
File & Meaning\\
\midrule
\texttt{bch\_associative.csv} & Every nonzero BCH word coefficient $c(w)$ through degree 12.\\
\texttt{bch\_right\_commutators.csv} & Dynkin-projection coefficients $c(w)/n$, with identically zero right brackets omitted.\\
\texttt{bch\_lyndon.csv} & Compact BCH commutator expansion through degree 12.\\
\texttt{zassenhaus\_associative.csv} & Associative coefficients of each $C_n$, $2\leq n\leq12$.\\
\texttt{zassenhaus\_lyndon.csv} & Compact commutator expansions of the same $C_n$.\\
\texttt{zassenhaus\_tree\_polynomials.json} & Each $C_n$, $2\leq n\leq10$, as a polynomial in the ordered $A_j$.\\
\texttt{strang\_lyndon.csv} & $\log(e^{X/2}e^Ye^{X/2})$ through degree 12.\\
\texttt{bernoulli\_plus.csv} & The coefficients $\bplus n=B_n^+/n!$ through $n=12$.\\
\bottomrule
\end{tabular}
\caption{Machine-readable companion expansions (exact rationals).}\label{tab:files}
\end{table}

The coefficient appendix (\cref{app:tables}) is generated from these
files by \texttt{code/make\_tables.py}, which independently recomputes every
printed degree-one to degree-six coefficient by the cut formula
\eqref{eq:wordcut} and by expansion of the displayed brackets, reconstructs
every printed Lyndon table from its brackets, and refuses to write the
appendix if any comparison fails. From the article's directory:
\begin{lstlisting}[language=bash]
python code/make_tables.py            # regenerate tex/appendix_tables.tex with checks
python code/run_all_verifiers.py      # rerun the three exact verifiers
latexmk -pdf -interaction=nonstopmode -halt-on-error bch_combined.tex
\end{lstlisting}
Python 3.10 or later with the standard library suffices; a standard TeX
installation with the packages named in the preamble compiles the source
without a bibliography processor, shell escape, or external files. The
verifiers accept larger cutoffs, but the number of words grows
exponentially with the degree; the finite coefficient formulas and
recurrences of this article, not the chosen degree limit, specify all
orders.

\appendix
\section{The ordered Poincar\'e--Birkhoff--Witt theorem, enveloping Hopf algebras, and the completed tensor product}\label{app:pbw}
This appendix proves the injectivity statement used in \cref{sec:universal},
constructs the Hopf structure on an arbitrary enveloping algebra, and shows
concretely why the coproduct of an infinite series needs a completed tensor
product.

\subsection{The universal enveloping algebra and its ordered normal forms}
Let $\h$ be a Lie algebra over $\kk$. Its universal enveloping algebra is
\begin{equation}\label{eq:Udef}
 U(\h)=T(\h)\big/\bigl\langle u\otimes v-v\otimes u-[u,v]:u,v\in\h\bigr\rangle,
\end{equation}
the quotient of the tensor algebra on the vector space $\h$ by the
two-sided ideal generated by the displayed relations. A linear map from
$\h$ to an associative algebra $A$ extends uniquely to an algebra
homomorphism on $T(\h)$, and this extension annihilates the ideal exactly
when the map preserves brackets (with $A$ given its commutator bracket);
this is the universal property. It remains to see that $\h\to U(\h)$ is
injective.

\begin{theorem}[Ordered PBW theorem]\label{thm:pbw}
Let $(e_i)_{i\in I}$ be a totally ordered basis of $\h$. The ordered
monomials $e_{i_1}e_{i_2}\cdots e_{i_m}$ with $i_1\leq i_2\leq\cdots\leq i_m$,
together with the empty monomial $1$, form a vector-space basis of $U(\h)$.
In particular the natural map $\h\to U(\h)$ is injective.
\end{theorem}
\begin{proof}
Work in $T(\h)$ with the basis of words in the letters $e_i$. Consider the
rewriting rule, applicable to any adjacent inverted pair $j>i$ inside a
word,
\begin{equation}\label{eq:rewrite}
 \cdots e_je_i\cdots\ \longmapsto\ \cdots e_ie_j\cdots+\cdots[e_j,e_i]\cdots,
\end{equation}
where $[e_j,e_i]$ is expanded as a finite linear combination of basis
letters. Assign to a word the pair (length, number of inversions), ordered
lexicographically with length first. The swapped word has the same length
and one inversion fewer, and every word in the bracket term is shorter. So
each application of the rule replaces a word by a linear combination of
words of strictly smaller measure. Since the measure takes values in a
well-ordered set, every sequence of rewritings of a fixed word terminates,
and the result is a linear combination of ordered words (words with no
inverted adjacent pair). This proves that ordered words span $U(\h)$.

We prove that the final result does not depend on the choices made
(confluence), by induction on the measure of the starting word $w$. If two
different first steps are available for $w$ at disjoint adjacent pairs,
performing the other one next in either case leads to the same linear
combination, whose words all have smaller measure than $w$; by the
induction hypothesis each of those words has a unique normal form, hence
so does the combination, and both first choices lead to it. If the two
available first steps overlap, the overlap is a factor $e_ke_je_i$ with
$k>j>i$ inside $w=p\,e_ke_je_i\,q$. Reducing the left pair first and then
continuing to order the three original letters gives, before any
normalization of the shorter terms,
\[
 p\bigl(e_ie_je_k+[e_j,e_i]e_k+e_j[e_k,e_i]+[e_k,e_j]e_i\bigr)q,
\]
and reducing the right pair first gives
\[
 p\bigl(e_ie_je_k+e_i[e_k,e_j]+[e_k,e_i]e_j+e_k[e_j,e_i]\bigr)q.
\]
(Each line records the sequence of swaps that brings $e_ke_je_i$ to
$e_ie_je_k$, collecting the bracket terms produced.) Their difference is
$p\,\Xi\,q$ with
\[
 \Xi=[e_j,e_i]e_k-e_k[e_j,e_i]+e_j[e_k,e_i]-[e_k,e_i]e_j+[e_k,e_j]e_i-e_i[e_k,e_j].
\]
Every word in $\Xi$ has length two, hence smaller measure than $e_ke_je_i$,
and the same holds after attaching $p$ and $q$. Applying the rule to each
length-two product $ab-ba$ with $a,b\in\h$ (which, expanded in the basis and
reduced, has normal form equal to the normal form of $[a,b]$; this is the
content of \eqref{eq:rewrite} for a single pair, applied bilinearly)
shows that $p\,\Xi\,q$ has the same normal form as
\[
 p\bigl([[e_j,e_i],e_k]+[e_j,[e_k,e_i]]+[[e_k,e_j],e_i]\bigr)q
 =p\bigl([[e_j,e_i],e_k]+[[e_i,e_k],e_j]+[[e_k,e_j],e_i]\bigr)q=0,
\]
by antisymmetry and the Jacobi identity in $\h$, where all normal forms of
the shorter words involved are unique by the induction hypothesis. Hence
the two first choices lead to combinations with the same normal form.
Since every rule has length two, there are no other overlapping patterns.
Thus any two complete reductions of $w$ agree: compare their first steps
as above, then apply the induction hypothesis to all words below $w$.
Extending by linearity yields a well-defined \emph{normal-form map}
$\nu:T(\h)\to T(\h)$ whose image is the span of ordered words and which
fixes ordered words.

Each rewriting step changes an element by an element of the ideal
$J$ in \eqref{eq:Udef}: it replaces $p\,e_je_i\,q$ by
$p(e_ie_j+[e_j,e_i])q$, and the difference is $p(e_je_i-e_ie_j-[e_j,e_i])q\in J$.
Hence $x-\nu(x)\in J$ for every $x$. Conversely, every generator
$r=uv-vu-[u,v]$ of $J$, multiplied on either side by words, has normal form
zero: by bilinearity it suffices to take $u=e_j$, $v=e_i$, and for $j>i$
the two sides $e_je_i$ and $e_ie_j+[e_j,e_i]$ of the rule have the same
normal form by construction, while for $j=i$ the relation is $0$ and for
$j<i$ it is the negative of a rule instance. Thus $\nu(J)=0$, and
$\nu$ descends to a linear map $U(\h)=T(\h)/J\to\mathrm{span}(\text{ordered words})$
that is surjective (ordered words are fixed) and injective (if $\nu(x)=0$
then $x=x-\nu(x)\in J$). Therefore the ordered words form a basis of
$U(\h)$. Words of length one are ordered, so the basis vectors $e_i$ remain
linearly independent in $U(\h)$, which proves injectivity of $\h\to U(\h)$.
\end{proof}
The proof uses only bilinearity, antisymmetry, and the Jacobi identity; it
is independent of the BCH theorem and of any matrix realization, so its use
in \cref{sec:universal} introduces no circularity.

\subsection{The free Lie algebra and the free associative algebra}
Let $F(X,Y)$ be the abstract free Lie algebra on two generators, the
quotient of the free nonassociative algebra on bracketed words by the
ideal generated by antisymmetry and Jacobi. A Lie homomorphism out of
$F(X,Y)$ is the same as a choice of two elements of the target: substitution
defines the map, and the imposed relations are precisely what makes it well
defined. Hence the universal property of $U(F(X,Y))$ says that an
associative homomorphism from it to any associative algebra is exactly a
choice of two elements, which is also the universal property of the free
associative algebra $\kk\langle X,Y\rangle=T(V)$; the mutually inverse
homomorphisms exchanging the generators identify the two, as in the proof
of \cref{prop:enveloping}. By \cref{thm:pbw}, $F(X,Y)$ injects into
$U(F(X,Y))\cong T(V)$, and its image is the Lie subalgebra $\cL(X,Y)$
generated by $X,Y$. This identifies the concrete $\cL(X,Y)$ used throughout
the article with the abstract free Lie algebra, and shows that an equality
of bracket polynomials verified as an equality of words in $T(V)$ is an
identity in every Lie algebra.

\subsection{The enveloping Hopf structure}
On the tensor algebra $T(\h)$ set, for $u\in\h$,
\begin{equation}\label{eq:UHopf}
 \Delta u=u\otimes1+1\otimes u,\qquad\eps(u)=0,\qquad S(u)=-u,
\end{equation}
and extend $\Delta,\eps$ as algebra homomorphisms and $S$ as an
antihomomorphism. For a defining relation $r=uv-vu-[u,v]$ of
\eqref{eq:Udef}, direct expansion gives
\[
 \Delta r=(uv-vu)\otimes1+1\otimes(uv-vu)-[u,v]\otimes1-1\otimes[u,v]
 =r\otimes1+1\otimes r,
\]
because the mixed terms $u\otimes v+v\otimes u-v\otimes u-u\otimes v$ cancel;
also $\eps(r)=0$ and $S(r)=vu-uv+[u,v]=-r$. Hence $\Delta(J)\subseteq
J\otimes T(\h)+T(\h)\otimes J$, $\eps(J)=0$, and $S(J)\subseteq J$, so the
three maps descend to $U(\h)$. Coassociativity and the counit identities
hold on generators, hence on all of $U(\h)$ by multiplicativity; the
antipode identities $m(S\otimes\id)\Delta=\eta\eps=m(\id\otimes S)\Delta$
hold on generators ($-u+u=0=\eps(u)$) and propagate to products by the
inductive argument of \cref{lem:antipode}, which used only that $\Delta$
is multiplicative and $S$ antimultiplicative. Thus $U(\h)$ is a Hopf
algebra, and under $U(\cL(X,Y))\cong T(V)$ this is exactly the unshuffle
Hopf algebra of \cref{sec:hopf}; degree completion gives its formal
power-series version. For a general $\h$, tensor length is not a grading
of $U(\h)$, since a bracket relation changes length; formal exponentials
are then taken in $U(\h)[[t]]$ with arguments $tx,ty$, which preserves every
coefficient identity while avoiding an inappropriate grading.

\subsection{Why the coproduct needs a completed tensor product}
Restrict to the one-generator subalgebra and consider
$F=\sum_{n\geq0}X^n$. By \eqref{eq:unshuffle},
\[
 \coeff{X^p\otimes X^q}\Delta F=\binom{p+q}{p}\qquad(p,q\geq0),
\]
since $X^{p+q}$ has $\binom{p+q}p$ subsets of positions of size $p$. The
leading $(N+1)\times(N+1)$ block of this coefficient matrix is
$C=(\binom{p+q}p)_{0\leq p,q\leq N}$. By Vandermonde's identity,
$\binom{p+q}p=\sum_k\binom pk\binom qk$ (compare coefficients of $t^p$ in
$(1+t)^p(1+t)^q=(1+t)^{p+q}$, using $\binom{q}{p-k}=\binom q{q-p+k}$ and
reindexing), so $C=PP^{\mathsf T}$ with $P_{pk}=\binom pk$ for $k\leq p$ and
$P_{pk}=0$ otherwise. Since $P$ is lower triangular with diagonal entries
$1$, $\det C=(\det P)^2=1$, so $C$ has rank $N+1$ for every $N$. On the
other hand, if $\Delta F$ were a finite sum $\sum_{j=1}^mA_j\otimes B_j$ of
separated tensors with $A_j,B_j\in\kk[[X]]$, its coefficient matrix would
be $\sum_j\alpha_j\beta_j^{\mathsf T}$ with $\alpha_j,\beta_j$ the
coefficient vectors of $A_j,B_j$, of rank at most $m$. Taking $N\geq m$
gives a contradiction. Thus $\Delta F$ does not lie in the algebraic tensor
product $\widehat T\otimes\widehat T$, while it does lie in the completed
tensor product used in \cref{sec:hopf}, whose elements are arbitrary
families of coefficients indexed by pairs of words.

\section{Coefficient certificates and additional complete degrees}\label{app:tables}
Every table in this appendix was regenerated from the exact data files by
\texttt{code/make\_tables.py}, which recomputed each printed degree-one to
degree-six coefficient by the finite cut formula \eqref{eq:wordcut}, expanded
the displayed brackets \eqref{eq:z1}--\eqref{eq:z6} and
\eqref{eq:c2}--\eqref{eq:c4} by \eqref{eq:bracketexpand}, reconstructed every
Lyndon table below from its brackets, and expanded the tree polynomials of
\cref{thm:tree} into words, comparing all results exactly before writing
the file.

\subsection{All nonzero associative BCH coefficients through degree six}
For each listed word $w$ the table gives $c(w)$ of \eqref{eq:wordcut}; every
unlisted word of degree at most six has coefficient zero, for both the cut
formula and the expansion of \eqref{eq:z1}--\eqref{eq:z6}. Among the
$2+4+8+16+32+64=126$ words there are $72$ nonzero entries. This is a complete
finite certificate for the displayed polynomial through degree six.
\begingroup\small
\begin{longtable}{@{}lrlrlr@{}}
\caption{All nonzero word coefficients $c(w)$ of $Z_1,\ldots,Z_6$.}\label{tab:words6}\\
\toprule
Word $w$ & $c(w)$ & Word $w$ & $c(w)$ & Word $w$ & $c(w)$\\
\midrule\endfirsthead
\multicolumn{6}{c}{\tablename\ \thetable\ (continued)}\\
\toprule
Word $w$ & $c(w)$ & Word $w$ & $c(w)$ & Word $w$ & $c(w)$\\
\midrule\endhead
\bottomrule\endfoot
\multicolumn{6}{@{}l}{\textit{Degree 1}}\\[2pt]
\texttt{X} & $1$ & \texttt{Y} & $1$ &  &  \\
\addlinespace[3pt]
\multicolumn{6}{@{}l}{\textit{Degree 2}}\\[2pt]
\texttt{XY} & $\tfrac{1}{2}$ & \texttt{YX} & $-\tfrac{1}{2}$ &  &  \\
\addlinespace[3pt]
\multicolumn{6}{@{}l}{\textit{Degree 3}}\\[2pt]
\texttt{XXY} & $\tfrac{1}{12}$ & \texttt{XYY} & $\tfrac{1}{12}$ & \texttt{YXY} & $-\tfrac{1}{6}$ \\
\texttt{XYX} & $-\tfrac{1}{6}$ & \texttt{YXX} & $\tfrac{1}{12}$ & \texttt{YYX} & $\tfrac{1}{12}$ \\
\addlinespace[3pt]
\multicolumn{6}{@{}l}{\textit{Degree 4}}\\[2pt]
\texttt{XXYY} & $\tfrac{1}{24}$ & \texttt{YXYX} & $\tfrac{1}{12}$ &  &  \\
\texttt{XYXY} & $-\tfrac{1}{12}$ & \texttt{YYXX} & $-\tfrac{1}{24}$ &  &  \\
\addlinespace[3pt]
\multicolumn{6}{@{}l}{\textit{Degree 5}}\\[2pt]
\texttt{XXXXY} & $-\tfrac{1}{720}$ & \texttt{XYXYY} & $-\tfrac{1}{120}$ & \texttt{YXYXY} & $\tfrac{1}{30}$ \\
\texttt{XXXYX} & $\tfrac{1}{180}$ & \texttt{XYYXX} & $-\tfrac{1}{120}$ & \texttt{YXYYX} & $-\tfrac{1}{120}$ \\
\texttt{XXXYY} & $\tfrac{1}{180}$ & \texttt{XYYXY} & $-\tfrac{1}{120}$ & \texttt{YXYYY} & $\tfrac{1}{180}$ \\
\texttt{XXYXX} & $-\tfrac{1}{120}$ & \texttt{XYYYX} & $\tfrac{1}{180}$ & \texttt{YYXXX} & $\tfrac{1}{180}$ \\
\texttt{XXYXY} & $-\tfrac{1}{120}$ & \texttt{XYYYY} & $-\tfrac{1}{720}$ & \texttt{YYXXY} & $-\tfrac{1}{120}$ \\
\texttt{XXYYX} & $-\tfrac{1}{120}$ & \texttt{YXXXX} & $-\tfrac{1}{720}$ & \texttt{YYXYX} & $-\tfrac{1}{120}$ \\
\texttt{XXYYY} & $\tfrac{1}{180}$ & \texttt{YXXXY} & $\tfrac{1}{180}$ & \texttt{YYXYY} & $-\tfrac{1}{120}$ \\
\texttt{XYXXX} & $\tfrac{1}{180}$ & \texttt{YXXYX} & $-\tfrac{1}{120}$ & \texttt{YYYXX} & $\tfrac{1}{180}$ \\
\texttt{XYXXY} & $-\tfrac{1}{120}$ & \texttt{YXXYY} & $-\tfrac{1}{120}$ & \texttt{YYYXY} & $\tfrac{1}{180}$ \\
\texttt{XYXYX} & $\tfrac{1}{30}$ & \texttt{YXYXX} & $-\tfrac{1}{120}$ & \texttt{YYYYX} & $-\tfrac{1}{720}$ \\
\addlinespace[3pt]
\multicolumn{6}{@{}l}{\textit{Degree 6}}\\[2pt]
\texttt{XXXXYY} & $-\tfrac{1}{1440}$ & \texttt{XYXYYY} & $\tfrac{1}{360}$ & \texttt{YXYYYX} & $-\tfrac{1}{360}$ \\
\texttt{XXXYXY} & $\tfrac{1}{360}$ & \texttt{XYYXXY} & $-\tfrac{1}{240}$ & \texttt{YYXXXX} & $\tfrac{1}{1440}$ \\
\texttt{XXXYYY} & $\tfrac{1}{360}$ & \texttt{XYYXYY} & $-\tfrac{1}{240}$ & \texttt{YYXXYX} & $\tfrac{1}{240}$ \\
\texttt{XXYXXY} & $-\tfrac{1}{240}$ & \texttt{XYYYXY} & $\tfrac{1}{360}$ & \texttt{YYXYXX} & $\tfrac{1}{240}$ \\
\texttt{XXYXYY} & $-\tfrac{1}{240}$ & \texttt{YXXXYX} & $-\tfrac{1}{360}$ & \texttt{YYXYYX} & $\tfrac{1}{240}$ \\
\texttt{XXYYXY} & $-\tfrac{1}{240}$ & \texttt{YXXYXX} & $\tfrac{1}{240}$ & \texttt{YYYXXX} & $-\tfrac{1}{360}$ \\
\texttt{XXYYYY} & $-\tfrac{1}{1440}$ & \texttt{YXXYYX} & $\tfrac{1}{240}$ & \texttt{YYYXYX} & $-\tfrac{1}{360}$ \\
\texttt{XYXXXY} & $\tfrac{1}{360}$ & \texttt{YXYXXX} & $-\tfrac{1}{360}$ & \texttt{YYYYXX} & $\tfrac{1}{1440}$ \\
\texttt{XYXXYY} & $-\tfrac{1}{240}$ & \texttt{YXYXYX} & $-\tfrac{1}{60}$ &  &  \\
\texttt{XYXYXY} & $\tfrac{1}{60}$ & \texttt{YXYYXX} & $\tfrac{1}{240}$ &  &  \\
\end{longtable}\endgroup

\subsection{Complete seventh and eighth BCH degrees}
The entries specify exactly $Z_n=\sum_wc_w\,L(w)$ for $n=7,8$ in the standard
Lyndon bracketing \eqref{eq:lyndon}; there are no omitted nonzero terms in
this convention. The CSV file continues these expansions through degree
twelve.

\begingroup\small
\renewcommand{\arraystretch}{1.3}
\begin{longtable}{@{}lr@{\hspace{3em}}lr@{}}
\caption{Complete BCH homogeneous polynomials of degrees seven and eight in Lyndon bracketing.}\label{tab:bch78}\\
\toprule
\multicolumn{2}{c}{$Z_7$} & \multicolumn{2}{c}{$Z_8$} \\
$w$ & coefficient of $L(w)$ & $w$ & coefficient of $L(w)$ \\
\midrule\endfirsthead
\multicolumn{4}{c}{\tablename\ \thetable\ (continued)}\\
\toprule
\multicolumn{2}{c}{$Z_7$} & \multicolumn{2}{c}{$Z_8$} \\
$w$ & coefficient of $L(w)$ & $w$ & coefficient of $L(w)$ \\
\midrule\endhead
\bottomrule\endfoot
\texttt{XXXXXXY} & $\tfrac{1}{30240}$ & \texttt{XXXXXXYY} & $\tfrac{1}{60480}$ \\
\texttt{XXXXXYY} & $-\tfrac{1}{5040}$ & \texttt{XXXXXYXY} & $-\tfrac{1}{15120}$ \\
\texttt{XXXXYXY} & $\tfrac{1}{10080}$ & \texttt{XXXXXYYY} & $-\tfrac{1}{10080}$ \\
\texttt{XXXXYYY} & $\tfrac{1}{3780}$ & \texttt{XXXXYXXY} & $\tfrac{1}{20160}$ \\
\texttt{XXXYXXY} & $\tfrac{1}{10080}$ & \texttt{XXXXYXYY} & $-\tfrac{1}{20160}$ \\
\texttt{XXXYXYY} & $\tfrac{1}{1680}$ & \texttt{XXXXYYXY} & $\tfrac{1}{2520}$ \\
\texttt{XXXYYXY} & $\tfrac{1}{1260}$ & \texttt{XXXXYYYY} & $\tfrac{23}{120960}$ \\
\texttt{XXXYYYY} & $\tfrac{1}{3780}$ & \texttt{XXXYXXYY} & $\tfrac{1}{4032}$ \\
\texttt{XXYXXYY} & $\tfrac{1}{2016}$ & \texttt{XXXYXYXY} & $-\tfrac{1}{10080}$ \\
\texttt{XXYXYXY} & $-\tfrac{1}{5040}$ & \texttt{XXXYXYYY} & $\tfrac{13}{30240}$ \\
\texttt{XXYXYYY} & $\tfrac{13}{15120}$ & \texttt{XXXYYXYY} & $\tfrac{1}{20160}$ \\
\texttt{XXYYXYY} & $\tfrac{1}{10080}$ & \texttt{XXXYYYXY} & $-\tfrac{1}{3024}$ \\
\texttt{XXYYYXY} & $-\tfrac{1}{1512}$ & \texttt{XXXYYYYY} & $-\tfrac{1}{10080}$ \\
\texttt{XXYYYYY} & $-\tfrac{1}{5040}$ & \texttt{XXYXYXYY} & $\tfrac{1}{2520}$ \\
\texttt{XYXYXYY} & $\tfrac{1}{1260}$ & \texttt{XXYXYYYY} & $-\tfrac{1}{4032}$ \\
\texttt{XYXYYYY} & $-\tfrac{1}{2016}$ & \texttt{XXYYXYYY} & $-\tfrac{1}{10080}$ \\
\texttt{XYYXYYY} & $-\tfrac{1}{5040}$ & \texttt{XXYYYYYY} & $\tfrac{1}{60480}$ \\
\texttt{XYYYYYY} & $\tfrac{1}{30240}$ & -- & -- \\
\end{longtable}
\endgroup

\subsection{Complete fifth and sixth Zassenhaus exponents}
These are the next two complete exponents after the terms displayed on the
source page. A coefficient multiplies $L(w)$, not $R(w)$. Every entry follows
from \eqref{eq:residual}, from \eqref{eq:frec}, and from
\eqref{eq:treeformula}; the three constructions agree exactly.

\begingroup\small
\renewcommand{\arraystretch}{1.3}
\begin{longtable}{@{}lr@{\hspace{3em}}lr@{}}
\caption{Complete fifth and sixth Zassenhaus exponents in Lyndon bracketing.}\label{tab:zass56}\\
\toprule
\multicolumn{2}{c}{$C_5$} & \multicolumn{2}{c}{$C_6$} \\
$w$ & coefficient of $L(w)$ & $w$ & coefficient of $L(w)$ \\
\midrule\endfirsthead
\multicolumn{4}{c}{\tablename\ \thetable\ (continued)}\\
\toprule
\multicolumn{2}{c}{$C_5$} & \multicolumn{2}{c}{$C_6$} \\
$w$ & coefficient of $L(w)$ & $w$ & coefficient of $L(w)$ \\
\midrule\endhead
\bottomrule\endfoot
\texttt{XXXXY} & $\tfrac{1}{120}$ & \texttt{XXXXXY} & $-\tfrac{1}{720}$ \\
\texttt{XXXYY} & $-\tfrac{1}{30}$ & \texttt{XXXXYY} & $\tfrac{1}{144}$ \\
\texttt{XXYXY} & $-\tfrac{1}{60}$ & \texttt{XXXYYY} & $-\tfrac{1}{72}$ \\
\texttt{XXYYY} & $\tfrac{1}{20}$ & \texttt{XXYYYY} & $\tfrac{1}{72}$ \\
\texttt{XYXYY} & $-\tfrac{1}{20}$ & \texttt{XYXYYY} & $-\tfrac{1}{72}$ \\
\texttt{XYYYY} & $-\tfrac{1}{30}$ & \texttt{XYYYYY} & $-\tfrac{1}{144}$ \\
\end{longtable}
\endgroup

\subsection{Tree polynomials for the Zassenhaus exponents}
Each $C_n$ below is written as a polynomial in the ordered raw polynomials
$A_j$ of \eqref{eq:An}, obtained by summing the weights $\omega(T)$ of
\cref{thm:tree} over trees with the same leaf sequence. Products are
noncommutative and are read from left to right; $A_2^2A_3$ and $A_3A_2^2$ are
different polynomials.
\begin{center}\small
\renewcommand{\arraystretch}{1.3}
\begin{tabular}{@{}r>{\raggedright\arraybackslash}p{0.9\textwidth}@{}}
\toprule
$n$ & $C_n$ in terms of the $A_j$\\
\midrule
2 & $A_{2}$\\
3 & $A_{3}$\\
4 & $A_{4}$ $-\tfrac{1}{2}A_{2}^{2}$\\
5 & $A_{5}$ $-A_{2}A_{3}$\\
6 & $A_{6}$ $+\tfrac{1}{3}A_{2}^{3}$ $-A_{2}A_{4}$ $-\tfrac{1}{2}A_{3}^{2}$\\
7 & $A_{7}$ $+\tfrac{1}{2}A_{2}^{2}A_{3}$ $-A_{2}A_{5}$ $-A_{3}A_{4}$ $+\tfrac{1}{2}A_{3}A_{2}^{2}$\\
8 & $A_{8}$ $-\tfrac{1}{4}A_{2}^{4}$ $+\tfrac{3}{4}A_{2}^{2}A_{4}$ $-A_{2}A_{6}$ $-A_{3}A_{5}$ $+A_{3}A_{2}A_{3}$ $-\tfrac{1}{2}A_{4}^{2}$ $+\tfrac{1}{4}A_{4}A_{2}^{2}$\\
9 & $A_{9}$ $-\tfrac{2}{3}A_{2}^{3}A_{3}$ $+A_{2}^{2}A_{5}$ $-A_{2}A_{7}$ $+\tfrac{1}{3}A_{3}^{3}$ $-A_{3}A_{6}$ $-\tfrac{1}{3}A_{3}A_{2}^{3}$ $+A_{3}A_{2}A_{4}$ $-A_{4}A_{5}$ $+A_{4}A_{2}A_{3}$\\
10 & $A_{10}$ $+\tfrac{1}{5}A_{2}^{5}$ $-\tfrac{2}{3}A_{2}^{3}A_{4}$ $-\tfrac{1}{4}A_{2}^{2}A_{3}^{2}$ $+A_{2}^{2}A_{6}$ $+\tfrac{1}{2}A_{2}A_{3}A_{5}$ $-\tfrac{1}{2}A_{2}A_{3}A_{2}A_{3}$ $-A_{2}A_{8}$ $+\tfrac{1}{2}A_{3}^{2}A_{4}$ $-\tfrac{1}{4}A_{3}^{2}A_{2}^{2}$ $-A_{3}A_{7}$ $-\tfrac{1}{2}A_{3}A_{2}^{2}A_{3}$ $+A_{3}A_{2}A_{5}$ $-A_{4}A_{6}$ $-\tfrac{1}{3}A_{4}A_{2}^{3}$ $+A_{4}A_{2}A_{4}$ $+\tfrac{1}{2}A_{4}A_{3}^{2}$ $-\tfrac{1}{2}A_{5}^{2}$ $+\tfrac{1}{2}A_{5}A_{2}A_{3}$\\
\bottomrule
\end{tabular}
\end{center}

\subsection{Sizes of the exported expansions}
\begin{table}[htbp]
\centering\small
\begin{tabular}{@{}rrrrr@{}}
\toprule
Degree & BCH words & BCH Lyndon entries & $C_n$ words & $C_n$ Lyndon entries\\
\midrule
1 & 2 & 2 & 0 & 0 \\
2 & 2 & 1 & 2 & 1 \\
3 & 6 & 2 & 6 & 2 \\
4 & 4 & 1 & 12 & 3 \\
5 & 30 & 6 & 30 & 6 \\
6 & 28 & 5 & 42 & 6 \\
7 & 126 & 18 & 126 & 18 \\
8 & 124 & 17 & 234 & 29 \\
9 & 390 & 55 & 510 & 56 \\
10 & 388 & 55 & 954 & 95 \\
11 & 2046 & 186 & 2046 & 186 \\
12 & 2044 & 185 & 3986 & 326 \\
\bottomrule
\end{tabular}
\caption{Numbers of nonzero coefficients in the supplied encodings. The
degree-one Zassenhaus factor $e^{Y}$ is separated from the exponents $C_n$,
$n\geq2$.}\label{tab:counts}
\end{table}

\FloatBarrier
\section{Claim-by-claim coverage of the source page}\label{app:coverage}
The following concordance refers to revision 1368909113 of \cite{wiki}. It
tracks mathematical content rather than historical prose; repeated displays
of the same identity are grouped, but different hypotheses and convergence
interpretations are distinguished. An entry marked \emph{corrected} or
\emph{qualified} gives the valid statement and identifies why the
unqualified original cannot be asserted; it does not mean that a false
unrestricted assertion has been proved. Theorems appearing only as names
in links are not statements made by the page; the Trotter and higher-order
splitting results are included as additional consequences.

\begingroup\small
\setlength{\tabcolsep}{4pt}
\renewcommand{\arraystretch}{1.18}
\begin{longtable}{@{}>{\raggedright\arraybackslash}p{0.21\textwidth}>{\raggedright\arraybackslash}p{0.40\textwidth}>{\raggedright\arraybackslash}p{0.34\textwidth}@{}}
\caption{Claim-by-claim mathematical coverage.}\label{tab:coverage}\\
\toprule
Source location & Formula or assertion & Proof or qualification here\\
\midrule\endfirsthead
\multicolumn{3}{c}{\tablename\ \thetable\ (continued)}\\
\toprule
Source location & Formula or assertion & Proof or qualification here\\
\midrule\endhead
\bottomrule\endfoot
Introduction & Exponential and logarithm series; $e^Xe^Y=e^Z$ with $Z$ a formal series of iterated commutators; leading terms. & \Cref{lem:explog,thm:wordcut,thm:formalbch,thm:dynkin}; \eqref{eq:z1}--\eqref{eq:z3}.\\
Introduction & Commutator interpretation, antisymmetry, Jacobi identity. & \Cref{sec:scope}, \eqref{eq:derivation}, by associative expansion.\\
Introduction & Small elements give a convergent BCH series; the exponential gives coordinates near the identity; local multiplication is determined by the bracket. & \Cref{thm:convergence,thm:local}.\\
Introduction; explicit forms & Lie correspondence; integration of Lie algebra homomorphisms and representations. & \Cref{thm:localcorr,thm:integration}; simple connectedness stated as the global hypothesis.\\
Dynkin formula & Complete block sum, factorial and total-degree denominators, right-nested convention. & \Cref{thm:dynkin}, \eqref{eq:dynkin}, \eqref{eq:rightbracket}.\\
Dynkin formula & Vanishing of summands whose final block ends in repeated letters. & Last assertion of \Cref{thm:dynkin}, $[a,a]=0$.\\
Dynkin formula & Every displayed BCH term through degree six. & \eqref{eq:z1}--\eqref{eq:z6}, \Cref{prop:lowdegree}, certificate \Cref{tab:words6}.\\
Dynkin formula & All nonlinear terms lie in the ideal generated by $[X,Y]$. & \Cref{thm:formalbch}.\\
Dynkin formula & $Z(Y,X)=-Z(-X,-Y)$; homogeneous interchange signs. & \Cref{prop:symmetry}, \eqref{eq:swap}, \eqref{eq:parity}.\\
Integral form & Poincar\'e integral and the function $\psi(x)=x\log x/(x-1)$. & \Cref{thm:poincare}, \eqref{eq:poincare}, \eqref{eq:psi}.\\
Integral form; note & $\psi(e^y)$ and the Bernoulli coefficients with $B_1=+\tfrac12$. & \eqref{eq:psibeta}, \Cref{prop:bernoulli}.\\
Matrix illustration & Tangent Lie algebra, matrix commutator, matrix exponential. & Opening of \Cref{sec:liegroups}.\\
Matrix illustration & Full associative logarithmic block sum. & \eqref{eq:blocks}, \Cref{thm:wordcut}; absolute convergence in \Cref{thm:convergence}.\\
Matrix illustration & Associative terms $z_1,\ldots,z_4$ and their expression by brackets. & \eqref{eq:assoc2}--\eqref{eq:assoc4}, \Cref{prop:lowdegree}.\\
Matrix illustration; note & Conditions $\norm X+\norm Y<\log2$ and $\norm Z<\log2$; Hilbert--Schmidt norm. & \Cref{thm:convergence}; discussion after \eqref{eq:mercatortail}; submultiplicativity of the HS norm proved in \Cref{sec:convergence}.\\
Matrix illustration & Trace of the BCH logarithm; higher terms traceless. & \Cref{prop:trace}; \emph{qualified} for other branches by \eqref{eq:tracebranch}.\\
Convergence issues & The $\mathfrak{sl}_2(\C)$ matrices, their product, and absence of a traceless logarithm; divergence of BCH. & \Cref{thm:sl2}; all logarithms classified, even conditional convergence excluded.\\
Convergence issues & Nonisomorphic groups with the same Lie algebra. & The $\R$ versus $S^1$ example after \Cref{thm:local}.\\
Convergence issues & Bound $\norm X+\norm Y<\tfrac12\log2$ in any submultiplicative norm. & \Cref{thm:convergence}(iii), distinguished from the grouped bound.\\
Special cases & Commuting variables. & \Cref{prop:commuting}.\\
Special cases; Campbell application & Central commutator: truncation and disentangling without smallness. & \Cref{thm:central}, including the differential proof of global validity.\\
Special cases & Expansion linear in $Y$, hyperbolic cotangent form, Bernoulli coefficients of $\ad_X^nY$. & \Cref{cor:linearY}; all higher $Y$-orders in \eqref{eq:bidegree}, \eqref{eq:Hqrec}, \eqref{eq:H2kernel}.\\
Special cases & $[X,Y]=sY$: no terms with two or more $Y$'s; closed form; excluded poles; removable $s=0$. & \Cref{thm:eigen}, \Cref{ex:resonance}; Taylor convergence versus global resummation distinguished.\\
Special cases & Eigenvector braiding and adjoint dilation. & \eqref{eq:eigenbraid}, \Cref{thm:campbell}.\\
Existence results & Formal exp/log bijection; homogeneous grouping. & \Cref{lem:explog,thm:wordcut}.\\
Existence results & Coproduct on noncommuting series. & \eqref{eq:coproduct}, \eqref{eq:unshuffle}; \emph{qualified}: completed tensor product, with the counterexample in \Cref{app:pbw}.\\
Existence results & Exponential is group-like iff exponent is primitive; group-like elements form a group. & \Cref{lem:grouplike}, with constant term $1$ in the definition.\\
Existence results & Friedrichs' characterization of primitive elements; BCH follows. & \Cref{lem:dsw,thm:friedrichs,thm:formalbch}.\\
Existence results & Free enveloping algebra is the free associative algebra; natural Hopf structure; validity in any characteristic-zero Lie algebra. & \Cref{prop:enveloping}, \Cref{sec:universal}, \Cref{app:pbw}.\\
Existence results & Existence of a recursive Lie-only construction. & \Cref{thm:homrec}, independent of the Hopf proof.\\
Zassenhaus formula & Ordered product; displayed $C_2,C_3,C_4$ including the left-nested $C_4$. & \Cref{thm:zassenhaus}, \Cref{prop:zassdisplayed}.\\
Zassenhaus formula & Compact product $e^{-tX}e^{t(X+Y)}=\prod e^{t^nC_n}$. & Convention $C_1=Y$, after \Cref{thm:zassenhaus}.\\
Zassenhaus formula & All later exponents are homogeneous Lie polynomials; recursive determination; the ellipsis. & \Cref{thm:zassenhaus}, \Cref{prop:zassdiff}, \Cref{thm:tree}, \Cref{tab:zass56}; convergence in \Cref{thm:zconv}.\\
Zassenhaus formula & Suzuki--Trotter consequence. & \Cref{sec:splitting}: \Cref{thm:trotter,prop:symmetric,thm:suzuki} with complete error series.\\
Campbell identity & $\Ad_{e^X}=e^{\ad_X}$, all iterated commutators, and the defining differential equation. & \Cref{thm:campbell}, \eqref{eq:adbinomial}.\\
Campbell identity & General conjugation and braiding identities; central reduction. & \eqref{eq:conjexp}, \eqref{eq:braiding}, \eqref{eq:centralconj}.\\
Campbell application & Differentiation of $e^{sX}e^{sY}$ and its solution. & Proof of \Cref{thm:central}.\\
Infinitesimal case & Full series for $e^{-X}\dd e^X$. & \eqref{eq:leftdexp}, \eqref{eq:infinitesimal}.\\
Infinitesimal case & Structure-constant expansion; matrices $M$ and $W$; quotient notation. & \eqref{eq:components}--\eqref{eq:thetaW}; \emph{qualified}: $(I-e^{-M})M^{-1}$ means $\phi(M)$.\\
Infinitesimal case & Identification of $M$ as a vielbein. & \emph{Corrected} by \Cref{prop:coframe}: $W$ is the coframe and $M$ is always singular.\\
Infinitesimal case & Pullback metric $g_{ij}=W^a{}_iW^b{}_jB_{ab}$ and the Killing form as metric. & \Cref{thm:pullback}, \eqref{eq:killing}; \emph{qualified}: Jacobian factors and nondegeneracy required; all orders in \Cref{thm:metric}.\\
Quantum mechanics & $[Q,P]=\ii\hbar I$ and centrality. & \eqref{eq:QP}, \eqref{eq:CCR} on the stated invariant core.\\
Quantum mechanics & Exponentiated canonical relation; exact symmetric three-factor form. & \Cref{thm:weyl}; \Cref{prop:ccrcounter} shows a bare dense-domain commutator is insufficient.\\
Quantum mechanics & $[a^\dagger,a]=-I$; normal ordering of the displacement operator. & \eqref{eq:creation}, \Cref{thm:displacement}, with explicit domain and closure.\\
Quantum mechanics & Product of two displacements, its phase, Heisenberg nilpotence. & \Cref{thm:Dcomposition}, \eqref{eq:Dproduct}, \eqref{eq:heisenberglaw}.\\
Quantum mechanics & Expansions hidden in the displacement exponentials. & \eqref{eq:Dall}--\eqref{eq:laguerre}.\\
\end{longtable}
\endgroup

\subsection*{Summary of essential qualifications}
The tensor product in the formal existence proof must be completed, and
group-like elements have constant term one. The trace identity refers to the
BCH logarithm, not to an arbitrary branch. Smallness is a sufficient
convergence condition, not a necessary one, and resummed validity of a
closed form does not imply convergence of its Taylor series; at nonzero
resonances $2\pi\ii k$ the closed form of shape $X+cY$ fails although a
logarithm exists. In exponential coordinates $M=\ad_X$ is singular and is
not the coframe; $W=\phi(M)$ is the relevant matrix, and its inverse exists
exactly off the resonances. Neither an arbitrary pullback nor the Killing
form is automatically a metric. The alternate Zassenhaus product starting
at $n=1$ needs $C_1=Y$. Finally, an unbounded commutator identity requires
domain and exponentiation hypotheses; concrete realizations, not formal
substitution alone, justify the quantum formulas.

\subsection*{Beyond the finite displays}
Every BCH coefficient is given by \eqref{eq:wordcut} and
\eqref{eq:bidegreeproj}; every bidegree by \eqref{eq:bidegree}; every power
of $Y$ by \eqref{eq:Hqrec} and \eqref{eq:H2kernel}; every Zassenhaus
exponent by \eqref{eq:treeformula}; every Maurer--Cartan component by
\eqref{eq:MW}--\eqref{eq:thetaW}; every invariant-metric coefficient by
\eqref{eq:metriceven}; every splitting error by \eqref{eq:multibch}; and
every oscillator matrix element by \eqref{eq:Dmatrix}. These statements
specify coefficients at arbitrary order; the accompanying tables are finite
illustrations of unlimited formulas, not their claimed endpoint.

\section{Formalization in Lean}\label{app:lean}
The repository directory \texttt{lean/} contains a Lean~4 library
(\texttt{BCH}, built against Mathlib) that formalizes some of the theorems
of this article. Each formal statement is designed to coincide with the
statement printed here, with no additional hypotheses; where the article's
wording needed to be made precise for this purpose, the article was
adjusted rather than the formal statement weakened. This appendix records
the correspondence and the conventions.

\subsection*{Setting and conventions}
A ``real or complex unital Banach algebra with a submultiplicative norm''
is rendered by the Mathlib hypotheses
\begin{center}
\texttt{[RCLike $\mathbb K$] [NormedRing $\mathbb A$]
[NormedAlgebra $\mathbb K$ $\mathbb A$] [CompleteSpace $\mathbb A$]},
\end{center}
where \texttt{RCLike $\mathbb K$} means $\kk=\R$ or $\kk=\C$,
\texttt{NormedRing} is a unital ring with $\norm{ab}\leq\norm a\norm b$ (no
normalization $\norm1=1$ is assumed, in agreement with
\cref{sec:convergence}), and \texttt{CompleteSpace} is completeness. The
Lean bracket $[\![X,Y]\!]$ is the ring commutator $XY-YX$ of
\cref{sec:scope}, scalars act by $t\bullet X$, and \texttt{exp} is
Mathlib's \texttt{NormedSpace.exp}, the sum of the series $\sum_nX^n/n!$.
Mathlib defines this exponential through the unique $\Q$-algebra structure
of $\mathcal A$; inside proofs that structure, and where needed an
$\R$-algebra structure, is obtained by restriction of scalars from $\kk$,
so that no extra instance hypothesis appears in any statement. When the
scalar field does not occur in a statement (\cref{prop:commuting} and the
group-commutator identity of \cref{thm:central}), it is an explicit
argument of the Lean theorem. The adjoint action $\ad_X$ is the bounded
operator \texttt{BCH.ad K X}, a continuous $\kk$-linear map
$\mathcal A\to\mathcal A$, and $e^{s\ad_X}$ is the exponential in the
Banach algebra of bounded operators.

\subsection*{Correspondence}
\begin{center}\small
\renewcommand{\arraystretch}{1.25}
\begin{tabular}{@{}p{0.34\textwidth}p{0.6\textwidth}@{}}
\toprule
Article & Lean declaration (namespace \texttt{BCH}) \\
\midrule
\Cref{prop:commuting} & \texttt{exp\_mul\_exp\_of\_lie\_eq\_zero} \\
\Cref{thm:central}, \eqref{eq:centralconj} & \texttt{central\_conj} \\
\Cref{thm:central}, \eqref{eq:centralt} & \texttt{central\_exp\_mul\_exp} \\
\Cref{thm:central}, \eqref{eq:centralbch}, four identities &
 \texttt{exp\_mul\_exp\_of\_central}, \texttt{exp\_add\_eq\_of\_central},
 \texttt{exp\_half\_mul\_exp\_mul\_exp\_half}, \texttt{exp\_group\_commutator} \\
\Cref{thm:campbell}, $\Ad_{e^X}=e^{\ad_X}$ & \texttt{exp\_mul\_mul\_exp\_neg\_eq\_exp\_ad} \\
\Cref{thm:campbell}, \eqref{eq:campbell} & \texttt{campbell},
 \texttt{exp\_smul\_ad\_apply\_eq\_tsum} (series form) \\
\Cref{thm:campbell}, norm bound & \texttt{tsum\_norm\_ad\_pow\_le} \\
\eqref{eq:conjexp}, \eqref{eq:braiding} & \texttt{exp\_mul\_exp\_mul\_exp\_neg},\newline
 \texttt{exp\_mul\_exp\_eq\_exp\_exp\_ad\_mul} \\
\Cref{thm:eigen}, $e^{tX}Ye^{-tX}=e^{st}Y$ & \texttt{conj\_of\_lie\_eq\_smul} \\
\Cref{thm:eigen}, \eqref{eq:eigenfactor} & \texttt{exp\_smul\_add\_smul\_eq} \\
\Cref{thm:eigen}, \eqref{eq:eigenbch} & \texttt{exp\_mul\_exp\_of\_lie\_eq\_smul} \\
\Cref{thm:eigen}, \eqref{eq:eigenbraid} & \texttt{exp\_mul\_exp\_mul\_exp\_neg\_of\_lie\_eq\_smul},
 \texttt{exp\_mul\_exp\_eq\_of\_lie\_eq\_smul} \\
\bottomrule
\end{tabular}
\end{center}
The functions $q_s$, $\phi$, and $\beta$ of \cref{thm:eigen} are the Lean
definitions \texttt{qs}, \texttt{phi}, \texttt{beta}, with the removable
values $\phi(0)=\beta(0)=1$ built in. The hypothesis
$s\notin2\pi\ii\Z\setminus\{0\}$ of \eqref{eq:eigenbch} is stated as
``$s=0$ or $e^{-s}\neq1$'', which is equivalent over $\C$ and vacuous over
$\R$; the article's statement carries this equivalent form.

\subsection*{Proof method in the formalization}
The article proves the exponential identities of \cref{sec:special} and
\cref{thm:campbell} by uniqueness for linear differential equations. The
formal proofs use an equivalent device that needs only the theorem that a
differentiable function on $\R$ or $\C$ with zero derivative is constant
(Mathlib's \texttt{is\_const\_of\_deriv\_eq\_zero}): an auxiliary
function is exhibited whose derivative vanishes identically and whose
values at $0$ and $1$ are the two sides of the identity. For instance,
\eqref{eq:centralconj} follows from the constancy of
$\sigma\mapsto e^{\sigma X}(Y+(s-\sigma)C)e^{-\sigma X}$,
\eqref{eq:centralt} from that of
$s\mapsto e^{-sY}e^{-sX}e^{s(X+Y)}e^{s^2D}$ with $D+D=C$,
and \eqref{eq:campbell} from that of
$\sigma\mapsto e^{-\sigma\ad_X}\bigl(e^{\sigma X}Ye^{-\sigma X}\bigr)$.
Statements with a general scalar parameter $t$ are deduced from the case
$t=1$ applied to $tX$ and $tY$.

The formal power-series parts of the article (\cref{sec:formal,sec:hopf,sec:dynkin}),
the convergence estimates of \cref{sec:convergence}, and the later sections
are not yet formalized; the finite exact computations of
\cref{sec:computation} remain the machine check of the coefficient
tables.

\begingroup\small
\begin{thebibliography}{99}
\setlength{\itemsep}{3pt}
\setlength{\parskip}{0pt}
\addcontentsline{toc}{section}{References}

\bibitem{wiki}
Wikipedia contributors,
\emph{Baker--Campbell--Hausdorff formula}, English Wikipedia, revision
1368909113 (11 August 2026), consulted 16 September 2026.
\href{https://en.wikipedia.org/w/index.php?oldid=1368909113}{Permanent revision 1368909113}.
This is the coverage target; every proof in the present article is supplied
independently.

\bibitem{dynkin}
E.~B. Dynkin,
Calculation of the coefficients in the Campbell--Hausdorff formula (in Russian),
\emph{Doklady Akademii Nauk SSSR} \textbf{57} (1947), 323--326.

\bibitem{eichler}
M. Eichler,
A new proof of the Baker--Campbell--Hausdorff formula,
\emph{Journal of the Mathematical Society of Japan} \textbf{20} (1968), 23--25.
\href{https://doi.org/10.2969/jmsj/02010023}{doi:10.2969/jmsj/02010023}.

\bibitem{menouspatras}
F. Menous and F. Patras,
Logarithmic derivatives and generalized Dynkin operators,
\emph{Journal of Algebraic Combinatorics} \textbf{38} (2013), 901--913.
\href{https://doi.org/10.1007/s10801-013-0431-3}{doi:10.1007/s10801-013-0431-3};
\href{https://arxiv.org/abs/1206.4990}{arXiv:1206.4990}.

\bibitem{casasmurua}
F. Casas and A. Murua,
An efficient algorithm for computing the Baker--Campbell--Hausdorff series
and some of its applications,
\emph{Journal of Mathematical Physics} \textbf{50} (2009), 033513.
\href{https://doi.org/10.1063/1.3078418}{doi:10.1063/1.3078418};
\href{https://arxiv.org/abs/0810.2656}{arXiv:0810.2656}.

\bibitem{casasmuruanadinic}
F. Casas, A. Murua, and M. Nadinic,
Efficient computation of the Zassenhaus formula,
\emph{Computer Physics Communications} \textbf{183} (2012), 2386--2391.
\href{https://doi.org/10.1016/j.cpc.2012.06.006}{doi:10.1016/j.cpc.2012.06.006};
\href{https://arxiv.org/abs/1204.0389}{arXiv:1204.0389}.

\bibitem{biagi}
S. Biagi, A. Bonfiglioli, and M. Matone,
On the Baker--Campbell--Hausdorff theorem: non-convergence and prolongation issues,
\emph{Linear and Multilinear Algebra} \textbf{68} (2020), 1310--1328;
published online 2018.
\href{https://doi.org/10.1080/03081087.2018.1540534}{doi:10.1080/03081087.2018.1540534};
\href{https://arxiv.org/abs/1805.10089}{arXiv:1805.10089}.

\bibitem{kimura}
T. Kimura,
Explicit description of the Zassenhaus formula,
\emph{Progress of Theoretical and Experimental Physics} \textbf{2017} (2017), 041A03.
\href{https://doi.org/10.1093/ptep/ptx044}{doi:10.1093/ptep/ptx044};
\href{https://arxiv.org/abs/1702.04681}{arXiv:1702.04681}.

\bibitem{wilcox}
R.~M. Wilcox,
Exponential operators and parameter differentiation in quantum physics,
\emph{Journal of Mathematical Physics} \textbf{8} (1967), 962--982.
\href{https://doi.org/10.1063/1.1705306}{doi:10.1063/1.1705306}.

\bibitem{trotter}
H.~F. Trotter,
On the product of semi-groups of operators,
\emph{Proceedings of the American Mathematical Society} \textbf{10} (1959), 545--551.
\href{https://doi.org/10.1090/S0002-9939-1959-0108732-6}{doi:10.1090/S0002-9939-1959-0108732-6}.

\bibitem{suzuki1990}
M. Suzuki,
Fractal decomposition of exponential operators with applications to
many-body theories and Monte Carlo simulations,
\emph{Physics Letters A} \textbf{146} (1990), 319--323.
\href{https://doi.org/10.1016/0375-9601(90)90962-N}{doi:10.1016/0375-9601(90)90962-N}.

\bibitem{suzuki1991}
M. Suzuki,
General theory of fractal path integrals with applications to many-body
theories and statistical physics,
\emph{Journal of Mathematical Physics} \textbf{32} (1991), 400--407.
\href{https://doi.org/10.1063/1.529425}{doi:10.1063/1.529425}.

\bibitem{glauber}
R.~J. Glauber,
Coherent and incoherent states of the radiation field,
\emph{Physical Review} \textbf{131} (1963), 2766--2788.
\href{https://doi.org/10.1103/PhysRev.131.2766}{doi:10.1103/PhysRev.131.2766}.

\bibitem{hall}
B.~C. Hall,
Holomorphic methods in analysis and mathematical physics, in
\emph{First Summer School in Analysis and Mathematical Physics}
(S. P\'erez-Esteva and C. Villegas-Blas, eds.), Contemporary Mathematics
\textbf{260}, American Mathematical Society, 2000, 1--59.
\href{https://arxiv.org/abs/quant-ph/9912054}{arXiv:quant-ph/9912054}.

\end{thebibliography}
\endgroup

\end{document}
